\chapter{Homogeneity, Legendre duality, and b-divisors}
\index{homogeneity!BJ}\index{Legendre transform!homogeneous psh function}\index{b-divisor@$b$-divisor}
\label{chap:bj-homogeneity-bdivisors}

The scaling action separates a psh function into homogeneous pieces.
At the homogeneous PL level, those pieces are recorded on every
birational model at once by Cartier b-divisors. General homogeneous psh
functions need not have modelwise finite divisor support; the numerical
correspondence below concerns their exceptional subclass.  The two descriptions are complementary: Legendre duality
recovers an inhomogeneous function from its homogeneous slices, while
b-divisors turn divisorial values into numerical classes.

We first work on an integral projective variety.  The b-divisor discussion
assumes normality, and the full historical correspondence with nef
b-divisor classes is stated under the smooth characteristic-zero hypotheses
used in \cite[\S6.5]{BJ22}.  The stronger realization needed for the general
envelope theorem is not invoked here; it is proved in
\Cref{chap:bj-cartier-realization}.

\section{Orientation}
\label{sec:ch35-orientation}

\subsection{Three descriptions of one object}
\label{sec:ch35-three-descriptions}

A $\theta$-psh function on the trivially valued analytification is an
extended-real function on a space which is a cone: $X^{\An}$ carries the
scaling action \eqref{eq:valuation-scaling}. A normalized link need not be
compact: the displayed link of a smooth curve in
\Cref{prop:ch08-trivial-curve-scaling}\,(3) is an infinite discrete set.
That example is not an assertion about every transverse slice or every X.  Three descriptions of
such a function are used in \Crefrange{chap:bj-envelope-statement}{chap:bj-envelope-consequences},
and this chapter constructs the translations among them.

\emph{The analytic description.}  A function
$\phi:X^{\An}\to\mathbf R\cup\{-\infty\}$ which is upper semicontinuous,
generically finite, and a decreasing limit of generically finite
Fubini--Study functions in classes converging to $\theta$.  This is
\Cref{def:bj-psh}.  It is the description in which the envelope property
\eqref{eq:bj-envelope-property} is stated, and it directly displays values at proper-support semivaluations.
The homogeneous slices also live on all of the analytification; their
divisorial data recover those values through the psh comparison theorem.

\emph{The homogeneous description.}  A one-parameter family
$(\widehat\phi_\lambda)_\lambda$ of \emph{homogeneous} $\theta$-psh
functions, concave and decreasing in the parameter $\lambda$, from which
$\phi$ is recovered by the pointwise formula
\eqref{eq:legendre-reconstruction-bj}.  Each member of the family is
determined by one homogeneous function on the cone, so all the
inhomogeneous information has been moved into the parameter.  This is the
description used by the reflection arguments of
\Cref{chap:bj-reflection}: homogeneous PL functions have Cartier
b-divisorial coefficients, on which modelwise positivity can be tested.
Legendre reconstruction then retains the inhomogeneous parameter.

\emph{The algebraic description.}  A compatible family of divisors on all
birational models of $X$, that is, a b-divisor, together with a numerical
class on $X$.  This is the description in which positivity becomes a
finite-dimensional condition on each model, namely nefness of
$\pi^*\theta+B_Y$, and in which the whole apparatus of
\Cref{app:intersection-bdivisors} applies.

The chapter proves that these descriptions determine each other, and it
records exactly where each translation is defined.  The last point is not
pedantry.  The forward translation from analytic to homogeneous data is
defined on all of $\PSH_X(\theta)$; the backward translation is defined
only on the image. Arbitrary concave families require a psh supremum and
compatibility at the parameter endpoint; the envelope property alone does
not supply both conditions. Similarly, homogeneous PL functions correspond
to Cartier b-divisors, with an additional nef criterion for psh membership.
The numerical correspondence for exceptional homogeneous psh functions is
stated here under smoothness and characteristic-zero hypotheses, which is why
\Cref{chap:bj-cartier-realization} proves a different and weaker
realization statement for use in arbitrary characteristic.

Three mechanisms drive the chapter.  Scaling orbits are one-variable
convex functions; finite polyhedral elimination turns their homogeneous
minorants back into Fubini--Study functions; and the coefficient pairing
turns homogeneous divisorial data into b-divisors.  The comparison theorem
then upgrades identities on divisorial points to identities on the whole
analytification.  The examples below also mark the limits of this package:
upper-semicontinuous functions need not be determined by their divisorial
values, and divisorial pointwise convergence need not make evaluation at a
proper-support point continuous.

\subsection{Standing hypotheses and notation}
\label{sec:ch35-hypotheses}

Throughout the chapter $k$ is an algebraically closed field equipped with
the trivial absolute value, and
\begin{equation}
\label{eq:ch35-standing}
 X\ \text{is an integral projective $k$-variety},\qquad n:=\dim X .
\end{equation}
Irreducibility is assumed from the outset; it is the hypothesis under which
\cite[\S6]{BJ22} works, and it costs nothing, because
\Cref{thm:bj-basic-closure}\,(6) restricts psh functions to irreducible
components.  We write $v_{\triv}:=v_{X,\triv}$ for the trivial valuation of
$k(X)$, $X^{\val}$ for the full-support locus
\eqref{eq:triv-full-support-locus}, and $X^{\divis}\subset X^{\val}$ for
the divisorial locus of \Cref{chap:trivial-analytification}, which by
convention contains $v_{\triv}$.  Numerical classes are real:
$N^1(X):=N^1(X)_{\mathbf R}$.

Two further hypotheses are imposed section by section and never silently.

\begin{enumerate}
\item \emph{Normality} is assumed for homogeneous Siu division and from
 \Cref{sec:ch35-bdivisors} onwards.
 It is what makes the birational model category of
 \Cref{sec:appF-models} available and what makes
 $\pi_*\cO_Y=\cO_X$ true, which is used in
 \Cref{thm:ch35-semiample-criterion}.
\item \emph{Smoothness and $\operatorname{char}k=0$} are assumed only in
 \Cref{sec:ch35-numerical-bclasses}, where numerical b-divisor classes are formed
 as an inverse limit over \emph{smooth} models.  This is the hypothesis
 under which \Cref{app:intersection-bdivisors} constructs
 $N^1_{\mathrm{Weil}}(\mathfrak X)$ at all; see the discussion preceding
 \eqref{eq:appF-numerical-limits}.
\end{enumerate}

We never form $(-\infty)-(-\infty)$.  Every inequality between
extended-real functions in this chapter is written in the relative form
$\phi\leq\psi+C$ rather than $\phi-\psi\leq C$, following ledger entry
(4) of \Cref{sec:extended-arithmetic} and
\Cref{rem:ledger-uses}.  Suprema and infima are defined in the complete extended real line,
by ledger entry (5); an upper bound is needed to exclude $+\infty$
from a supremum. They interact with finite
translations and positive scalings as recorded in
\Cref{lem:extended-real-safe} and \Cref{lem:safe-supremum}; those two lemmas
are used so often below that we cite them only where the reader might
hesitate.

\subsection{Dependencies}
\label{sec:ch35-dependencies}

The chapter uses: the analytification, the scaling action, the divisorial
locus and its density from \Cref{chap:trivial-analytification}; flag ideals
and the test-configuration description of PL functions from
\Cref{chap:bj-test-configurations}; the Fubini--Study and psh classes from
\Cref{chap:bj-fs-pl-psh}; the one-variable Legendre duality of
\Cref{subsec:legendre-finite} and \Cref{sec:appD-test-curves}; the
extended-real ledger of \Cref{sec:extended-arithmetic}; and the b-divisor
formalism of \Cref{app:intersection-bdivisors}.  Nothing in
\Crefrange{chap:bj-envelope-history}{chap:bj-envelope-consequences} is used,
with one exception which is flagged wherever it occurs: the correspondence
between nef b-divisor \emph{classes} and homogeneous psh functions in
\Cref{sec:ch35-nef-bclasses} rests, through
\Cref{thm:ch35-cartier-psh-criterion}, on the envelope property for ample
classes, which in the smooth characteristic-zero domain of that section was
already available historically.  The general envelope theorem is not used
anywhere in this chapter, and this chapter is used in its proof.

\section{The scaling action on functions}
\label{sec:ch35-scaling}

\subsection{The action and its normalization}
\label{sec:ch35-action}

Recall from \eqref{eq:valuation-scaling} that $\mathbf R_{>0}$ acts
continuously on $X^{\An}$ by $v\mapsto tv$, where
$(tv)(f)=t\,v(f)$ with $t\cdot(+\infty)=+\infty$.  The fixed points are
the trivial semivaluations $v_{Y,\triv}$ attached to irreducible
subvarieties $Y\subset X$, by \eqref{eq:trivial-endpoints} and the
discussion following it.

The valuation scaling $v\mapsto tv$ induces an action on extended-real
functions:
\begin{equation}
\label{eq:function-scaling-action}
                  (t\cdot\phi)(v):=t\,\phi(t^{-1}v),
                  \qquad t>0.
\end{equation}
The prefactor is chosen so that section functions behave well:
\begin{equation}
\label{eq:scaling-fs-weights}
 t\cdot\left(\frac1m\max_i\{\log|s_i|+\lambda_i\}\right)
   =\frac1m\max_i\{\log|s_i|+t\lambda_i\}.
\end{equation}

Two remarks on \eqref{eq:function-scaling-action} are needed before anything
is proved with it.

First, the formula makes sense for extended-real $\phi$: if
$\phi(t^{-1}v)=-\infty$ then $t\,\phi(t^{-1}v)=-\infty$ by ledger entry
(3) of \Cref{sec:extended-arithmetic}, since $t>0$.  No case distinction
is required, and this is exactly the use of that entry recorded in
\Cref{rem:ledger-uses}.

Second, the outer factor $t$ is not a normalization one is free to drop.
Without it, the map $\phi\mapsto\phi(t^{-1}\cdot)$ would still be a group
action, but it would not in general preserve the Fubini--Study class for a
fixed line bundle. For rational scales it instead changes the bundle class
by the same scalar factor as the section slopes.  Identity \eqref{eq:scaling-fs-weights} is the reason: the outer
$t$ cancels the $t^{-1}$ that appears in the section values, leaving the
weights $\lambda_i$ as the only quantities that move.  Since the weights
are exactly the inhomogeneous part of a presentation
\eqref{eq:gf-fs-presentation}, the action \eqref{eq:function-scaling-action}
is the one that dilates the inhomogeneity while fixing the algebraic data.

\begin{lemma}
\label{lem:ch35-action-basic}
Let $\phi,\phi',\phi_i:X^{\An}\to\mathbf R\cup\{-\infty\}$ and
$s,t>0$, $c\in\mathbf R$, $a>0$.  Then, with all expressions defined
by the ledger of \Cref{sec:extended-arithmetic}:
\begin{enumerate}
\item $1\cdot\phi=\phi$ and $s\cdot(t\cdot\phi)=(st)\cdot\phi$;
\item $\phi\leq\phi'$ implies $t\cdot\phi\leq t\cdot\phi'$;
\item $t\cdot(\phi+c)=t\cdot\phi+tc$ and
 $t\cdot(a\phi)=a(t\cdot\phi)$;
\item $t\cdot\bigl(\sup_i\phi_i\bigr)=\sup_i(t\cdot\phi_i)$ and
 $t\cdot\bigl(\inf_i\phi_i\bigr)=\inf_i(t\cdot\phi_i)$; in particular the
 action commutes with finite maxima and with decreasing limits;
\item if $\phi$ is upper semicontinuous, so is $t\cdot\phi$;
\item \eqref{eq:scaling-fs-weights} holds, and more generally
 $t\cdot\phi\in\cH^{\mathrm{gf}}_{\mathbf R}(L)$ whenever
 $\phi\in\cH^{\mathrm{gf}}_{\mathbf R}(L)$.
\end{enumerate}
\par
\end{lemma}

\begin{proof}
(1)  For $t=1$ there is nothing to prove.  In general
\[
 \bigl(s\cdot(t\cdot\phi)\bigr)(v)
 =s\,(t\cdot\phi)(s^{-1}v)
 =s\,t\,\phi(t^{-1}s^{-1}v)
 =(st)\,\phi\bigl((st)^{-1}v\bigr),
\]
where the middle equality uses associativity of the scaling action on
$X^{\An}$ and the outer equality uses $s(t\,r)=(st)r$ in $\Rbar$ for
$s,t>0$ and $r\in\mathbf R\cup\{-\infty\}$, which is ledger entry (3).

(2)  Multiplication by $t>0$ is an order isomorphism of
$\mathbf R\cup\{-\infty\}$; this is part of
\Cref{lem:extended-real-safe}\,(2).

(3)  The first identity is
$t(\phi(t^{-1}v)+c)=t\phi(t^{-1}v)+tc$, valid also when
$\phi(t^{-1}v)=-\infty$ because both sides are then $-\infty$.  The
second is associativity of positive scalar multiplication.

(4)  Both identities are \Cref{lem:extended-real-safe}\,(2) applied
pointwise at $t^{-1}v$: multiplication by $t>0$ commutes with suprema
and infima in $\Rbar$.  A decreasing net is an infimum, by
\Cref{lem:monotone-net-converges}.

(5)  The map $v\mapsto t^{-1}v$ is a homeomorphism of $X^{\An}$, being
continuous with continuous inverse $v\mapsto tv$ by
\eqref{eq:valuation-scaling}.  Composition with a homeomorphism preserves
upper semicontinuity, and multiplication by $t>0$ is an increasing
homeomorphism of $\Rbar$.

(6)  For a nonzero section $s$ of $mL$, the function
$\log|s|(v)=-v(s)$ of \eqref{eq:section-log-value} satisfies
$\log|s|(t^{-1}v)=t^{-1}\log|s|(v)$, because $v\mapsto v(s)$ is
positively homogeneous.  Hence
\[
 \begin{aligned}
 t\cdot\Bigl(\tfrac1m\max_i\{\log|s_i|+\lambda_i\}\Bigr)(v)
 &=\tfrac tm\max_i\{t^{-1}\log|s_i|(v)+\lambda_i\}\\
 &=\tfrac1m\max_i\{\log|s_i|(v)+t\lambda_i\},
 \end{aligned}
\]
which is \eqref{eq:scaling-fs-weights}; the middle step uses
\Cref{lem:extended-real-safe}\,(2) once more, since a maximum is a supremum
over a finite index set.  The right-hand side is again a presentation of the
form \eqref{eq:gf-fs-presentation} with the same sections, so the generic
finiteness condition of \Cref{def:gf-fs-functions} is unchanged.
\end{proof}

\begin{warning}{The action is not by translation}
Assertion (3) of \Cref{lem:ch35-action-basic} says that
$t\cdot(\phi+c)=t\cdot\phi+tc$, not $t\cdot\phi+c$.  For $\phi\in\PSH_X(\theta)$ its supremum is finite and
$\sup(t\cdot\phi)=t\sup\phi$, since scaling is a bijection of
$X^{\An}$. Consequently the normalization $N\phi:=\phi-\sup\phi$
\emph{does} commute with the action:
\[
 N(t\cdot\phi)=t\cdot\phi-t\sup\phi=t\cdot(N\phi).
\]
A fixed added constant scales by $t$, while subtracting the corresponding
supremum is equivariant.
\end{warning}

\subsection{Homogeneous functions}
\label{sec:ch35-homogeneous-functions}

\begin{definition}[Homogeneous function]
\label{def:homogeneous-function}
A function $\psi:X^{\An}\to\mathbf R\cup\{-\infty\}$ is homogeneous if
\begin{equation}
\label{eq:homogeneous-function}
                         \psi(tv)=t\psi(v)
\end{equation}
for every $t>0$, with $t(-\infty)=-\infty$.
\end{definition}

\begin{lemma}
\label{lem:ch35-homogeneity-reformulations}
For $\psi:X^{\An}\to\mathbf R\cup\{-\infty\}$ the following are
equivalent:
\begin{enumerate}
\item $\psi$ is homogeneous in the sense of
 \Cref{def:homogeneous-function};
\item $t\cdot\psi=\psi$ for every $t>0$;
\item $t\cdot\psi=\psi$ for every $t$ in some subset of
 $\mathbf R_{>0}$ generating a dense subgroup, and $\psi$ is continuous
 along scaling orbits.
\end{enumerate}
Moreover a homogeneous $\psi$ satisfies
$\psi(v_{Y,\triv})\in\{0,-\infty\}$ for every irreducible
$Y\subseteq X$, and $\psi(v)=0$ at every $v$ at which $\psi$ is bounded on the
scaling orbit through $v$; equivalently, a homogeneous $\psi$ that is
nonzero at $v$ is unbounded on the orbit of $v$.  This is local and
cannot be globalized: $\psi=\log|x|$ on $\mathbf P^1$ is homogeneous and
bounded on the orbit through every $q\neq[0:1]$, yet $\psi\not\equiv0$.
\end{lemma}

\begin{proof}
(1)$\Leftrightarrow$(2).  By definition
$(t\cdot\psi)(v)=t\psi(t^{-1}v)$; this equals $\psi(v)$ for all $v$ and
all $t>0$ if and only if $\psi(t^{-1}v)=t^{-1}\psi(v)$ for all $v,t$,
which is \eqref{eq:homogeneous-function} after renaming $t^{-1}$ as
$t$.  The manipulation is legitimate at points where $\psi=-\infty$
because both sides are then $-\infty$.

(2)$\Rightarrow$(3) also gives orbit continuity: the function
$t\mapsto\psi(tv)$ is either the real linear function $t\psi(v)$
or the constant $-\infty$ function.  For (3)$\Rightarrow$(2), the set
$\{t>0:t\cdot\psi=\psi\}$ is a subgroup of $\mathbf R_{>0}$, by
\Cref{lem:ch35-action-basic}\,(1); if it is dense and
$t\mapsto\psi(tv)$ is continuous for each $v$, then the identity
$\psi(tv)=t\psi(v)$ passes to the closure, the two sides being continuous
in $t$ for fixed $v$.

For the value at a fixed point, $tv_{Y,\triv}=v_{Y,\triv}$ gives
$\psi(v_{Y,\triv})=t\psi(v_{Y,\triv})$ for every $t>0$; the only
elements of $\mathbf R\cup\{-\infty\}$ fixed by every positive dilation are
$0$ and $-\infty$.

Finally, suppose $\psi(v_0)=:c\neq0$ for some $v_0$.  If
$c\in\mathbf R\setminus\{0\}$ then $\psi(tv_0)=tc$ is unbounded as
$t\to+\infty$; if $c=-\infty$ then $\psi\equiv-\infty$ on the orbit.
In either case $\psi$ is unbounded on that orbit, which proves the
contrapositive of the last assertion.
\end{proof}

A finite homogeneous function vanishes at the fixed point
$v_{X,\triv}$.  Nonzero homogeneous functions cannot be bounded on a full
scaling orbit.  This explains why homogeneous PL functions, defined below
on $X^{\val}$, are generally not ordinary bounded PL functions on
$X^{\An}$.

\begin{definition}[Homogeneous FS and psh functions]
\label{def:homogeneous-fs-psh}
For $L\in\Pic(X)_{\mathbf Q}$, set
\[
 \cH_{\mathrm{hom}}(L):=\cH^{\mathrm{gf}}_0(L).
\]
Thus its elements have presentations
\begin{equation}
\label{eq:homogeneous-fs-presentation}
                  \psi=\frac1m\max_i\log|s_i|.
\end{equation}
For $\theta\in N^1(X)$, let
\begin{equation}
\label{eq:homogeneous-psh-space}
 \PSH_{\mathrm{hom},X}(\theta)
 :=\{\psi\in\PSH_X(\theta)\mid\psi\text{ is homogeneous}\}.
\end{equation}
\end{definition}

Thus $\cH_{\mathrm{hom}}(L)$ is the subclass of
\Cref{def:gf-fs-functions} in which all weights vanish, and the containment
$\cH_{\mathrm{hom}}(L)\subset\PSH_{\mathrm{hom},X}(c_1(L))$ is immediate
from \Cref{def:bj-psh} with the constant net.  Zero weights make
\eqref{eq:homogeneous-fs-presentation} homogeneous.  Conversely, a
homogeneous generically finite Fubini--Study function admits a weight-zero
presentation even if its original weights are real. Indeed its root value
is zero, so $\max_i\lambda_i=0$. In
\eqref{eq:scaling-fs-weights}, let $t\to+\infty$: negative-weight terms
vanish from the maximum, while homogeneity keeps the left side fixed.
The remaining zero-weight sections therefore present the same function,
including at points where every such section vanishes. A nonminimal
presentation may still contain dominated terms with negative weights.

\begin{proposition}
\label{prop:ch35-homogeneous-normalization}
Let $\theta\in N^1(X)$ and $\psi\in\PSH_{\mathrm{hom},X}(\theta)$.
Then
\begin{equation}
\label{eq:ch35-homogeneous-normalized}
                     \psi(v_{\triv})=0,\qquad \psi\leq0
                     \quad\text{on }X^{\An},
\end{equation}
and $\PSH_{\mathrm{hom},X}(\theta)\subset\PSH^{\sup}_X(\theta)$.
Moreover each $\PSH_{\mathrm{hom},X}(\theta)$ is stable under finite
maxima, while, as $\theta$ varies, the collection of pairs
$(\theta,\psi)$ with
$\psi\in\PSH_{\mathrm{hom},X}(\theta)$ is a graded convex cone under
positive scalar multiplication and sums, whenever the result is generically
finite.  At an endpoint of a convex combination, a zero-coefficient term is
omitted, never evaluated as $0\cdot(-\infty)$.
\end{proposition}

\begin{proof}
A $\theta$-psh function is finite at $v_{\triv}$ by generic finiteness in
\Cref{def:bj-psh}, so \Cref{lem:ch35-homogeneity-reformulations} forces
$\psi(v_{\triv})=0$.  A $\theta$-psh function is decreasing for the
order \eqref{eq:semivaluation-order} and satisfies the maximum principle
$\sup_{X^{\An}}\psi=\psi(v_{\triv})$; this is
\cite[Proposition~4.12]{BJ22}, and for our purposes it follows already from
the presentation \eqref{eq:gf-fs-presentation}: for a generically finite
Fubini--Study function $\phi=\frac1m\max_i\{\log|s_i|+\lambda_i\}$ one has
$\log|s_i|(v)=-v(s_i)\leq0=\log|s_i|(v_{\triv})$, so
$\sup\phi=\phi(v_{\triv})=\frac1m\max_i\lambda_i$, and the property
survives decreasing limits.  Hence $\sup\psi=\psi(v_{\triv})=0$, which is
\eqref{eq:ch35-homogeneous-normalized} and the inclusion into
\eqref{eq:normalized-psh-space}.

For the cone statement, \Cref{thm:bj-basic-closure}\,(2)--(3) gives the psh part
of each assertion, and homogeneity is preserved because
$t\cdot$ commutes with finite maxima and positive linear combinations by
\Cref{lem:ch35-action-basic}\,(3)--(4).
\end{proof}

\begin{warning}{Homogeneous means $\leq0$, not bounded}
\Cref{prop:ch35-homogeneous-normalization} says that a homogeneous psh
function is nonpositive.  It does not say it is bounded: by
\Cref{lem:ch35-homogeneity-reformulations} a nonzero homogeneous function is
unbounded below along some scaling ray, and it may equal $-\infty$ at a
proper-support point.  Consequently $\PSH_{\mathrm{hom},X}(\theta)$ is
\emph{not} contained in $C^0(X^{\An})$ and is not contained in the space
$\cE^\infty$ of bounded potentials.  This is the single most common source
of a false step in this part: an estimate proved for bounded psh functions
does not transfer to homogeneous ones without an argument.
\end{warning}

\subsection{Convexity along orbits}
\label{sec:ch35-convexity}

The next proposition is the analytic engine of the whole chapter.  Once it
is available, every Legendre statement below is the one-variable duality of
\Cref{prop:concave-biconjugation} applied at a fixed analytic point.

\begin{proposition}
\label{prop:ch35-scaling-convex}
Let $\theta\in N^1(X)$, $\phi\in\PSH_X(\theta)$, and $v\in X^{\An}$.
Then
\begin{equation}
\label{eq:ch35-orbit-convex}
 t\longmapsto(t\cdot\phi)(v)=t\,\phi(t^{-1}v)
\end{equation}
is a convex function of $t\in\mathbf R_{>0}$ with values in
$\mathbf R\cup\{-\infty\}$; it is nonincreasing if $\phi\leq0$, and
more generally
\begin{equation}
\label{eq:ch35-orbit-slope-bound}
 (t\cdot\phi)(v)\leq t\sup\phi
 \qquad\text{for every }t>0 .
\end{equation}
Furthermore $s\mapsto\phi(sv)$ is convex and nonincreasing on
$[0,+\infty)$.
\end{proposition}

\begin{proof}
\emph{Step 1: the Fubini--Study case.}
Let $\phi=\frac1m\max_i\{\log|s_i|+\lambda_i\}$ belong to
$\cH^{\mathrm{gf}}_{\mathbf R}(L)$.  By
\eqref{eq:scaling-fs-weights},
\[
 (t\cdot\phi)(v)=\frac1m\max_i\{\log|s_i|(v)+t\lambda_i\},
\]
which for fixed $v$ is the maximum of finitely many affine functions of
$t$, with the convention that a term with
$\log|s_i|(v)=-\infty$ contributes the constant $-\infty$.  A maximum of
finitely many affine functions with values in
$\mathbf R\cup\{-\infty\}$ is convex: on the set where at least one term
is finite the assertion is the usual one, and if all terms are
$-\infty$ at some $v$ the function is identically $-\infty$ in
$t$, which is convex.  If moreover $\phi\leq0$, then
$\sup\phi=\frac1m\max_i\lambda_i\leq0$ by the computation in the proof of
\Cref{prop:ch35-homogeneous-normalization}, so every slope
$\lambda_i/m$ that is attained in the maximum is nonpositive and
$t\mapsto(t\cdot\phi)(v)$ is nonincreasing.  Estimate
\eqref{eq:ch35-orbit-slope-bound} holds because
$\log|s_i|\leq0$ pointwise, whence
$(t\cdot\phi)(v)\leq\frac tm\max_i\lambda_i=t\sup\phi$.

\emph{Step 2: decreasing limits.}
Let $\phi_a\searrow\phi$ be a decreasing net as in
\eqref{eq:bj-psh-regularization}.  By
\Cref{lem:ch35-action-basic}\,(4), $t\cdot\phi_a\searrow t\cdot\phi$
pointwise for each fixed $t$.  Thus, for fixed $v$, the functions
$t\mapsto(t\cdot\phi_a)(v)$ are convex and decrease pointwise in $t$ to
$t\mapsto(t\cdot\phi)(v)$.  A pointwise limit of a decreasing net of
convex $\mathbf R\cup\{-\infty\}$-valued functions on an interval is
convex: for $0<\sigma<1$ and
$t=(1-\sigma)t_0+\sigma t_1$ one has
\[
 (t\cdot\phi_a)(v)
 \leq(1-\sigma)(t_0\cdot\phi_a)(v)+\sigma(t_1\cdot\phi_a)(v),
\]
and both sides converge monotonically, the right-hand side being a sum of
two decreasing nets and hence, by
\Cref{lem:extended-real-safe}\,(5) applied whenever one summand is finite
and by inspection when both are $-\infty$, converging to
$(1-\sigma)(t_0\cdot\phi)(v)+\sigma(t_1\cdot\phi)(v)$.  The bound \eqref{eq:ch35-orbit-slope-bound} passes to the limit:
the maximum principle and convergence at the trivial valuation give
\[
 \sup\phi_a=\phi_a(v_{\triv})\searrow
 \phi(v_{\triv})=\sup\phi .
\]
To prove monotonicity when $\phi\leq0$, put
$F(t):=(t\cdot\phi)(v)$. A convex function on $(0,+\infty)$ bounded
above by a finite linear function is either finite everywhere or
identically $-\infty$. Indeed, if it is $-\infty$ at one parameter,
express any other parameter as a strict convex combination of that
parameter and a third one; convexity forces its value to be $-\infty$.
In the finite case $F(t)\leq0$. If $F(t_2)>F(t_1)$ for $t_2>t_1$,
the positive secant slope and convexity force $F(t)\to+\infty$ as
$t\to+\infty$, a contradiction. The identically $-\infty$ case is
immediate. This proves monotonicity without assuming the approximants
are themselves nonpositive. The endpoint cases
$\sigma\in\{0,1\}$ are kept separate, so no coefficient $0$ is ever
multiplied by $-\infty$.

\emph{Step 3: the second parametrization.}
Fix $v$ and put $g(s):=\phi(sv)$ for $s>0$.  Then
$(t\cdot\phi)(v)=t\,g(t^{-1})$.  For the Fubini--Study case,
$g(s)=\frac1m\max_i\{-s\,v(s_i)+\lambda_i\}$ is again a finite maximum of
affine functions of $s$, hence convex, and its slopes
$-v(s_i)/m$ are nonpositive because $v(s_i)\geq0$; so $g$ is
nonincreasing.  Decreasing limits preserve both properties as in Step 2.
The exact endpoint statement is
\Cref{thm:bj-basic-closure}\,\cref{itm:bj-closure:radial}:
$g(0):=\phi(v_{s(v),\triv})$ makes $g$ convex and nonincreasing on
$[0,+\infty)$.
\end{proof}

\begin{remark}[Where convexity comes from]
\label{rem:ch35-convexity-source}
Convexity in \eqref{eq:ch35-orbit-convex} is not a soft consequence of
plurisubharmonicity.  It is inherited, through the definition of
$\PSH_X(\theta)$ as a class of decreasing limits, from the elementary fact
that a maximum of affine functions is convex, and the affine functions in
question are indexed by the sections in a Fubini--Study presentation.  This
is the same mechanism that makes tropical objects convex, and it is why the
Legendre transform below is a one-variable operation at
each fixed analytic point even though $\PSH_X(\theta)$ is
infinite-dimensional.
\end{remark}

\begin{example}[The orbit of a Fubini--Study function on $\mathbf P^1$]
\label{ex:ch35-p1-orbit}
Let $X=\mathbf P^1_k$, $L=\cO(1)$ with homogeneous coordinates
$x,y$, and let
$\phi=\max\{\log|x|+a,\log|y|+b\}$ be the function
\eqref{eq:p1-fs} of \Cref{ex:p1-fs-function}.  Write $p=[0:1]$ and
$v=s\,\ord_p$.  Trivializing $L$ near $p$ by $y$, one has
$\log|y|(v)=0$ and $\log|x|(v)=-s$, so
\[
 \phi(s\,\ord_p)=\max\{-s+a,\ b\},
 \qquad
 (t\cdot\phi)(\ord_p)=\max\{-1+ta,\ tb\}.
\]
Both displayed functions are maxima of two affine functions, in $s$ and
in $t$ respectively, in accordance with
\Cref{prop:ch35-scaling-convex}.  The first is nonincreasing in $s$, as it
must be; the second is nonincreasing in $t$ precisely when
$\max\{a,b\}\leq0$, that is, when $\sup\phi\leq0$.  With $a=0$ and
$b=-1$, for instance,
\[
 \phi(s\,\ord_p)=\max\{-s,-1\},
 \qquad
 (t\cdot\phi)(\ord_p)=\max\{-1,-t\},
\]
and the break occurs at $s=1$, respectively $t=1$.
\end{example}

\section{Homogeneous PL and Fubini--Study functions}
\label{sec:ch35-hom-pl-fs}

\subsection{Ideal functions and the homogeneous PL cone}
\label{sec:ch35-hom-pl}

If $\mathfrak b\subset\cO_X$ is a nonzero ideal, then
$\log|\mathfrak b|=-v(\mathfrak b)$ is homogeneous.  The cone
\begin{equation}
\label{eq:positive-homogeneous-pl}
 \PL_{\mathrm{hom}}^+(X)
 :=\{m^{-1}\log|\mathfrak b|\mid m\geq1,\ 0\ne\mathfrak b\subset\cO_X\}
\end{equation}
is stable under addition, positive rational scaling, and finite maxima.
The vector space it spans on $X^{\val}$ is
$\PL_{\mathrm{hom}}(X)$. Its real span is denoted
$\PL_{\mathrm{hom},\mathbf R}(X)$; it consists of finite real linear
combinations of ideal functions on $X^{\val}$. Except for zero, its elements are unbounded
along some ray and hence do not belong to the ordinary $\PL(X)$.

Each of these assertions is proved in the next proposition, together with
the one fact that makes $\PL_{\mathrm{hom}}(X)$ usable as a linear space
rather than a mere cone.

\begin{proposition}
\label{prop:ch35-hom-pl-semiring}
Let $\mathfrak b,\mathfrak b'\subset\cO_X$ be nonzero coherent ideals.
Then, on all of $X^{\An}$,
\begin{equation}
\label{eq:ch35-ideal-tropical}
 \log|\mathfrak b\mathfrak b'|
   =\log|\mathfrak b|+\log|\mathfrak b'|,
 \qquad
 \log|\mathfrak b+\mathfrak b'|
   =\max\{\log|\mathfrak b|,\log|\mathfrak b'|\},
\end{equation}
and each $\log|\mathfrak b|$ is homogeneous, nonpositive, upper
semicontinuous, and finite on $X^{\val}$.  Consequently:
\begin{enumerate}
\item $\PL_{\mathrm{hom}}^+(X)$ is stable under sums, under finite maxima,
 and under multiplication by $\mathbf Q_{>0}$, and it contains $0$;
\item every element of the $\mathbf Q$-vector space
 $\PL_{\mathrm{hom}}(X)$ spanned by
 $\PL_{\mathrm{hom}}^+(X)$ inside $C^0(X^{\val})$ is a difference of two
 elements of $\PL_{\mathrm{hom}}^+(X)$;
\item $\PL_{\mathrm{hom}}(X)\cap\PL(X)=\{0\}$;
\item $\PL_{\mathrm{hom}}^+(X)$ is stable under the scaling action, which
 acts on it trivially.
\end{enumerate}
\par
\end{proposition}

\begin{proof}
The two identities in \eqref{eq:ch35-ideal-tropical} are
\eqref{eq:ideal-tropical-identities} of
\Cref{chap:trivial-analytification} after multiplication by $-1$, which
exchanges minima and maxima and sums with sums; the third identity there
says that $\log|\mathfrak b|(v)=-\infty$ exactly when
$s(v)\subset V(\mathfrak b)$, so $\log|\mathfrak b|$ is finite at every
point of full support, in particular on $X^{\val}$, because
$\mathfrak b\neq0$.  Homogeneity is
$(tv)(\mathfrak b)=t\,v(\mathfrak b)$, which follows from
\eqref{eq:ideal-evaluation-triv} because the minimum defining
$v(\mathfrak b)$ is taken over a fixed finite generating set at the centre
and the centre is unchanged by scaling.  Nonpositivity is
$v(\mathfrak b)\geq0$, valid because $v$ is nonnegative on the local
ring at its centre.  Continuity of $v\mapsto v(\mathfrak b)$ into
$[0,+\infty]$ is part of \Cref{chap:trivial-analytification}; upper
semicontinuity of $\log|\mathfrak b|$ into
$[-\infty,0]$ follows.

(1)  Stability under sums and finite maxima is
\eqref{eq:ch35-ideal-tropical} applied to products and sums of ideals, after
passing to a common denominator: if
$\psi=m^{-1}\log|\mathfrak b|$ and
$\psi'=(m')^{-1}\log|\mathfrak b'|$, then
$\psi+\psi'=(mm')^{-1}\log|\mathfrak b^{m'}(\mathfrak b')^{m}|$ and
$\max\{\psi,\psi'\}=(mm')^{-1}\log|\mathfrak b^{m'}+(\mathfrak b')^{m}|$.
For $c=p/q\in\mathbf Q_{>0}$, $c\psi=(qm)^{-1}\log|\mathfrak b^{p}|$.
Finally $0=\log|\cO_X|$.

(2)  By definition $\PL_{\mathrm{hom}}(X)$ is the set of finite
$\mathbf Q$-linear combinations $\sum_ic_i\psi_i$ with
$\psi_i\in\PL_{\mathrm{hom}}^+(X)$.  Splitting the coefficients according
to sign and using (1) writes such a combination as
$\psi'-\psi''$ with $\psi',\psi''\in\PL_{\mathrm{hom}}^+(X)$.  The
difference is unambiguous on $X^{\val}$, where both terms are finite; it
is for this reason that $\PL_{\mathrm{hom}}(X)$ is defined as a space of
functions on $X^{\val}$ and not on $X^{\An}$.

(3)  A function in $\PL(X)$ is bounded, being continuous on the compact
space $X^{\An}$ by \Cref{def:bj-pl-function} and
\Cref{thm:triv-berkovich-compactness}.  A homogeneous function that is
bounded on a scaling orbit vanishes on it, by
\Cref{lem:ch35-homogeneity-reformulations}.  Since every point lies on an
orbit, the intersection is $\{0\}$.

(4)  Immediate from homogeneity and
\Cref{lem:ch35-homogeneity-reformulations}.
\end{proof}

\begin{warning}{Fractional ideals and differences of ideal functions}
For a nonzero coherent fractional ideal $\mathfrak a$, its logarithm on
$X^{\val}$ belongs to $\PL_{\mathrm{hom}}(X)$. Put
$\mathfrak b:=\{f\in\cO_X:f\mathfrak a\subset\cO_X\}$ and
$\mathfrak b':=\mathfrak b\mathfrak a\subset\cO_X$; both are nonzero
coherent ideals. Multiplicativity gives
$\log|\mathfrak a|=\log|\mathfrak b'|-\log|\mathfrak b|$.
This is \cite[Example~6.4]{BJ22}, a one-way inclusion. It does not say
that every difference, or every rational multiple, is the logarithm of
one fractional ideal on X.

For example, on normal $X=\mathbf P^2$ fix a closed point $q$ and set
$\chi:=-\log|\mathfrak m_q|$. This function vanishes at every prime
divisor of X, but is positive at the exceptional valuation of the blowup
of q. If $\chi=\log|\mathfrak a|$, all height-one orders of every local
generator of $\mathfrak a$ would be nonnegative. Normality would imply
$\mathfrak a\subset\cO_X$, hence $\log|\mathfrak a|\leq0$, a
contradiction. General differences are instead realized by Cartier
b-divisors on a common model in \Cref{prop:cartier-homogeneous-pl}.

Nor may the displayed difference be used at a proper-support point where
both ideal functions equal $-\infty$. The linear space is defined on
$X^{\val}$, where both are finite; any psh extension to $X^{\An}$ is
handled separately by divisorial determination.
\end{warning}

\begin{example}[Homogeneous PL functions on the trivially valued line]
\label{ex:ch35-p1-hom-pl}
Let $X=\mathbf P^1_k$, so that $X^{\An}$ is the star-shaped cone of
\Cref{ex:triv-p1-tree}, with one ray $\{t\ord_p\}_{t\in[0,\infty]}$ for
each $p\in X(k)$.  Every nonzero ideal of $\cO_X$ has the form
$\mathfrak b=\cO_X(-\sum_ia_ip_i)$ with $a_i\in\mathbf Z_{>0}$ and the
$p_i$ distinct closed points, so
\[
 \log|\mathfrak b|(t\,\ord_p)=
 \begin{cases}
  -a_it,&p=p_i,\\
  0,&p\notin\{p_1,\ldots,p_r\}.
 \end{cases}
\]
Hence $\PL_{\mathrm{hom}}^+(X)$ consists of the functions that are linear
with nonpositive rational slope on finitely many rays and zero on all the
others, and $\PL_{\mathrm{hom}}(X)$ is the same description with slopes of
arbitrary sign.  In particular the only element that is bounded is $0$,
confirming \Cref{prop:ch35-hom-pl-semiring}\,(3), and the value at the
endpoint $v_{p_i,\triv}$ of a ray with $a_i>0$ is $-\infty$.
\end{example}

\subsection{Homogeneous Fubini--Study functions}
\label{sec:ch35-hom-fs}

\begin{proposition}\label{prop:hom-fs-ideal}
For a $\mathbf Q$-line bundle $L$, a function $\psi$ lies in
$\cH_{\mathrm{hom}}(L)$ if and only if
\begin{equation}
\label{eq:hom-fs-ideal}
                    \psi=m^{-1}\log|\mathfrak b|
\end{equation}
for a nonzero ideal $\mathfrak b$ such that $mL$ is a line bundle and
$mL\otimes\mathfrak b$ is globally generated.
\end{proposition}

\begin{proof}
The sections in \eqref{eq:homogeneous-fs-presentation}, expressed in local
frames, generate an ideal $\mathfrak b$, and evaluation of a finite set of
generators gives
$\max_i\log|s_i|=\log|\mathfrak b|$.  Conversely, choose global
generators of $mL\otimes\mathfrak b$ and reverse the construction.
Nonzeroness of $\mathfrak b$ is exactly generic finiteness.

In detail, suppose first $\psi\in\cH_{\mathrm{hom}}(L)$ and write
$\psi=\frac1m\max_{1\leq i\leq r}\log|s_i|$ with
$0\neq s_i\in H^0(X,mL)$ and $mL$ an honest line bundle, as in
\eqref{eq:homogeneous-fs-presentation}.  Fix a point $\xi\in X$ and a
local trivialization $e$ of $mL$ near $\xi$; write
$s_i=f_ie$ with $f_i\in\cO_{X,\xi}$.  Let
$\mathfrak b\subset\cO_X$ be the ideal locally generated by the $f_i$.
This is independent of the trivialization, since a change of frame
multiplies all $f_i$ by one unit, and it is coherent because it is
generated by the images of the finite-dimensional space spanned by the
$s_i$ under the evaluation map
$H^0(X,mL)\otimes\cO_X\to mL$, twisted by $-mL$.  It is nonzero because
some $s_i$ is nonzero and $X$ is integral.  Now for
$v\in X^{\An}$ with centre $c(v)=\xi$ we have, by
\eqref{eq:section-log-value} and \eqref{eq:ideal-evaluation-triv},
\[
 \max_i\log|s_i|(v)=-\min_iv(f_i)=-v(\mathfrak b)
 =\log|\mathfrak b|(v),
\]
which is \eqref{eq:hom-fs-ideal}.  Moreover
$mL\otimes\mathfrak b$ is globally generated, because the $s_i$
themselves are global sections of $mL\otimes\mathfrak b$ generating it by
construction.

Conversely, assume \eqref{eq:hom-fs-ideal} with
$mL\otimes\mathfrak b$ globally generated and $\mathfrak b\neq0$.
Choose finitely many global sections $\sigma_1,\ldots,\sigma_r$ of
$mL\otimes\mathfrak b$ generating it, and let $s_i$ be their images
under the inclusion $mL\otimes\mathfrak b\hookrightarrow mL$.  Then the
$s_i$ are nonzero sections of $mL$, they generate
$\mathfrak b$ in every local frame, and the computation above reversed
gives $\frac1m\max_i\log|s_i|=m^{-1}\log|\mathfrak b|=\psi$.  Since
$\mathfrak b\neq0$, not all $s_i$ vanish identically, which is the
generic finiteness condition of \Cref{def:gf-fs-functions}.  All weights are
zero, so $\psi\in\cH^{\mathrm{gf}}_0(L)=\cH_{\mathrm{hom}}(L)$.
\end{proof}

\begin{corollary}
\label{cor:ch35-ample-exhausts}
For every ample $L\in\Pic(X)_{\mathbf Q}$,
\begin{equation}
\label{eq:ch35-ample-exhausts}
 \mathbf Q_{>0}\cdot\cH_{\mathrm{hom}}(L)=\PL^+_{\mathrm{hom}}(X),
\end{equation}
and $\cH_{\mathrm{hom}}(L)$ spans the $\mathbf Q$-vector space
$\PL_{\mathrm{hom}}(X)$.  Moreover
\[
 \PL^+_{\mathrm{hom}}(X)
   =\bigcup_{L\ \mathrm{ample}}\cH_{\mathrm{hom}}(L).
\]
\end{corollary}

\begin{proof}
The inclusion $\subseteq$ in \eqref{eq:ch35-ample-exhausts} is
\Cref{prop:hom-fs-ideal} together with
\Cref{prop:ch35-hom-pl-semiring}\,(1).  For $\supseteq$, let
$\psi=m^{-1}\log|\mathfrak b|\in\PL^+_{\mathrm{hom}}(X)$.  Since $L$ is
ample, there is $q\geq1$ such that $qmL$ is an honest line bundle and
$qmL\otimes\mathfrak b$ is globally generated; this is Serre vanishing and
global generation applied to the coherent sheaf $\mathfrak b$.  Then
\Cref{prop:hom-fs-ideal} gives
$(qm)^{-1}\log|\mathfrak b|\in\cH_{\mathrm{hom}}(L)$, so
$\psi=q\cdot(qm)^{-1}\log|\mathfrak b|$ lies in
$\mathbf Q_{>0}\cH_{\mathrm{hom}}(L)$.  The spanning assertion follows
from \Cref{prop:ch35-hom-pl-semiring}\,(2), and the last assertion is the
same argument with $L$ allowed to vary, using that
$\cH_{\mathrm{hom}}(L)\subseteq\cH_{\mathrm{hom}}(L')$ whenever
$L'-L$ is ample, by \Cref{lem:gf-fs-operations}\,(1) applied to a
homogeneous Fubini--Study function for $L'-L$.
\end{proof}

\begin{example}[A point ideal on the projective line]
\label{ex:ch35-point-ideal-p1}
Let $X=\mathbf P^1_k$, let $p\in X$ be a closed point, and take
$L=\mathcal O_X(1)$.  The ideal sheaf
$\mathfrak m_p=\mathcal O_X(-p)$ satisfies
\[
       L\otimes\mathfrak m_p\simeq\mathcal O_X,
\]
so it is globally generated.  Thus \Cref{prop:hom-fs-ideal} gives the
homogeneous Fubini--Study function
\[
                 \psi=\log|\mathfrak m_p|.
\]
For the divisorial valuation $v_c=c\,\ord_p$, with $c>0$, one has
\[
       \psi(v_c)=-v_c(\mathfrak m_p)=-c,
       \qquad
       \psi(tv_c)=t\psi(v_c)\quad(t>0).
\]
At the trivial valuation the value is zero.  Hence this one-dimensional
calculation displays both conventions at once: homogeneous functions vanish
at the vertex of the valuation cone, but a nonzero one is unbounded below
along a divisorial ray.

Concretely, in homogeneous coordinates with $p=[0:1]$ one may take
$m=1$ and the single section $s=x\in H^0(X,\cO(1))$; then
$\psi=\log|x|$, and the presentation
\eqref{eq:homogeneous-fs-presentation} has one term and no weight.  This is
the function obtained from \eqref{eq:p1-fs} by deleting the section $y$,
as noted at the end of \Cref{ex:p1-fs-function}, and it equals $-\infty$
at the proper-support point $v_{p,\triv}$.
\end{example}

\begin{example}[A monomial homogeneous function]
\label{ex:ch35-monomial-hom-fs}
Let $X=\mathbf P^2_k$ with homogeneous coordinates $x,y,z$, let
$L=\cO(1)$, and put
$\psi=\tfrac12\max\{\log|x^2|,\ \log|xy|,\ \log|y^2|\}$.  The three
displayed sections of $2L=\cO(2)$ have common zero locus the single point
$q=[0:0:1]$, and in a local frame they generate the ideal
$\mathfrak b=\mathfrak m_q^2$; thus $2L\otimes\mathfrak b$ is globally
generated and \Cref{prop:hom-fs-ideal} gives
$\psi=\tfrac12\log|\mathfrak m_q^2|\in\cH_{\mathrm{hom}}(L)$.
On the chart $z=1$ with affine coordinates $u=x/z$,
$w=y/z$, and for the monomial valuation $v_\alpha$ of
\eqref{eq:quasimonomial-valuation} with
$v_\alpha(u)=\alpha_1$, $v_\alpha(w)=\alpha_2$, one gets
\[
 \psi(v_\alpha)=-\tfrac12\min\{2\alpha_1,\ \alpha_1+\alpha_2,\ 2\alpha_2\}
 =-\min\{\alpha_1,\alpha_2\},
\]
because $\alpha_1+\alpha_2\geq2\min\{\alpha_1,\alpha_2\}$.  Thus
$\psi=\log|\mathfrak m_q|$ on the monomial cone, in agreement with the
identity $\tfrac12\log|\mathfrak b^2|=\log|\mathfrak b|$ of
\eqref{eq:ch35-ideal-tropical}.  The example shows that a homogeneous
Fubini--Study presentation is not unique and that the associated ideal is
well defined only up to powers and integral closure; the function
$\psi$ is the same for $\mathfrak b=\mathfrak m_q$ with $m=1$,
provided one checks that $\cO(1)\otimes\mathfrak m_q$ is globally
generated, which it is, by the sections $x$ and $y$.
\end{example}

\section{Homogenization and Legendre slices}
\label{sec:ch35-homogenization}

\begin{definition}[Homogenization]
\label{def:bj-homogenization}
For any $\phi:X^{\An}\to\mathbf R\cup\{-\infty\}$, define
\begin{equation}
\label{eq:bj-homogenization}
                 \widehat\phi(v):=\inf_{t>0}(t\cdot\phi)(v).
\end{equation}
\end{definition}

The function $\widehat\phi$ is the largest homogeneous function below
$\phi$.  Indeed, it is homogeneous by the substitution $t=su$, and
$\widehat\phi\leq\phi$ by taking $t=1$.  If $\psi\leq\phi$ is
homogeneous, then $\psi=t\cdot\psi\leq t\cdot\phi$ for every $t$, hence
$\psi\leq\widehat\phi$.

We record the verification, because the second step is the one place where
the ledger of \Cref{sec:extended-arithmetic} is used without comment
elsewhere.

\begin{lemma}
\label{lem:ch35-homogenization-maximal}
Let $\phi,\phi':X^{\An}\to\mathbf R\cup\{-\infty\}$.  Then:
\begin{enumerate}
\item $\widehat\phi$ is homogeneous and $\widehat\phi\leq\phi$;
\item if $\psi$ is homogeneous and $\psi\leq\phi$, then
 $\psi\leq\widehat\phi$;
\item $\phi\leq\phi'$ implies $\widehat\phi\leq\widehat{\phi'}$, and
 $\widehat{\widehat\phi}=\widehat\phi$;
\item if $(\phi_a)$ is a decreasing net with pointwise limit $\phi$,
 then $\widehat{\phi_a}$ is decreasing with pointwise limit
 $\widehat\phi$;
\item for $s>0$ one has $\widehat{s\cdot\phi}=\widehat\phi$, and
 $\widehat{a\phi}=a\widehat\phi$ for $a>0$.
\end{enumerate}
\par
\end{lemma}

\begin{proof}
(1)  For $s>0$,
\[
 \widehat\phi(sv)=\inf_{t>0}t\,\phi(t^{-1}sv)
 =\inf_{u>0}su\,\phi(u^{-1}v)=s\,\widehat\phi(v),
\]
by the substitution $t=su$ and
\Cref{lem:extended-real-safe}\,(2), which allows the factor $s>0$ to be
pulled out of an infimum in $\Rbar$.  Taking $t=1$ in
\eqref{eq:bj-homogenization} gives $\widehat\phi\leq\phi$.

(2)  If $\psi$ is homogeneous then $t\cdot\psi=\psi$ for every $t>0$
by \Cref{lem:ch35-homogeneity-reformulations}, so
$\psi=t\cdot\psi\leq t\cdot\phi$ by
\Cref{lem:ch35-action-basic}\,(2).  Taking the infimum over $t$ gives
$\psi\leq\widehat\phi$.

(3)  Monotonicity follows from
\Cref{lem:ch35-action-basic}\,(2) and monotonicity of infima.  Idempotence
follows from (1) and (2): $\widehat\phi$ is homogeneous and lies below
$\phi$, so it lies below $\widehat\phi$ by (2), and above by (1).

(4)  Each $\widehat{\phi_a}$ is homogeneous, and the net is decreasing by
(3).  Let $\chi:=\inf_a\widehat{\phi_a}$, an infimum of homogeneous
functions and therefore homogeneous, again by
\Cref{lem:extended-real-safe}\,(2).  Since $\phi\leq\phi_a$ we get
$\widehat\phi\leq\widehat{\phi_a}$ for all $a$, so
$\widehat\phi\leq\chi$.  Conversely
$\chi\leq\widehat{\phi_a}\leq\phi_a$ for all $a$, so $\chi\leq\phi$,
and (2) gives $\chi\leq\widehat\phi$.  Thus $\chi=\widehat\phi$.  This
is \cite[Lemma~6.13]{BJ22}.

(5)  By \Cref{lem:ch35-action-basic}\,(1),
$\widehat{s\cdot\phi}=\inf_{t>0}\bigl(t\cdot(s\cdot\phi)\bigr)
=\inf_{t>0}\bigl((ts)\cdot\phi\bigr)=\widehat\phi$, the last equality
because $t\mapsto ts$ is a bijection of $\mathbf R_{>0}$.  The second
identity is \Cref{lem:ch35-action-basic}\,(3) and
\Cref{lem:extended-real-safe}\,(2).
\end{proof}

At a fixed trivial semivaluation $v_{Y,\triv}$,
\begin{equation}
\label{eq:homogenization-fixed-point}
 \widehat\phi(v_{Y,\triv})=
 \begin{cases}
 0,&\phi(v_{Y,\triv})\geq0,\\
 -\infty,&\phi(v_{Y,\triv})<0.
 \end{cases}
\end{equation}
This is why one checks the supremum before homogenizing.

Indeed, $t^{-1}v_{Y,\triv}=v_{Y,\triv}$, so the quantity being minimized
in \eqref{eq:bj-homogenization} is $t\,\phi(v_{Y,\triv})$; its infimum
over $t>0$ is $0$ if $\phi(v_{Y,\triv})\geq0$ and $-\infty$
otherwise.  This is \cite[Example~6.12]{BJ22}.  Applying it at
$Y=X$ yields the dichotomy that governs the rest of the section: for a
decreasing function, minimality of $v_{\triv}$ in
\eqref{eq:semivaluation-order} gives
$\sup\phi=\phi(v_{\triv})$.  In particular this applies to psh and
generically finite Fubini--Study functions.

\begin{corollary}
\label{cor:ch35-homogenization-alternative}
Let $\phi:X^{\An}\to\mathbf R\cup\{-\infty\}$ be decreasing for the order
\eqref{eq:semivaluation-order}.  Then $\widehat\phi$ is decreasing, and
\[
 \widehat\phi\not\equiv-\infty
 \iff\sup_{X^{\An}}\phi\geq0 .
\]
In particular this applies to every $\phi\in\PSH_X(\theta)$ and to every
$\phi\in\cH^{\mathrm{gf}}_{\mathbf R}(L)$, for which
$\sup\phi=\phi(v_{\triv})$.
\end{corollary}

\begin{proof}
If $v\geq w$ then $t^{-1}v\geq t^{-1}w$ for every $t>0$, because the
order \eqref{eq:semivaluation-order} is tested on ideals and scaling
multiplies all ideal values by $t^{-1}$.  Hence each $t\cdot\phi$ is
decreasing and so is the infimum.

If $\sup\phi<0$, then $\phi(v_{\triv})\leq\sup\phi<0$, and
\eqref{eq:homogenization-fixed-point} gives
$\widehat\phi(v_{\triv})=-\infty$.  Since $\widehat\phi$ is decreasing
and $v_{\triv}$ is the \emph{minimum} of $X^{\An}$ for
\eqref{eq:semivaluation-order} on an integral $X$ --- this is the left
half of \eqref{eq:trivial-endpoints}, together with $v(\mathfrak a)\geq0
=v_{\triv}(\mathfrak a)$ for every nonzero $\mathfrak a$ --- one has
$\widehat\phi\leq\widehat\phi(v_{\triv})$ everywhere, and it follows that
$\widehat\phi\equiv-\infty$.  Conversely, if $\sup\phi\geq0$ then
$\phi(v_{\triv})\geq0$ when $\phi$ is decreasing, whence
$\widehat\phi(v_{\triv})=0$ by
\eqref{eq:homogenization-fixed-point}, so $\widehat\phi\not\equiv-\infty$.
This is \cite[Lemma~6.14]{BJ22}.
\end{proof}

\subsection{The elimination lemma}
\label{sec:ch35-elimination}

The content of \Cref{prop:homogenization-gf} is that the infimum over
$t>0$ of a finite maximum of affine functions is again a finite maximum of
\emph{linear} functions, with rational coefficients when the data are
rational.  We isolate this as a statement about $\mathbf R^r$, where it is
elementary convex geometry, and then transport it.

\begin{lemma}
\label{lem:ch35-elimination}
Let $\lambda_1,\ldots,\lambda_r\in\mathbf R$ with
$\max_i\lambda_i\geq0$, and define
\begin{equation}
\label{eq:ch35-elimination-g}
 g(x):=\inf_{t>0}\max_{1\leq i\leq r}\{x_i+\lambda_it\},
 \qquad x\in\mathbf R^r .
\end{equation}
Then:
\begin{enumerate}
\item $g(x)\geq\min_ix_i$; in particular $g$ is real valued;
\item if $\max_i\lambda_i=0$, then $g(x)=\max_{\lambda_i=0}x_i$; if
 $\min_i\lambda_i\geq0$, then $g(x)=\max_ix_i$;
\item in general there is a finite set $W$ contained in the simplex
\begin{equation}
\label{eq:ch35-elimination-simplex}
 \sigma:=\Bigl\{w\in\mathbf R^r_{\geq0}\ \Bigm|\ \sum_iw_i=1\Bigr\}
\end{equation}
 such that $g(x)=\max_{w\in W}\langle w,x\rangle$ for all
 $x\in\mathbf R^r$;
\item if all $\lambda_i\in\mathbf Q$, then $W$ may be chosen inside
 $\sigma\cap\mathbf Q^r$.
\end{enumerate}
\par
\end{lemma}

\begin{proof}
\emph{(1).}  Choose $i_0$ with $\lambda_{i_0}=\max_i\lambda_i\geq0$.
For every $t>0$,
$\max_i\{x_i+\lambda_it\}\geq x_{i_0}+\lambda_{i_0}t\geq x_{i_0}\geq\min_ix_i$.
Taking the infimum gives (1).  Since also
$g(x)\leq\max_i\{x_i+\lambda_i\}<+\infty$ by taking $t=1$, $g$ is real
valued.

\emph{(2).}  Suppose $\max_i\lambda_i=0$.  For $t>0$ and any $i$ with
$\lambda_i=0$ we get $\max_j\{x_j+\lambda_jt\}\geq x_i$, so
$g(x)\geq\max_{\lambda_i=0}x_i$.  Conversely, if $\lambda_i<0$ then
$x_i+\lambda_it\leq\max_{\lambda_j=0}x_j$ once
$t\geq(x_i-\max_{\lambda_j=0}x_j)/(-\lambda_i)$; taking $t$ larger than
the finitely many such thresholds gives
$\max_j\{x_j+\lambda_jt\}=\max_{\lambda_j=0}x_j$, whence
$g(x)\leq\max_{\lambda_j=0}x_j$.  If $\min_i\lambda_i\geq0$, then
$t\mapsto\max_i\{x_i+\lambda_it\}$ is nondecreasing and its infimum over
$t>0$ is its limit as $t\downarrow0$, namely $\max_ix_i$.

\emph{(3).}  Set
$f(x,t):=\max_i\{x_i+\lambda_it\}$ on
$\mathbf R^r\times\mathbf R_{\geq0}$.  It is convex, being a maximum of
finitely many affine functions, and positively homogeneous of degree one,
because each $x_i+\lambda_it$ is linear in $(x,t)$.  Its epigraph
\[
 \mathcal E:=\{(x,t,s)\in\mathbf R^r\times\mathbf R_{\geq0}\times\mathbf R
 \mid s\geq f(x,t)\}
\]
is a rational-if-the-$\lambda_i$-are polyhedral convex cone, being the
intersection of the $r$ half-spaces $s\geq x_i+\lambda_it$ with
$t\geq0$.  The image of $\mathcal E$ under the projection
$(x,t,s)\mapsto(x,s)$, which forgets $t$, is again a polyhedral convex
cone, by Fourier--Motzkin elimination, and it is rational if
$\mathcal E$ is.  Since $f(x,t)$ is continuous at $t=0$ and
$\max_i\lambda_i\geq0$, the infimum over $t>0$ equals the minimum over
$t\geq0$; hence that image is exactly the epigraph of $g$.  By (1),
$g$ is real valued.  Hence $g$ is a positively homogeneous, convex,
piecewise linear function on $\mathbf R^r$.

A positively homogeneous convex function is the support function of the
closed convex set
\[
 P:=\{w\in\mathbf R^r\mid\langle w,x\rangle\leq g(x)
 \ \text{for all }x\in\mathbf R^r\},
\]
by \Cref{thm:appD-fenchel-moreau} applied to $g$, or directly by
\Cref{ex:appD-indicator-support}: the conjugate of a positively homogeneous
proper convex function is the indicator of $P$, and biconjugation gives
$g=h_P$ in the notation of \eqref{eq:appD-support}.  Since $g$ is
piecewise linear with finitely many pieces, $P$ is a polytope, and a
support function is the maximum over the vertex set $W$ of $P$.

It remains to prove $P\subset\sigma$.  Each $x\mapsto\max_i\{x_i+\lambda_it\}$
is nondecreasing in every coordinate of $x$, hence so is $g$; therefore
$\langle w,\cdot\rangle\leq g$ forces $w\geq0$, since otherwise
$w_j<0$ for some $j$ and letting $x_j\to-\infty$ with the other
coordinates fixed would make $\langle w,x\rangle\to+\infty$ while
$g(x)$ stays bounded above by
$\max_{i\neq j}\{x_i+\lambda_i\}$ whenever $r\geq2$, and the case
$r=1$ is trivial.  Moreover, adding $c\in\mathbf R$ to all coordinates
adds $c$ to $f(x,t)$ and hence to $g$:
$g(x+c\mathbf 1)=g(x)+c$.  If $w\in P$, then for all $x$ and all
$c$,
\[
 \langle w,x\rangle+c\sum_iw_i=\langle w,x+c\mathbf 1\rangle
 \leq g(x)+c ,
\]
so $c(\sum_iw_i-1)\leq g(x)-\langle w,x\rangle$ for every
$c\in\mathbf R$; letting $c\to\pm\infty$ forces $\sum_iw_i=1$.  Thus
$P\subset\sigma$ and $W\subset\sigma$.

\emph{(4).}  If all $\lambda_i$ are rational, then $\mathcal E$ is a
rational polyhedral cone, so its projection is rational and $g$ is a
rational piecewise linear function; then $P$ is a rational polytope and
its vertices lie in $\mathbf Q^r$.
\end{proof}

\begin{remark}[Reading the elimination lemma tropically]
\label{rem:ch35-elimination-tropical}
Formula \eqref{eq:ch35-elimination-g} is a max-plus expression, and the
lemma says that eliminating the variable $t$ from a max-plus polynomial
with a nonnegative exponent produces a max-plus polynomial with the same
variables and exponents in the simplex $\sigma$.  The exponents
$w\in W$ are precisely the barycentric weights that occur when the
inhomogeneous system $(s_i,\lambda_i)$ is replaced by the homogeneous
system of monomials $\prod_is_i^{bw_i}$; this is the substance of the
proof of \Cref{prop:homogenization-gf} below and it is the reason that
homogenization of a Fubini--Study function does not leave the algebraic
category.
\end{remark}

\begin{proposition}\label{prop:homogenization-gf}
If $\phi\in\cH^{\mathrm{gf}}_{\mathbf Q}(L)$, then
$\widehat\phi\in\cH_{\mathrm{hom}}(L)$ when
$\sup\phi\geq0$, while $\widehat\phi\equiv-\infty$ when
$\sup\phi<0$.  Consequently every homogeneous psh function has a
decreasing homogeneous GF regularization.
\end{proposition}

\begin{proof}[Polyhedral elimination]
Write $m\phi=\max_i\{\log|s_i|+\lambda_i\}$.  Formula
\eqref{eq:scaling-fs-weights} gives
\[
 m\widehat\phi
  =\inf_{t>0}\max_i\{\log|s_i|+t\lambda_i\}.
\]
The right side is a rational, convex, homogeneous piecewise linear function
of the finitely many variables $\log|s_i|$; it is therefore a finite
maximum of rational convex combinations
$\sum_i a_i\log|s_i|$.  Clearing denominators turns each such combination
into the logarithm of a monomial in the $s_i$, giving
\eqref{eq:homogeneous-fs-presentation}; terms with zero coefficients
are omitted. If all relevant weights are
negative, the infimum is $-\infty$.  Homogenization commutes with
decreasing nets by its maximal characterization, which proves the last
assertion.  This is \cite[Theorem~6.15 and Corollary~6.17]{BJ22}.

We now carry out the four steps in detail.

\emph{Step 1: the alternative.}  A generically finite Fubini--Study
function is decreasing, by the computation in the proof of
\Cref{prop:ch35-homogeneous-normalization}, so
\Cref{cor:ch35-homogenization-alternative} applies: if
$\sup\phi<0$ then $\widehat\phi\equiv-\infty$, and there is nothing more
to prove.  Assume therefore $\sup\phi\geq0$.  Fix a presentation
\begin{equation}
\label{eq:ch35-gf-presentation-fixed}
 \phi=\frac1m\max_{1\leq i\leq r}\{\log|s_i|+\lambda_i\},
 \qquad 0\neq s_i\in H^0(X,mL),\ \lambda_i\in\mathbf Q,
\end{equation}
as in \eqref{eq:gf-fs-presentation}.  Since
$\sup\phi=\frac1m\max_i\lambda_i$, the hypothesis of
\Cref{lem:ch35-elimination} is met: $\max_i\lambda_i\geq0$.

\emph{Step 2: elimination.}  By \eqref{eq:scaling-fs-weights},
\[
 m\,\widehat\phi(v)=\inf_{t>0}\max_i\{\log|s_i|(v)+t\lambda_i\}.
\]
Suppose first that $v\in X^{\An}$ is such that all the numbers
$x_i:=\log|s_i|(v)$ are finite.  Then the displayed quantity is
$g(x)$ for the function $g$ of \eqref{eq:ch35-elimination-g}, and
\Cref{lem:ch35-elimination}\,(3)--(4) supply a finite set
$W\subset\sigma\cap\mathbf Q^r$ with
\begin{equation}
\label{eq:ch35-elimination-applied}
 m\,\widehat\phi(v)=\max_{w\in W}\sum_{i:w_i>0}w_i\log|s_i|(v).
\end{equation}
If some $x_i=-\infty$, the same identity holds with both sides read in
$\mathbf R\cup\{-\infty\}$.  Indeed, replace every such coordinate by
$-N$.  For each $N$ the finite-valued identity above holds.  As
$N\to+\infty$, both sides decrease: on the right, the finite maximum
loses precisely the terms whose support meets
$\{i:x_i=-\infty\}$; on the left,
\[
 \inf_N\inf_{t>0}f_N(t)=\inf_{t>0}\inf_N f_N(t),
\]
because both sides are the infimum over the product index set.  The latter
is exactly the extended-real expression in
\eqref{eq:ch35-elimination-applied}.  Thus the same identity holds at the
boundary faces as well.

\emph{Step 3: monomials.}  Choose $b\in\mathbf Z_{>0}$ with
$bw\in\mathbf N^r$ for every $w\in W$, possible because $W$ is finite
and rational.  For $w\in W$ put
\[
 s_w:=\prod_{i=1}^rs_i^{bw_i}\in H^0(X,mbL),
\]
which is nonzero because $X$ is integral and each $s_i\neq0$, and which
has degree $\sum_ibw_i=b$ in the $s_i$, matching the line bundle
$mbL$ since $\sum_iw_i=1$.  Multiplicativity of $v\mapsto v(\cdot)$
gives $\log|s_w|=\sum_{i:w_i>0}bw_i\log|s_i|$, so
\eqref{eq:ch35-elimination-applied} becomes
\[
 \widehat\phi=\frac1{mb}\max_{w\in W}\log|s_w| .
\]
This is a presentation \eqref{eq:homogeneous-fs-presentation} with all
weights zero, and it is generically finite because each $s_w\neq0$.
Hence $\widehat\phi\in\cH^{\mathrm{gf}}_0(L)=\cH_{\mathrm{hom}}(L)$.

\emph{Step 4: the psh consequence.}  Let
$\psi\in\PSH_{\mathrm{hom},X}(\theta)$, and let
$\phi_a\searrow\psi$ be a decreasing net with
$\phi_a\in\cH^{\mathrm{gf}}_{\mathbf Q}(L_a)$ and
$c_1(L_a)\to\theta$, as furnished by \Cref{def:bj-psh} and
\cite[Theorem~4.15]{BJ22}, which allows the approximants to be taken with
rational weights.  Then $\sup\phi_a\geq\sup\psi=0$ by
\Cref{prop:ch35-homogeneous-normalization}, so
$\widehat{\phi_a}\in\cH_{\mathrm{hom}}(L_a)$ by Steps 1--3.  By
\Cref{lem:ch35-homogenization-maximal}\,(4), the net
$\widehat{\phi_a}$ is decreasing with pointwise limit
$\widehat\psi=\psi$, the last equality by
\Cref{lem:ch35-homogenization-maximal}\,(3) since $\psi$ is homogeneous.
Thus $\psi$ has a decreasing regularization by homogeneous Fubini--Study
functions in classes tending to $\theta$.  If $\theta=c_1(L)$ with
$L$ ample, \cite[Theorem~4.15\,(iii)]{BJ22} allows $L_a=L$ for all
$a$.
\end{proof}

\begin{corollary}
\label{cor:ch35-homogenization-psh}
Let $\theta\in N^1(X)$ and $\phi\in\PSH_X(\theta)$ with
$\sup\phi\geq0$.  Then
$\widehat\phi\in\PSH_{\mathrm{hom},X}(\theta)$, and $\widehat\phi$ is
the pointwise limit of a decreasing net $\psi_a\in\cH_{\mathrm{hom}}(L_a)$
with $c_1(L_a)\to\theta$.  If $\theta$ is the class of an ample
$\mathbf Q$-line bundle $L$, one may take $L_a=L$ throughout.
\end{corollary}

\begin{proof}
Choose a decreasing regularization $\phi_a\searrow\phi$ with
$\phi_a\in\cH^{\mathrm{gf}}_{\mathbf Q}(L_a)$ and
$c_1(L_a)\to\theta$, using \cite[Theorem~4.15]{BJ22} to make the weights
rational; in the ample case the same theorem gives $L_a=L$.  Then
$\sup\phi_a\geq\sup\phi\geq0$, so
$\psi_a:=\widehat{\phi_a}\in\cH_{\mathrm{hom}}(L_a)$ by
\Cref{prop:homogenization-gf}, and $\psi_a\searrow\widehat\phi$ by
\Cref{lem:ch35-homogenization-maximal}\,(4).  The limit is finite at
$v_{\triv}$, being $0$ by
\eqref{eq:homogenization-fixed-point}, hence generically finite; so
\Cref{thm:bj-basic-closure}\,(5) gives
$\widehat\phi\in\PSH_X(\theta)$, and it is homogeneous by
\Cref{lem:ch35-homogenization-maximal}\,(1).  This is
\cite[Corollary~6.17]{BJ22}.
\end{proof}

\begin{warning}{The hypothesis $\sup\phi\geq0$ is not cosmetic}
If $\sup\phi<0$, then $\widehat\phi\equiv-\infty$, which is not a psh
function: \Cref{def:bj-psh} excludes it by generic finiteness.  Since
$\sup(\phi+c)=\sup\phi+c$, the hypothesis can always be arranged by
adding a constant, but the resulting homogenization then depends on the
constant: by \Cref{lem:ch35-action-basic}\,(3),
$\widehat{\phi+c}=\widehat\phi_{-c}$ in the notation
\eqref{eq:homogeneous-legendre-slice} below, and these are genuinely
different homogeneous functions.  The one-parameter family of Legendre
slices exists precisely to keep track of that dependence.
\end{warning}

\begin{example}[Homogenizing a Fubini--Study function on $\mathbf P^1$]
\label{ex:ch35-homogenize-p1}
Continue \Cref{ex:ch35-p1-orbit} with
$\phi=\max\{\log|x|+a,\log|y|+b\}$ on $X=\mathbf P^1$,
$L=\cO(1)$, and $a,b\in\mathbf Q$.  Then $m=1$,
$r=2$, $(\lambda_1,\lambda_2)=(a,b)$, and
$\sup\phi=\max\{a,b\}$.

If $a=b$, then $\phi=\max\{\log|x|,\log|y|\}+a\equiv a$, since $x,y$
have no common zero.  Suppose first $a\geq0$.  Then
$\min_i\lambda_i=a\geq0$ and \Cref{lem:ch35-elimination}\,(2) applies
\emph{directly}, with no subtraction, giving
$\widehat\phi=\max\{\log|x|,\log|y|\}=0$; the direct computation agrees,
$\widehat\phi(v)=\inf_{t>0}t\,\phi(t^{-1}v)=\inf_{t>0}ta=0$.  If instead
$a<0$, then $\sup\phi=a<0$ and \Cref{prop:homogenization-gf} gives
$\widehat\phi\equiv-\infty$; again directly,
$\inf_{t>0}ta=-\infty$.  So the homogenization of a Fubini--Study
function with equal weights $a$ is the zero function when $a\geq0$ and
$-\infty$ when $a<0$.  One may \emph{not} normalize by subtracting
$a$ first: homogenization does not commute with the addition of a
constant, by the warning preceding this example.

If $a>b\geq0$, then $\min_i\lambda_i\geq0$.  Hence
\Cref{lem:ch35-elimination}\,(2) gives
\[
 \widehat\phi=\max\{\log|x|,\log|y|\}=0
\]
again.

The interesting case is $a>0>b$.  Here
$g(x_1,x_2)=\inf_{t>0}\max\{x_1+at,\ x_2+bt\}$.  The two affine functions
of $t$ cross at $t^*=(x_2-x_1)/(a-b)$, which is positive exactly when
$x_2>x_1$; since $a>0>b$, the maximum is decreasing to the left of the
crossing and increasing to its right, so
\[
 g(x)=
 \begin{cases}
  \dfrac{-b\,x_1+a\,x_2}{a-b},&x_2>x_1,\\[2mm]
  x_1,&x_2\leq x_1 .
 \end{cases}
\]
Both branches are linear with nonnegative coefficients summing to one, and
$g(x)=\max\bigl\{x_1,\ \tfrac{-b}{a-b}x_1+\tfrac{a}{a-b}x_2\bigr\}$, so
$W=\{(1,0),(\tfrac{-b}{a-b},\tfrac{a}{a-b})\}\subset\sigma\cap\mathbf Q^2$,
confirming \Cref{lem:ch35-elimination}\,(3)--(4).  Write
$\tfrac{-b}{a-b}=p/q$ in lowest terms with integers $0<p<q$.
Then $\tfrac a{a-b}=(q-p)/q$, and Step 3 of the proof of
\Cref{prop:homogenization-gf} produces the monomial sections
$s_1=x^q$ and $s_2=x^{p}y^{q-p}$ in $H^0(X,\cO(q))$ and
\[
 \widehat\phi=\tfrac1q\max\{\log|x^q|,\ \log|x^py^{q-p}|\}
 =\max\bigl\{\log|x|,\ \tfrac pq\log|x|+\tfrac{q-p}q\log|y|\bigr\}.
\]
Evaluating on the ray $s\,\ord_{[0:1]}$, where $\log|x|=-s$ and
$\log|y|=0$, gives $\widehat\phi(s\,\ord_{[0:1]})=-\tfrac pqs$, a linear
function of $s$ with slope $-p/q=b/(a-b)$, whereas
$\phi(s\,\ord_{[0:1]})=\max\{-s+a,b\}$ is only piecewise affine.  The
homogenization has replaced the broken graph by the linear function through
the origin that lies below it and has the largest possible slope; this is
exactly the geometric content of
\Cref{lem:ch35-homogenization-maximal}\,(2).
\end{example}

\section{Determination by divisorial points}
\label{sec:ch35-determination}

Everything that follows --- the Legendre calculus, the b-divisor
correspondence, and the reflection arguments of
\Cref{chap:bj-reflection} --- rests on the ability to promote an inequality
verified on the divisorial locus to an inequality on the whole
analytification.  We prove that promotion here, isolate its single cited
input, and show by example that it fails without plurisubharmonicity.

\subsection{Why density is not enough}
\label{sec:ch35-density-insufficient}

\Cref{thm:triv-divisorial-density} says that $X^{\divis}$ is dense in
$X^{\An}$.  Density gives the following implication, and no more.

\begin{lemma}
\label{lem:ch35-density-gives}
Let $D\subset Z$ be dense in a topological space $Z$, let
$f:Z\to\Rbar$ be lower semicontinuous and $g:Z\to\Rbar$ upper
semicontinuous.  If $f\leq g$ on $D$, then $f\leq g$ on $Z$.
\end{lemma}

\begin{proof}
Fix $z\in Z$ and a net $(z_a)$ in $D$ with $z_a\to z$, which exists
by density.  Then
\[
 f(z)\leq\liminf_af(z_a)\leq\liminf_ag(z_a)
 \leq\limsup_ag(z_a)\leq g(z),
\]
by \Cref{prop:net-semicontinuity} applied to $f$ and to $g$.
\end{proof}

The hypothesis of \Cref{lem:ch35-density-gives} is the wrong way round for
our purposes.  A psh function is upper semicontinuous, not lower
semicontinuous, and we wish to compare it \emph{from below} against another
upper-semicontinuous function.  The following example shows that the
implication then fails outright, so
\Cref{thm:ch35-divisorial-comparison} below is not a corollary of density.

\begin{example}[An usc function not determined on $X^{\divis}$]
\label{ex:ch35-usc-not-determined}
Let $X=\mathbf P^1_k$ and fix a closed point $p$.  Recall from
\Cref{ex:triv-p1-tree} that $v_{p,\triv}$, the trivial semivaluation with
support $p$, is the endpoint of the ray through $p$ and does not belong
to $X^{\divis}$.  Define
\[
 f:=\one_{\{v_{p,\triv}\}},
 \qquad
 g:=0 .
\]
Then $f$ is upper semicontinuous: the set $\{f\geq c\}$ is
$\{v_{p,\triv}\}$ for $0<c\leq1$, which is closed because
$X^{\An}$ is Hausdorff by \Cref{thm:triv-berkovich-compactness}, it is
$X^{\An}$ for $c\leq0$, and it is empty for $c>1$.  Both $f$ and
$g$ are upper semicontinuous, and $f\leq g$ on
$X^{\divis}$, indeed $f=0=g$ there.  Yet $f(v_{p,\triv})=1>0=g(v_{p,\triv})$.

Thus two upper-semicontinuous functions agreeing on the dense set
$X^{\divis}$ need not agree on $X^{\An}$.  What fails for $f$ is
plurisubharmonicity: a $\theta$-psh function is decreasing and hence
attains its supremum at $v_{\triv}$, by
\cite[Proposition~4.12]{BJ22}, whereas
$\sup f=1>0=f(v_{\triv})$.  Equivalently, $f$ is not decreasing for the
order \eqref{eq:semivaluation-order}: by \eqref{eq:trivial-endpoints} the
point $v_{\triv}$ is the \emph{minimum} of $X^{\An}$, so
$v_{p,\triv}\geq v_{\triv}$, whereas
$f(v_{p,\triv})=1>0=f(v_{\triv})$.
\end{example}

\subsection{The cited input}
\label{sec:ch35-cited-input}

One statement is imported.  It is proved in
\Cref{chap:bj-test-configurations} at the level of test configurations, and
it was already used in the proof of
\Cref{thm:triv-divisorial-density}; we restate it here with its exact
hypotheses so that the logical dependence of this section is visible.

\begin{definition}[Rees valuations of a flag ideal]
\label{def:ch35-rees-valuations}
Let $\mathfrak a$ be a flag ideal on $X_{\triv}=X\times\mathbf A^1$, in
the sense of \Cref{def:bj-flag-ideal}, and let
$\cX\to X_{\triv}$ be the relative integral closure of the blowup of $\mathfrak a$
inside the unchanged generic product, an integrally closed test configuration.
Full normalization is equivalent only when $X$ is normal.  The finite set
\begin{equation}
\label{eq:ch35-rees-set}
 \Sigma_{\mathfrak a}
 :=\{v_E\mid E\ \text{an irreducible component of}\ \cX_0\}
 \subset X^{\divis}
\end{equation}
of divisorial points obtained from the components of the central fibre by
Gauss extension \eqref{eq:triv-gauss-section} is the set of \emph{Rees
valuations} of $\mathfrak a$.
\end{definition}

\begin{proposition}
\label{prop:ch35-effectivity}
Let $\cX$ be an integrally closed test configuration for $X$ and let
$D$ be a vertical $\mathbf Q$-Cartier divisor on $\cX$.  Then
$D\geq0$ if and only if $\phi_D(v_E)\geq0$ for every irreducible
component $E$ of $\cX_0$.
\end{proposition}

\begin{proof}
This is \cite[Lemmas~1.23 and~2.8]{BJ22}; the argument is the one recalled
in the proof of \Cref{thm:triv-divisorial-density}, namely that
$\sigma_{\cX}(v_E)=\ord_E(\cX_0)^{-1}\ord_E$, so that nonnegativity of
$\phi_D$ at the $v_E$ is exactly the containment
$\cO_{\cX}(-D)\subset\cO_{\cX}$, which for an integrally closed
$\cX$ is equivalent to effectivity of $D$.  Relative integral closedness is used and cannot be dropped. The criterion
is the vertical Serre regularity statement of \Cref{chap:bj-test-configurations};
it permits a nonnormal generic fibre and is not an assertion that the
whole test configuration must be normal.
\end{proof}

\subsection{Attainment of suprema on a finite set}
\label{sec:ch35-attainment}

\begin{lemma}
\label{lem:ch35-finite-attainment}
Let $\mathfrak a$ be a flag ideal on $X_{\triv}$ with associated PL
function $\phi_{\mathfrak a}$, and let
$\Sigma_{\mathfrak a}\subset X^{\divis}$ be its set of Rees valuations.
Then
\begin{equation}
\label{eq:ch35-finite-attainment}
 \sup_{X^{\An}}(\chi-\phi_{\mathfrak a})
 =\max_{\Sigma_{\mathfrak a}}(\chi-\phi_{\mathfrak a})
\end{equation}
for every $\chi\in\PL^+(X)$.  Consequently, for every
$\eta\in\PL(X)$ there is a finite set $\Sigma\subset X^{\divis}$ such
that
\begin{equation}
\label{eq:ch35-finite-attainment-general}
 \sup_{X^{\An}}(\chi-\eta)=\max_\Sigma(\chi-\eta)
 \qquad\text{for every }\chi\in\PL^+(X).
\end{equation}
\end{lemma}

\begin{proof}
\emph{Step 1: reduction.}  Write $\chi=m^{-1}\phi_{\mathfrak a'}$ for a
flag ideal $\mathfrak a'$ and $m\in\mathbf Z_{>0}$, which is possible by
the definition \eqref{eq:flag-ideal-function} of $\PL^+(X)$.  Then
$\chi-\phi_{\mathfrak a}=m^{-1}(\phi_{\mathfrak a'}-\phi_{\mathfrak a^m})$
by the semiring identities \eqref{eq:flag-semiring}, and the Rees valuations
of $\mathfrak a^m$ coincide with those of $\mathfrak a$, because their blowups agree and so do their relative integral closures.  Since
$m^{-1}>0$ commutes with suprema by
\Cref{lem:safe-supremum}, we may assume $m=1$.

\emph{Step 2: the value of the maximum.} Let $\cX$ be the relative
integral closure of the blowup of $\mathfrak a$, with central components
$(E_i)$. Write $\mathfrak a\cO_{\cX}=\cO_{\cX}(-E)$, where $E$ is a
vertical Cartier divisor, possibly of either sign for a fractional flag
ideal. Thus $\sigma(v)(E)=\sigma(v)(\mathfrak a)$ and
\[
 c:=\max_{\Sigma_{\mathfrak a}}(\phi_{\mathfrak a'}-\phi_{\mathfrak a})
 =\max_i\frac{\ord_{E_i}(E)-\ord_{E_i}(\mathfrak a')}
                    {\ord_{E_i}(\cX_0)}\in\mathbf Q.
\]

\emph{Step 3: ideal containment.} Choose an integer $r>0$ with
$rc\in\mathbf Z$. The coherent fractional ideal
\[
 \mathfrak c:=(\mathfrak a')^r\cO_{\cX}(rE-rc\,\cX_0)
\]
has nonnegative order at every $E_i$, by the definition of $c$, and is
regular away from the central fibre. Apply the vertical Serre regularity
criterion of \cite[Lemma~1.23]{BJ22} locally to each generator, as recalled
in \Cref{chap:bj-test-configurations}. This gives
$\mathfrak c\subseteq\cO_{\cX}$. The Cartier-divisor effectivity
criterion alone would not establish containment of a noninvertible ideal.
Evaluating the containment through $\sigma(v)$ gives
\[
 r\,\sigma(v)(\mathfrak a')-r\,\sigma(v)(\mathfrak a)+rc\geq0,
\]
since $\sigma(v)(\cX_0)=1$. Equivalently
$\phi_{\mathfrak a'}(v)-\phi_{\mathfrak a}(v)\leq c$ for every $v$.
The flag functions are finite, so these differences are defined.
Equality of the supremum with $c$ follows because the finite Rees set
is contained in $X^{\An}$. This proves
\eqref{eq:ch35-finite-attainment}; compare \cite[Lemma~2.13]{BJ22}.

\emph{Step 4: general $\eta\in\PL(X)$.}  By
\Cref{def:bj-pl-function} and the semiring identities
\eqref{eq:flag-semiring}, we may write
$\eta=r(\phi_{\mathfrak a}-\rho)$ with $r\in\mathbf Q_{>0}$,
$\mathfrak a$ a flag ideal, and $\rho\in\PL^+(X)$.  Then, for
$\chi\in\PL^+(X)$,
\[
 \chi-\eta=r\bigl((r^{-1}\chi+\rho)-\phi_{\mathfrak a}\bigr),
\]
and $r^{-1}\chi+\rho\in\PL^+(X)$ because $\PL^+(X)$ is stable under
sums and under multiplication by positive rationals.  Applying
\eqref{eq:ch35-finite-attainment} with $\Sigma:=\Sigma_{\mathfrak a}$ and
using \Cref{lem:safe-supremum} once more gives
\eqref{eq:ch35-finite-attainment-general}.
\end{proof}

\begin{lemma}
\label{lem:ch35-attainment-psh}
Let $\eta\in\PL(X)$ and let $\Sigma\subset X^{\divis}$ be as in
\Cref{lem:ch35-finite-attainment}.  Then
\begin{equation}
\label{eq:ch35-attainment-psh}
 \sup_{X^{\An}}(\phi-\eta)=\max_\Sigma(\phi-\eta)
\end{equation}
for every $\theta\in N^1(X)$ and every $\phi\in\PSH_X(\theta)$.
\end{lemma}

\begin{proof}
Choose $\omega\in\Amp(X)$ with $\omega-\theta$ nef.  Then
$0\in\PSH_X(\omega-\theta)$ by \eqref{eq:psh-existence-cones}, and
\Cref{thm:bj-basic-closure}\,(3) gives
$\PSH_X(\theta)\subseteq\PSH_X(\omega)$; so we may assume
$\theta=\omega$ is ample.  By \cite[Theorem~4.15\,(iii)]{BJ22} there is a
decreasing net $(\phi_a)$ in
$\cH_{\mathrm{dom}}(\omega)\subset\PL(X)\cap\PSH_X(\omega)$ with
$\phi_a\searrow\phi$ pointwise, and
$\cH_{\mathrm{dom}}(\omega)\subset\PL^+(X)$ by
\cite[(3.2)]{BJ22}.

Fix $v\in X^{\An}$.  For each $a$,
\Cref{lem:ch35-finite-attainment} gives
$\phi_a(v)-\eta(v)\leq\max_\Sigma(\phi_a-\eta)$.  Since
$\phi_a\searrow\phi$ pointwise and $\Sigma$ is finite,
$\max_\Sigma(\phi_a-\eta)\searrow\max_\Sigma(\phi-\eta)$; the interchange
of a decreasing limit with a maximum over a finite set is legitimate in
$\Rbar$ by \Cref{lem:extended-real-safe}\,(1) applied finitely often.
Letting $a$ increase gives
$\phi(v)-\eta(v)\leq\max_\Sigma(\phi-\eta)$, where the difference is
defined because $\eta$ is real valued.  Taking the supremum over
$v$ proves $\leq$ in \eqref{eq:ch35-attainment-psh}, and $\geq$ is
trivial.
\end{proof}

\subsection{The comparison theorem}
\label{sec:ch35-comparison}

\begin{lemma}
\label{lem:ch35-usc-inf-continuous}
Let $Z$ be a compact Hausdorff space and let
$h:Z\to[-\infty,C]$ be upper semicontinuous, $C\in\mathbf R$.  Then the
family
\[
 \mathcal G:=\{g\in C^0(Z)\mid g\geq h\}
\]
is nonempty, downward directed, and satisfies
$\inf_{g\in\mathcal G}g=h$ pointwise.  In particular $h$ is the limit of
a decreasing net of continuous functions.
\end{lemma}

\begin{proof}
The constant $C$ belongs to $\mathcal G$, which is therefore nonempty,
and $\min\{g_1,g_2\}\in\mathcal G$ for $g_1,g_2\in\mathcal G$, so
$\mathcal G$ is downward directed.  Since every $g\in\mathcal G$ satisfies
$g\geq h$, one has $\inf\mathcal G\geq h$.

For the reverse inequality fix $z\in Z$ and $c>h(z)$; we may assume
$c\leq C$, replacing $c$ by $\min\{c,C\}$ if necessary, which is still
$>h(z)$ unless $h(z)=C$, in which case there is nothing to prove.  By
upper semicontinuity there is an open $U\ni z$ with $h<c$ on $U$.  A
compact Hausdorff space is normal, so Urysohn's lemma provides a continuous
$u:Z\to[0,1]$ with $u(z)=1$ and $u=0$ on $Z\setminus U$.  Put
\[
 g:=C+(c-C)u\in C^0(Z).
\]
On $Z\setminus U$ we have $g=C\geq h$.  On $U$, $u\in[0,1]$ and
$c\leq C$ give $g\in[c,C]$, so $g\geq c>h$.  Hence
$g\in\mathcal G$, and $g(z)=c$.  Letting $c\downarrow h(z)$ gives
$\inf_{\mathcal G}g\leq h(z)$.
\end{proof}

\begin{theorem}
\label{thm:ch35-divisorial-comparison}
Let $\theta\in N^1(X)$, let $\phi\in\PSH_X(\theta)$, and let
$\eta:X^{\An}\to\mathbf R\cup\{-\infty\}$ be upper semicontinuous.  Then
\begin{equation}
\label{eq:ch35-divisorial-comparison}
 \phi\leq\eta\ \text{on }X^{\divis}
 \iff
 \phi\leq\eta\ \text{on }X^{\An}.
\end{equation}
Equivalently, $\phi$ is the smallest upper-semicontinuous extension of
$\phi|_{X^{\divis}}$ to $X^{\An}$.
\end{theorem}

\begin{proof}
The implication $\Leftarrow$ is trivial.  We prove $\Rightarrow$ in
three steps of increasing generality.

\emph{Step 1: $\eta\in\PL(X)$.}  Let $\Sigma\subset X^{\divis}$ be the
finite set provided by \Cref{lem:ch35-finite-attainment} for $\eta$.  By
\Cref{lem:ch35-attainment-psh},
\[
 \sup_{X^{\An}}(\phi-\eta)=\max_\Sigma(\phi-\eta)\leq0,
\]
the last inequality because $\Sigma\subset X^{\divis}$ and
$\phi\leq\eta$ there.  Hence $\phi\leq\eta$ on $X^{\An}$.

\emph{Step 2: $\eta\in C^0(X^{\An})$.}  By
\eqref{eq:pl-stone-weierstrass}, $\PL(X)$ is uniformly dense in
$C^0(X^{\An})$, so we may pick
$\eta_j\in\PL(X)$ with $\|\eta_j-\eta\|_\infty\leq1/j$.  Then
$\phi\leq\eta\leq\eta_j+1/j$ on $X^{\divis}$, and
$\eta_j+1/j\in\PL(X)$; Step 1 gives
$\phi\leq\eta_j+1/j$ on $X^{\An}$.  Letting $j\to\infty$ gives
$\phi\leq\eta$ on $X^{\An}$.

\emph{Step 3: $\eta$ upper semicontinuous.}  The function $\eta$ is
bounded above on the compact space $X^{\An}$ by
\Cref{cor:usc-attains-max}; if $\eta\equiv-\infty$ then
the assumed divisorial inequality contradicts the finiteness of $\phi$ on
$X^{\divis}$, so this case does not occur.  By
\Cref{lem:ch35-usc-inf-continuous} there is a decreasing net
$(g_a)$ in $C^0(X^{\An})$ with $g_a\searrow\eta$ pointwise.  For each
$a$ we have $\phi\leq\eta\leq g_a$ on $X^{\divis}$, so Step 2 gives
$\phi\leq g_a$ on $X^{\An}$.  Taking the infimum over $a$ gives
$\phi\leq\eta$ on $X^{\An}$.

For the final sentence, let $\eta$ be any upper-semicontinuous function on
$X^{\An}$ with $\eta=\phi$ on $X^{\divis}$.  Then
$\phi\leq\eta$ on $X^{\An}$ by the theorem, so $\phi$ is the smallest
such extension.  This is \cite[Theorem~4.22]{BJ22}.
\end{proof}

\begin{corollary}
\label{cor:ch35-divisorial-determination}
Let $\theta,\theta'\in N^1(X)$, $\phi\in\PSH_X(\theta)$,
$\phi'\in\PSH_X(\theta')$.  Then:
\begin{enumerate}
\item if $\phi\leq\phi'$ on $X^{\divis}$, then $\phi\leq\phi'$ on
 $X^{\An}$;
\item if $\phi=\phi'$ on $X^{\divis}$, then $\phi=\phi'$ on
 $X^{\An}$;
\item the topology on $\PSH_X(\theta)$ of pointwise convergence on
 $X^{\divis}$ is Hausdorff;
\item if $(\phi_a)$ is a net in $\PSH_X(\theta)$ which is decreasing on
 $X^{\divis}$, then it is decreasing on $X^{\An}$, and if its pointwise
 limit on $X^{\divis}$ is the restriction of some
 $\phi\in\PSH_X(\theta)$, then $\phi_a\to\phi$ pointwise on
 $X^{\An}$.
\end{enumerate}
\par
\end{corollary}

\begin{proof}
(1)  A $\theta'$-psh function is upper semicontinuous by
\Cref{def:bj-psh}, so \Cref{thm:ch35-divisorial-comparison} applies with
$\eta=\phi'$.

(2)  Apply (1) twice.

(3)  Two distinct elements of $\PSH_X(\theta)$ differ at some point of
$X^{\divis}$ by (2), and the evaluation maps at those points separate
them; a topology of pointwise convergence for a separating family of maps
into the Hausdorff space $\Rbar$ is Hausdorff.

(4)  Monotonicity of the net on $X^{\An}$ is (1) applied to each pair
$a\leq a'$.  Let $\chi$ be the pointwise limit on $X^{\An}$; it is
upper semicontinuous by \Cref{lem:usc-inf}, agrees with $\phi$ on
$X^{\divis}$, but neither global inequality follows formally from
density alone.  Precisely: $\chi\leq\phi_a$ for all $a$, and
$\chi=\phi$ on $X^{\divis}$, so $\phi\leq\chi$ on $X^{\An}$ by
\Cref{thm:ch35-divisorial-comparison} applied with $\eta=\chi$; on the
other hand $\chi$ is generically finite because
$\chi(v_{\triv})=\phi(v_{\triv})$ is finite, so
\Cref{thm:bj-basic-closure}\,(4) gives $\chi\in\PSH_X(\theta)$ and
$\chi\leq\phi$ on $X^{\An}$ by
\Cref{thm:ch35-divisorial-comparison} with the roles exchanged.  Hence
$\chi=\phi$.  This is \cite[Corollaries~4.23 and~4.25]{BJ22}.
\end{proof}

\begin{corollary}
\label{cor:ch35-zero-class}
$\PSH_{\mathrm{hom},X}(0)=\{0\}$.
\end{corollary}

\begin{proof}
The function $0$ is $0$-psh and homogeneous, so
$\{0\}\subseteq\PSH_{\mathrm{hom},X}(0)$.  Conversely let
$\psi\in\PSH_{\mathrm{hom},X}(0)$, so that $\psi\leq0$ and
$\psi(v_{\triv})=0$ by
\Cref{prop:ch35-homogeneous-normalization}.  For every integer
$m\geq1$, \Cref{thm:bj-basic-closure}\,(3) gives
$m\psi\in\PSH_X(m\cdot0)=\PSH_X(0)$, and $m\psi$ is homogeneous.  The
sequence $(m\psi)_m$ is decreasing, because $\psi\leq0$, and its
pointwise limit $\chi:=\inf_m m\psi$ satisfies
$\chi(v_{\triv})=0$, hence is generically finite; so
$\chi\in\PSH_X(0)$ by \Cref{thm:bj-basic-closure}\,(4).

By \Cref{thm:bj-basic-closure}\,(7), $\chi$ is finite at every point of
$X^{\divis}$.  But $\chi(v)=-\infty$ at any $v$ with
$\psi(v)<0$.  Hence $\psi=0$ on $X^{\divis}$, and
\Cref{cor:ch35-divisorial-determination}\,(2) gives $\psi=0$ on
$X^{\An}$.  Compare \cite[Corollary~4.24]{BJ22}, which proves the stronger
statement $\PSH_X(0)=\mathbf R$ for connected $X$.
\end{proof}

\begin{warning}{Determination is not continuity of evaluation}
\Cref{cor:ch35-divisorial-determination} says that the restriction map
$\PSH_X(\theta)\to\Rbar^{X^{\divis}}$ is injective.  It does not say that
the inverse is continuous, and it does not say that convergence on
$X^{\divis}$ implies convergence at every analytic point.  The distinction
matters at proper-support points, where a limit of psh functions may drop.
\end{warning}

\begin{example}[Evaluation is not continuous]
\label{ex:ch35-evaluation-discontinuous}
Let $X=\mathbf P^1_k$, $L=\cO(1)$ with homogeneous coordinates
$x,y$, and $p=[0:1]$, so that $p$ is the zero of $x$.  For
$m\geq1$ set
\[
 \phi_m:=\max\bigl\{\tfrac1m\log|x|,\ -1\bigr\}
 \in\cH_{\mathbf Q}\bigl(\tfrac1mL\bigr)
 \subset\PSH_X\bigl(\tfrac1m c_1(L)\bigr)
 \subset\PSH_X(c_1(L)),
\]
the last inclusion by \eqref{eq:psh-existence-cones} because
$(1-\tfrac1m)c_1(L)$ is nef and $0$ is psh for it.
For the asserted FS membership, the two sections x,y of L give
$\phi_m=m^{-1}\max\{\log|x|,\log|y|-m\}$: near p the
first term is nonpositive and the second is $-m$, while off the p-ray
the first is zero and dominates.

On the ray through $p$ one has
$\phi_m(s\,\ord_p)=\max\{-s/m,-1\}\to0$ as $m\to\infty$, for each fixed
$s<\infty$; on every other ray $\phi_m=\max\{0,-1\}=0$; and
$\phi_m(v_{\triv})=0$.  Hence $\phi_m\to0$ pointwise on
$X^{\divis}$, that is, $\phi_m\to0$ in $\PSH_X(c_1(L))$.

At the proper-support point $v_{p,\triv}$, however,
$\log|x|(v_{p,\triv})=-\infty$ because $x$ lies in the support ideal,
so
\[
 \phi_m(v_{p,\triv})=\max\{-\infty,-1\}=-1
 \qquad\text{for every }m .
\]
Thus $\phi_m\to0$ in $\PSH_X(c_1(L))$ while
$\lim_m\phi_m(v_{p,\triv})=-1\neq0$.  Evaluation at
$v_{p,\triv}$ is therefore not continuous.  It is upper semicontinuous, in
accordance with \cite[Corollary~4.29]{BJ22}: the limit $0$ dominates
$\limsup_m\phi_m=-1$ at that point.

The sequence $(\phi_m)$ is monotone \emph{increasing}, and is not
decreasing; had it been decreasing,
\Cref{cor:ch35-divisorial-determination}\,(4) would have forced pointwise
convergence everywhere.  Indeed $\phi_m$ increases with $m$, and
increasing nets are exactly what the envelope property is about.
\end{example}

\begin{remark}[Where the theorem is used]
\label{rem:ch35-determination-uses}
\Cref{thm:ch35-divisorial-comparison} is invoked at the following places in
this part, and each use should be read as a use of
plurisubharmonicity, not of density.  In
\Cref{thm:ch35-homogeneous-siu} it upgrades
\eqref{eq:ch35-siu-slice-bound} and the final identity from
$X^{\divis}$ to $X^{\An}$.  In
\Cref{thm:bj-real-cartier-realization} it extends the ordering of the
rational stages and identifies the realization.  In
\Cref{thm:bj-homogeneous-reflection} and
\Cref{thm:bj-carrier-reflection} it converts equalities of divisorial traces
into equalities of functions.  In
\Cref{thm:ch35-cartier-psh-criterion} and
\Cref{thm:bj-nef-bdivisor-correspondence} it identifies the coefficient
function of a b-divisor with an honest psh function on the analytification.
In each case the smaller function is psh and the larger one is merely upper
semicontinuous, which is the exact shape of
\eqref{eq:ch35-divisorial-comparison}.
\end{remark}

\section{Legendre slices, test curves, and carriers}
\label{sec:ch35-legendre}

\subsection{The slices}
\label{sec:ch35-slices}

For $\lambda\in\mathbf R$, define the homogeneous Legendre slice
\begin{equation}
\label{eq:homogeneous-legendre-slice}
 \widehat\phi_\lambda
   :=\inf_{t>0}\{t\cdot\phi-t\lambda\}.
\end{equation}
It is the largest homogeneous $\psi$ satisfying
$\psi+\lambda\leq\phi$.

\begin{lemma}
\label{lem:ch35-slice-basic}
Let $\phi,\phi':X^{\An}\to\mathbf R\cup\{-\infty\}$, let
$\lambda,\mu,c\in\mathbf R$, and let $s,a>0$.  Then:
\begin{enumerate}
\item $\widehat\phi_\lambda=\widehat{\phi-\lambda}$ whenever
 $\phi-\lambda$ is defined, and in all cases
 $\widehat\phi_\lambda$ is homogeneous and is the largest homogeneous
 function $\psi$ with $\psi+\lambda\leq\phi$;
\item $\widehat\phi_0=\widehat\phi$, and $\lambda\mapsto\widehat\phi_\lambda$
 is nonincreasing;
\item $\widehat{\phi+c}_{\,\lambda}=\widehat\phi_{\lambda-c}$ and
 $\widehat{a\phi}_{\,\lambda}=a\,\widehat\phi_{\lambda/a}$;
\item $\widehat{s\cdot\phi}_{\,\lambda}=\widehat\phi_{\lambda/s}$;
\item $\phi\leq\phi'$ implies
 $\widehat\phi_\lambda\leq\widehat{\phi'}_\lambda$, and
 $\phi_a\searrow\phi$ implies
 $\widehat{(\phi_a)}_\lambda\searrow\widehat\phi_\lambda$;
\item for each $v\in X^{\An}$, the map
 $\lambda\mapsto\widehat\phi_\lambda(v)$ is concave and nonincreasing on
 $\mathbf R$, with values in $\mathbf R\cup\{-\infty\}$.
\end{enumerate}
\par
\end{lemma}

\begin{proof}
(1)  If $\phi-\lambda$ is defined, that is, if $\lambda$ is finite,
then $t\cdot(\phi-\lambda)=t\cdot\phi-t\lambda$ by
\Cref{lem:ch35-action-basic}\,(3); taking the infimum over $t$ gives
$\widehat{\phi-\lambda}=\widehat\phi_\lambda$.  Since
$\lambda\in\mathbf R$ always, this is the general case, and the extremal
description follows from
\Cref{lem:ch35-homogenization-maximal}\,(1)--(2) applied to
$\phi-\lambda$: a homogeneous $\psi$ satisfies $\psi+\lambda\leq\phi$
if and only if $\psi\leq\phi-\lambda$, a legitimate rearrangement because
$\lambda$ is finite, and this is ledger entry (2) of
\Cref{sec:extended-arithmetic}.

(2)  Immediate from (1) and monotonicity of
$\lambda\mapsto\phi-\lambda$ together with
\Cref{lem:ch35-homogenization-maximal}\,(3).

(3)  $\widehat{\phi+c}_{\,\lambda}=\widehat{\phi+c-\lambda}
=\widehat{\phi-(\lambda-c)}=\widehat\phi_{\lambda-c}$.  For the second,
$\widehat{a\phi}_{\,\lambda}=\widehat{a\phi-\lambda}
=\widehat{a(\phi-\lambda/a)}=a\widehat{\phi-\lambda/a}
=a\widehat\phi_{\lambda/a}$, using
\Cref{lem:ch35-homogenization-maximal}\,(5).

(4)  By \Cref{lem:ch35-action-basic}\,(1) and (3),
\[
 \widehat{s\cdot\phi}_{\,\lambda}
 =\inf_{t>0}\{t\cdot(s\cdot\phi)-t\lambda\}
 =\inf_{t>0}\{(ts)\cdot\phi-t\lambda\}
 =\inf_{u>0}\{u\cdot\phi-u(\lambda/s)\}
 =\widehat\phi_{\lambda/s},
\]
substituting $u=ts$.

(5)  Both follow from (1) and
\Cref{lem:ch35-homogenization-maximal}\,(3)--(4).

(6)  For fixed $v$ and $t>0$, the map
$\lambda\mapsto(t\cdot\phi)(v)-t\lambda$ is affine and nonincreasing, with
values in $\mathbf R\cup\{-\infty\}$.  An infimum of a family of
nonincreasing affine functions is concave and nonincreasing; concavity is
the statement that, for $0<\sigma<1$ and
$\lambda=(1-\sigma)\lambda_0+\sigma\lambda_1$,
\[
 (t\cdot\phi)(v)-t\lambda
 =(1-\sigma)\bigl[(t\cdot\phi)(v)-t\lambda_0\bigr]
 +\sigma\bigl[(t\cdot\phi)(v)-t\lambda_1\bigr]
 \geq(1-\sigma)\widehat\phi_{\lambda_0}(v)+\sigma\widehat\phi_{\lambda_1}(v),
\]
after which one takes the infimum over $t$.  The manipulation is valid in
$\mathbf R\cup\{-\infty\}$: if $(t\cdot\phi)(v)=-\infty$ all three
displayed quantities are $-\infty$, and if one of
$\widehat\phi_{\lambda_j}(v)$ is $-\infty$ then the right-hand side is
$-\infty$ and the inequality is trivial.  The endpoint cases are
tautological; treating them separately avoids the undefined product
$0\cdot(-\infty)$.
\end{proof}

\subsection{The trichotomy and the endpoint}
\label{sec:ch35-trichotomy}

\begin{proposition}
\label{prop:ch35-slice-trichotomy}
Let $\theta\in N^1(X)$, $\phi\in\PSH_X(\theta)$, and
$b:=\sup\phi=\phi(v_{\triv})\in\mathbf R$.  Then:
\begin{enumerate}
\item if $\lambda\leq\inf_{X^{\An}}\phi$, then $\widehat\phi_\lambda=0$;
\item if $\lambda\leq b$, then
 $\widehat\phi_\lambda\in\PSH_{\mathrm{hom},X}(\theta)$, and it admits a
 decreasing regularization by functions in $\cH_{\mathrm{hom}}(L_a)$ with
 $c_1(L_a)\to\theta$;
\item if $\lambda>b$, then $\widehat\phi_\lambda\equiv-\infty$.
\end{enumerate}
\par
\end{proposition}

\begin{proof}
(1)  If $\lambda\leq\inf\phi$, then $0+\lambda\leq\phi$ and $0$ is
homogeneous, so $0\leq\widehat\phi_\lambda$ by the extremal
characterization in \Cref{lem:ch35-slice-basic}\,(1).  On the other hand
$\widehat\phi_\lambda\leq0$ always, because
$\widehat\phi_\lambda(v_{\triv})\in\{0,-\infty\}$ by homogeneity and
$\widehat\phi_\lambda$ is decreasing by
\Cref{cor:ch35-homogenization-alternative} applied to $\phi-\lambda$.

(2)  By \Cref{lem:ch35-slice-basic}\,(1),
$\widehat\phi_\lambda=\widehat{\phi-\lambda}$, and
$\phi-\lambda\in\PSH_X(\theta)$ with
$\sup(\phi-\lambda)=b-\lambda\geq0$.  Now apply
\Cref{cor:ch35-homogenization-psh}.

(3)  If $\lambda>b$, then $\sup(\phi-\lambda)=b-\lambda<0$ and
\Cref{cor:ch35-homogenization-alternative} gives
$\widehat{\phi-\lambda}\equiv-\infty$.
\end{proof}

\begin{lemma}
\label{lem:ch35-slice-continuity}
Let $\theta\in N^1(X)$, $\phi\in\PSH_X(\theta)$ and $b=\sup\phi$.
Then for every $v\in X^{\divis}$ the map
$\lambda\mapsto\widehat\phi_\lambda(v)$ is real valued, concave and
nonincreasing on $(-\infty,b]$, and continuous there; in particular
\begin{equation}
\label{eq:ch35-endpoint-continuity}
 \widehat\phi_b=\inf_{\lambda<b}\widehat\phi_\lambda
 =\lim_{\lambda\uparrow b}\widehat\phi_\lambda
\end{equation}
pointwise on $X^{\divis}$, and hence, by
\Cref{cor:ch35-divisorial-determination}, on all of $X^{\An}$.
\end{lemma}

\begin{proof}
Finiteness on $X^{\divis}$ is \Cref{thm:bj-basic-closure}\,(7) applied to
$\widehat\phi_\lambda\in\PSH_{\mathrm{hom},X}(\theta)$, which is available
for $\lambda\leq b$ by \Cref{prop:ch35-slice-trichotomy}\,(2).  Concavity
and monotonicity are \Cref{lem:ch35-slice-basic}\,(6).  Applying
\Cref{lem:concave-on-body}\,(3) on each compact interval
$\Delta=[c,b]\subset\mathbf R$, whose relative interior is $(c,b)$, and
letting $c\to-\infty$ shows that
$\lambda\mapsto\widehat\phi_\lambda(v)$ is continuous on
$(-\infty,b)$; the hypotheses of that lemma are met because the function
is finite there and concave, hence proper and, being finite and concave on
an open interval, upper semicontinuous.

Only the endpoint $\lambda=b$ remains.  Put
$\psi:=\inf_{\lambda<b}\widehat\phi_\lambda$, a decreasing limit of
functions in $\PSH_{\mathrm{hom},X}(\theta)$ by
\Cref{lem:ch35-slice-basic}\,(2).  Each of them vanishes at
$v_{\triv}$ by \Cref{prop:ch35-homogeneous-normalization}, hence so does
$\psi$; in particular $\psi$ is generically finite and
\Cref{thm:bj-basic-closure}\,(4) gives $\psi\in\PSH_X(\theta)$.  It is
homogeneous, being an infimum of homogeneous functions.  Thus
$\psi\in\PSH_{\mathrm{hom},X}(\theta)$.

Monotonicity gives $\widehat\phi_b\leq\psi$.  For the reverse
inequality, note that for every $\lambda<b$ we have
$\psi\leq\widehat\phi_\lambda\leq\phi-\lambda$, whence
$\psi\leq\phi-\lambda$ for all $\lambda<b$ and therefore
$\psi\leq\phi-b$.  Here the subtraction is legitimate because
$\lambda,b\in\mathbf R$.  Since $\psi$ is homogeneous, the extremal
characterization of \Cref{lem:ch35-slice-basic}\,(1) gives
$\psi\leq\widehat\phi_b$.  Hence $\psi=\widehat\phi_b$, which is
\eqref{eq:ch35-endpoint-continuity} and the asserted continuity at the
endpoint.  This is \cite[Lemma~6.23]{BJ22}.
\end{proof}

\begin{theorem}\label{thm:homogeneous-legendre-decomposition}
Let $\phi\in\PSH_X(\theta)$ and $b=\sup\phi$.  Then:
\begin{enumerate}
\item $\widehat\phi_\lambda\in
      \PSH_{\mathrm{hom},X}(\theta)$ for $\lambda\leq b$, and it is
      identically $-\infty$ for $\lambda>b$;
\item $\lambda\mapsto\widehat\phi_\lambda(v)$ is decreasing and concave,
      and is continuous on $(-\infty,b]$ on divisorial points;
\item pointwise on all of $X^{\An}$,
\begin{equation}
\label{eq:legendre-reconstruction-bj}
                   \phi=\sup_{\lambda\leq b}
                         \{\widehat\phi_\lambda+\lambda\}.
\end{equation}
\end{enumerate}
\par
\end{theorem}

\begin{proof}[Scalar Legendre duality]
For fixed $v$, the map $t\mapsto(t\cdot\phi)(v)$ is convex; this follows
first from the finite maximum \eqref{eq:scaling-fs-weights}, then by
decreasing regularization.  Its negative Legendre transform is
$\lambda\mapsto\widehat\phi_\lambda(v)$.  Fenchel--Moreau duality at
$t=1$ gives \eqref{eq:legendre-reconstruction-bj}.  Concavity and
monotonicity are immediate from the infimum of affine functions of
$\lambda$.  The psh statement follows by applying
\Cref{prop:homogenization-gf} to a decreasing GF regularization
of $\phi-\lambda$, using the decreasing-limit axiom and the value zero at
$v_{X,\triv}$ to exclude the identically $-\infty$ alternative at the
endpoint.  See \cite[Lemma~6.23 and Theorem~6.24]{BJ22}.

In detail, (1) is \Cref{prop:ch35-slice-trichotomy}\,(2)--(3), and (2) is
\Cref{lem:ch35-slice-basic}\,(6) together with
\Cref{lem:ch35-slice-continuity}.  We prove (3).

Fix $v\in X^{\An}$ and set
\[
 F(t):=(t\cdot\phi)(v)=t\,\phi(t^{-1}v),\qquad t>0 .
\]
By \Cref{prop:ch35-scaling-convex}, $F$ is convex on
$\mathbf R_{>0}$ with values in $\mathbf R\cup\{-\infty\}$, and by
\eqref{eq:ch35-orbit-slope-bound} it satisfies $F(t)\leq bt$. Two cases
occur.

\emph{Case A: $F\equiv-\infty$ on $\mathbf R_{>0}$.}  Then
$\phi(t^{-1}v)=-\infty$ for every $t>0$; taking $t=1$ gives
$\phi(v)=-\infty$.  On the other hand
$\widehat\phi_\lambda(v)\leq F(1)-\lambda=-\infty$ for every
$\lambda$, so both sides of
\eqref{eq:legendre-reconstruction-bj} equal $-\infty$ at $v$, the
right-hand side by \Cref{lem:extended-real-safe}\,(1).

\emph{Case B: $F(t_0)>-\infty$ for some $t_0>0$.} The orbit alternative
proved in \Cref{prop:ch35-scaling-convex} shows that $F$ is finite on all
of $(0,+\infty)$. A finite convex function has a supporting line at the
interior point 1: choose a slope $\lambda$ between its left and right
derivatives there. Then
\[
 F(t)\geq F(1)+\lambda(t-1)\qquad(t>0).
\]
Since $F(t)\leq bt$, letting $t\to+\infty$ in this inequality gives
$\lambda\leq b$. The infimum defining the slice is attained at $t=1$:
\[
 \widehat\phi_\lambda(v)=\inf_{t>0}\{F(t)-t\lambda\}
 =F(1)-\lambda.
\]
Consequently $\widehat\phi_\lambda(v)+\lambda=F(1)=\phi(v)$, whereas
all slice competitors are bounded above by $\phi(v)$ by their extremal
characterization. This proves \eqref{eq:legendre-reconstruction-bj}.
The supporting-line argument is the interior-point biconjugation step of
\Cref{thm:appD-fenchel-moreau}; no closedness of an extension at $t=0$
is required.

Since $v$ was arbitrary, \eqref{eq:legendre-reconstruction-bj} holds
pointwise on all of $X^{\An}$, including at proper-support
semivaluations, and no upper-semicontinuous regularization has been applied
to the right-hand side.
\end{proof}

\begin{warning}{No star, and why that matters}
The right-hand side of \eqref{eq:legendre-reconstruction-bj} is a supremum
of a family of $\theta$-psh functions, and such a supremum is in general
neither upper semicontinuous nor psh: that is the content of the envelope
problem, \Cref{def:bj-envelope-property}.  What
\Cref{thm:homogeneous-legendre-decomposition} asserts is that for
\emph{this particular} family the supremum is already equal to
$\phi$ pointwise, hence automatically upper semicontinuous and psh.  The
reason is visible in the proof: each competitor satisfies
$\widehat\phi_\lambda+\lambda\leq\phi$ by construction, so the supremum is
squeezed between a family of minorants and $\phi$, and scalar
biconjugation shows there is no gap.  No envelope theorem is used, and this
is why \Cref{chap:bj-reflection} may use the decomposition inside the proof
of the envelope theorem itself.
\end{warning}

\begin{remark}[The family is a concave test curve]
\label{rem:ch35-test-curve}
Use the ambient ordered convex cone $\mathcal C$ of all nonpositive
homogeneous extended-real functions on $X^{\An}$, with the pointwise
order. The fixed-class set
$\PSH_{\mathrm{hom},X}(\theta)\cup\{-\infty\}$ is convex inside it;
arbitrary positive scalar multiplication changes the numerical class.  By
\Cref{lem:ch35-slice-basic}\,(6) and
\Cref{prop:ch35-slice-trichotomy}, the family
\[
 \lambda\longmapsto\psi_\lambda:=\widehat\phi_\lambda,
 \qquad\lambda\in\mathbf R,
\]
is nonincreasing, concave in the sense of
\eqref{eq:appD-test-curve-concavity}, and pointwise upper semicontinuous in
$\lambda$: at every analytic point it is an infimum of affine functions
of $\lambda$, while at divisorial points it is even continuous up to the
finite endpoint by \Cref{lem:ch35-slice-continuity}.  It is therefore a
concave test curve in the
sense of \Cref{def:appD-test-curve}, and
\eqref{eq:legendre-reconstruction-bj} is the value at $t=1$ of its formal
Legendre ray \eqref{eq:appD-test-curve-ray}
\[
 \ell_t=\sup_{\lambda\in\mathbf R}\{\psi_\lambda+t\lambda\}.
\]
Indeed $\ell_t=t\cdot\phi$ for every $t>0$: apply
\eqref{eq:legendre-reconstruction-bj} to $t\cdot\phi$, whose slices are
$\widehat{t\cdot\phi}_{\,\mu}=\widehat\phi_{\mu/t}$ by
\Cref{lem:ch35-slice-basic}\,(4), and substitute $\mu=t\lambda$.  Thus the
scaling orbit of $\phi$ \emph{is} the Legendre ray of its slice family,
which is the precise form of the analogy with the Ross--Witt Nystr\"om
correspondence between test curves and geodesic rays
\cite{RWN14}; see \Cref{chap:rays-test-curves} for the complex
counterpart and \Cref{chap:rays-filtrations} for the comparison.
\end{remark}

\begin{proposition}
\label{prop:ch35-phimax}
Let $\phi\in\PSH_X(\theta)$, $b=\sup\phi$, and set
\begin{equation}
\label{eq:ch35-phimax}
 \widehat\phi_{\max}:=\widehat\phi_b\in\PSH_{\mathrm{hom},X}(\theta).
\end{equation}
Then:
\begin{enumerate}
\item $\widehat\phi_{\max}$ is the largest homogeneous function
 $\psi$ with $\psi\leq\phi-b$, and
 $\widehat{(\phi+c)}_{\max}=\widehat\phi_{\max}$ for every
 $c\in\mathbf R$;
\item $\widehat\phi_{\max}=0$ if and only if $\phi$ is constant;
\item for every $v\in X^{\An}$,
\begin{equation}
\label{eq:ch35-gateaux}
 \widehat\phi_{\max}(v)
 =\lim_{s\downarrow0}\frac{\phi(sv)-\phi(v_{\triv})}{s},
\end{equation}
the limit existing in $\mathbf R\cup\{-\infty\}$ and being monotone.
\end{enumerate}
\par
\end{proposition}

\begin{proof}
(1)  The first assertion is \Cref{lem:ch35-slice-basic}\,(1) at
$\lambda=b$.  For the second,
$\sup(\phi+c)=b+c$ and
$\widehat{(\phi+c)}_{\,b+c}=\widehat\phi_{(b+c)-c}=\widehat\phi_b$ by
\Cref{lem:ch35-slice-basic}\,(3).

(2)  If $\phi$ is constant, then $\phi-b=0$ and the largest homogeneous
minorant of $0$ is $0$.  Conversely, if
$\widehat\phi_{\max}=0$, then $0\leq\phi-b$, that is,
$\phi\geq b=\sup\phi$, so $\phi\equiv b$.

(3)  Replacing $\phi$ by $\phi-b$ changes neither side, by (1) and by
$\phi(v_{\triv})=b$; so assume $b=0$, that is,
$\phi\leq0$ and $\phi(v_{\triv})=0$.  By
\Cref{prop:ch35-scaling-convex}, $t\mapsto(t\cdot\phi)(v)$ is then
nonincreasing, so
\[
 \widehat\phi_{\max}(v)=\widehat\phi_0(v)=\widehat\phi(v)
 =\inf_{t>0}(t\cdot\phi)(v)=\lim_{t\to+\infty}t\,\phi(t^{-1}v)
 =\lim_{s\downarrow0}\frac{\phi(sv)}{s},
\]
with $s=t^{-1}$; the middle equality is
\Cref{lem:monotone-net-converges}.  Since $\phi(v_{\triv})=0$, this is
\eqref{eq:ch35-gateaux}.  This is \cite[Lemma~6.29]{BJ22}.
\end{proof}

\begin{example}[A Legendre pair computed on $\mathbf P^1$]
\label{ex:ch35-legendre-pair-p1}
Let $X=\mathbf P^1_k$, $L=\cO(1)$, $\theta=c_1(L)$, and
\[
 \phi=\max\{\log|x|,\ \log|y|-1\}
 \in\cH_{\mathbf Q}(L)\subset\PSH_X(\theta),
\]
the function \eqref{eq:p1-fs} with $a=0$, $b=-1$.  Write
$p=[0:1]$ and $q=[1:0]$.  In the trivialization by $y$ near $p$ one
has $\log|x|(s\,\ord_p)=-s$ and $\log|y|=0$; in the trivialization by
$x$ near $q$ one has $\log|y|(s\,\ord_q)=-s$ and $\log|x|=0$.
Hence
\[
 \phi(s\,\ord_p)=\max\{-s,-1\},
 \qquad
 \phi(s\,\ord_q)=\max\{0,-s-1\}=0,
\]
and $\phi\equiv0$ on every other ray.  Thus $\sup\phi=0=b$,
$\inf\phi=-1$.

\emph{The slices.}  Fix $\lambda\leq0$.  A homogeneous function
$\psi$ on $X^{\An}$ is determined by its slopes
$c_r:=-\psi(\ord_r)\geq0$ on each finite ray, with the psh endpoint
values fixed by their radial limits, and
$\psi+\lambda\leq\phi$ reads
\[
 -c_ps+\lambda\leq\max\{-s,-1\}\ \ (s\geq0),
 \qquad
 -c_rs+\lambda\leq0\ \ (r\neq p,\ s\geq0).
\]
For $r\neq p$ the constraint at $s=0$ gives $\lambda\leq0$, which
holds, and for $s>0$ it is implied; the largest admissible function therefore has
$c_r=0$, so $\widehat\phi_\lambda$ vanishes off the ray
through $p$.  On that ray, the constraint at $s\geq1$ reads
$-c_ps+\lambda\leq-1$, that is,
$c_p\geq(\lambda+1)/s$ for all $s\geq1$, i.e.
$c_p\geq\lambda+1$ when $\lambda\geq-1$ and no constraint when
$\lambda<-1$; the constraint for $0\leq s\leq1$ reads
$-c_ps+\lambda\leq-s$, that is,
$(1-c_p)s\leq-\lambda$ for $0\leq s\leq1$, which is implied by
$1-c_p\leq-\lambda$, i.e. $c_p\geq1+\lambda$, the same condition.
Taking the smallest admissible $c_p$, that is, the largest
$\psi$,
\[
 \widehat\phi_\lambda(s\,\ord_p)=-\max\{0,\ 1+\lambda\}\,s,
 \qquad
 \widehat\phi_\lambda=0\ \text{on all other rays},
\]
for $\lambda\leq0$; and $\widehat\phi_\lambda\equiv-\infty$ for
$\lambda>0$.  In particular
\[
 \widehat\phi_\lambda=0\ (\lambda\leq-1),
 \qquad
 \widehat\phi_{\max}=\widehat\phi_0=\log|\mathfrak m_p|
\]
in the notation of \Cref{ex:ch35-point-ideal-p1}.  The function
$\lambda\mapsto\widehat\phi_\lambda(\ord_p)=-\max\{0,1+\lambda\}$ is
concave, nonincreasing, continuous on $(-\infty,0]$, and identically zero
for $\lambda\leq-1=\inf\phi$, as
\Cref{prop:ch35-slice-trichotomy}\,(1) and
\Cref{lem:ch35-slice-continuity} require.

\emph{The reconstruction.}  On the ray through $p$,
\[
 \sup_{\lambda\leq0}\{\widehat\phi_\lambda(s\,\ord_p)+\lambda\}
 =\sup_{\lambda\leq0}\{-\max\{0,1+\lambda\}s+\lambda\} .
\]
For $\lambda\leq-1$ the expression is $\lambda$, maximized at
$\lambda=-1$ with value $-1$; for $-1\leq\lambda\leq0$ it is
$-(1+\lambda)s+\lambda$, an affine function of $\lambda$ with slope
$1-s$, maximized at $\lambda=0$ with value $-s$ if $s\leq1$ and at
$\lambda=-1$ with value $-1$ if $s\geq1$.  The overall supremum is
$\max\{-s,-1\}=\phi(s\,\ord_p)$, confirming
\eqref{eq:legendre-reconstruction-bj}.  On every other ray both sides are
$0$.

\emph{The Legendre ray.}  Consistently with
\Cref{rem:ch35-test-curve}, one computes
$\ell_t=\sup_\lambda\{\widehat\phi_\lambda+t\lambda\}$ on the ray through
$p$ to be $\max\{-s,-t\}$, which is exactly
$(t\cdot\phi)(s\,\ord_p)=t\,\phi(t^{-1}s\,\ord_p)
=t\max\{-s/t,-1\}=\max\{-s,-t\}$.
\end{example}

\begin{example}[A test curve with an infinite value]
\label{ex:ch35-test-curve-infinite}
Keep $X=\mathbf P^1$ and $L=\cO(1)$ but take
$\phi=\log|x|\in\cH_{\mathrm{hom}}(L)$, which is already homogeneous, with
$\sup\phi=0$ and $\phi(v_{p,\triv})=-\infty$ at $p=[0:1]$.  Then
$t\cdot\phi=\phi$ for all $t$, so
\[
 \widehat\phi_\lambda=\inf_{t>0}\{\phi-t\lambda\}=
 \begin{cases}
  \phi,&\lambda\leq0,\\
  -\infty,&\lambda>0,
 \end{cases}
\]
and the reconstruction \eqref{eq:legendre-reconstruction-bj} reads
$\phi=\sup_{\lambda\leq0}\{\phi+\lambda\}=\phi$, attained at
$\lambda=0$.  At the proper-support point $v_{p,\triv}$ every entry of
the family equals $-\infty$ and the identity is the tautology
$-\infty=-\infty$.  This is the smallest example in which
\eqref{eq:legendre-reconstruction-bj} would be destroyed by rewriting it as
a statement about $\phi-\widehat\phi_\lambda$: that difference is
undefined at $v_{p,\triv}$ by ledger entry (4) of
\Cref{sec:extended-arithmetic}.
\end{example}

\subsection{Carrier functions}
\label{sec:ch35-carriers}

A homogeneous carrier is a function
$p\in\PSH_{\mathrm{hom},X}(\eta)$ controlling the radial tail of an
inhomogeneous function:
\begin{equation}
\label{eq:homogeneous-carrier}
                         \phi\leq p+C.
\end{equation}
The inequality is meaningful at proper-support points: if $p=-\infty$,
it forces $\phi=-\infty$ there.  It is not rewritten as
$\phi-p\leq C$, since $(-\infty)-(-\infty)$ is undefined.

\begin{definition}[Carrier]
\label{def:ch35-carrier}
Let $\theta,\eta\in N^1(X)$, let $\phi\in\PSH_X(\theta)$, and let
$p\in\PSH_{\mathrm{hom},X}(\eta)$.  We say that $p$ is a
\emph{carrier} for $\phi$ if \eqref{eq:homogeneous-carrier} holds for some
finite constant $C\in\mathbf R$.
\end{definition}

The constant belongs to the inequality, not to $p$: by
\Cref{lem:ch35-homogeneity-reformulations}, adding a nonzero constant to a
homogeneous function destroys homogeneity.  The class $\eta$ is allowed to
differ from $\theta$; in \Cref{chap:bj-reflection} it is
$\theta$ itself, whereas in the applications of
\Cref{chap:bj-finite-descent} it is an ample class chosen for convenience.

\begin{lemma}
\label{lem:ch35-carrier-tail}
Let $\phi\in\PSH_X(\theta)$ and let
$p\in\PSH_{\mathrm{hom},X}(\eta)$ satisfy
$\phi\leq p+C$ with $C\in\mathbf R$.  Then
\begin{equation}
\label{eq:ch35-slice-below-carrier}
 \widehat\phi_\lambda\leq p
 \qquad\text{for every }\lambda\leq C .
\end{equation}
Consequently, for $\lambda\leq\min\{C,\sup\phi\}$,
\begin{equation}
\label{eq:ch35-tail-decay}
 \widehat\phi_\lambda+\lambda\leq p+\lambda,
\end{equation}
so the part of the family \eqref{eq:legendre-reconstruction-bj} with
$\lambda\leq-M$ is dominated by $p-M$ for every
$M\geq-\min\{C,\sup\phi\}$.
\end{lemma}

\begin{proof}
Fix $\lambda\leq C$ and $t>0$.  Applying
\Cref{lem:ch35-action-basic}\,(2) and (3) to
$\phi\leq p+C$ gives
\[
 t\cdot\phi\leq t\cdot(p+C)=t\cdot p+tC=p+tC,
\]
the last equality because $p$ is homogeneous.  Hence
\[
 t\cdot\phi-t\lambda\leq p+t(C-\lambda),
\]
an inequality in $\mathbf R\cup\{-\infty\}$; the subtraction of the finite
number $t\lambda$ is legitimate by ledger entry (2) of
\Cref{sec:extended-arithmetic}, and if $p(v)=-\infty$ at some $v$, then
$\phi\leq p+C$ forces $\phi(t^{-1}v)=-\infty$ and both sides are
$-\infty$.  Taking the infimum over $t>0$ and using
$C-\lambda\geq0$, let $t\downarrow0$: the right-hand side tends to
$p(v)$ when $p(v)$ is finite, and equals $-\infty$ otherwise.  Thus
\[
 \widehat\phi_\lambda(v)\leq\inf_{t>0}\{p(v)+t(C-\lambda)\}=p(v),
\]
which is \eqref{eq:ch35-slice-below-carrier}.  Inequality
\eqref{eq:ch35-tail-decay} is \eqref{eq:ch35-slice-below-carrier} plus the
finite constant $\lambda$.
\end{proof}

\begin{remark}[What the carrier is for]
\label{rem:ch35-carrier-purpose}
The reconstruction \eqref{eq:legendre-reconstruction-bj} is a supremum over
an unbounded parameter interval $(-\infty,b]$.
\Cref{lem:ch35-carrier-tail} says that a carrier truncates that interval:
outside a compact window $[-M,b]$ the competitors are dominated by the
single function $p-M$.  Since $\PSH_X(\theta)$ is stable under finite
maxima but not under arbitrary suprema before the envelope theorem, this
truncation is what turns \eqref{eq:legendre-reconstruction-bj} into a
statement that can be manipulated with the closure properties of
\Cref{thm:bj-basic-closure} alone.  The construction is carried out in
\Cref{lem:bj-compact-legendre-window} and used in
\Cref{thm:bj-carrier-reflection}; we do not repeat it here.
\end{remark}

\begin{example}[The tautological carrier of a bounded function]
\label{ex:ch35-bounded-carrier}
If $\phi\in\PSH_X(\theta)$ is bounded, say $\phi\geq-M$ and
$\phi\leq b$, then $p=0\in\PSH_{\mathrm{hom},X}(\eta)$ is a carrier for
every nef $\eta$, with $C=b$; indeed $0$ is $\eta$-psh for nef
$\eta$ by \eqref{eq:psh-existence-cones}.  In that case
\eqref{eq:ch35-slice-below-carrier} says $\widehat\phi_\lambda\leq0$,
which we already know.  A carrier carries information only when $\phi$ is
unbounded, and the interesting case is $p=\log|s_E|$ for an effective
divisor $E$.
\end{example}

\begin{example}[A divisorial carrier]
\label{ex:ch35-divisorial-carrier}
Let $E\geq0$ be an effective $\mathbf Q$-Cartier divisor on $X$ and
let $\log|s_E|$ be its section function
\eqref{eq:divisor-section-pole}, an element of
$\cH_{\mathrm{hom}}(E)\subset\PSH_{\mathrm{hom},X}([E])$ by
\Cref{ex:gf-proper-support-pole}.  A function
$\phi\in\PSH_X(\theta)$ admits $\log|s_E|$ as a carrier if and only if
$\phi\leq\log|s_E|+C$ for some $C$, that is, if and only if the
singularity of $\phi$ along $E$ is at least that of the divisor.

On $\mathbf P^1$ with $E=p=[0:1]$, so that $\log|s_E|=\log|x|$, the
function $\phi=\max\{\log|x|,\log|y|-1\}$ of
\Cref{ex:ch35-legendre-pair-p1} does \emph{not} admit $\log|x|$ as a
carrier: at the proper-support point $v_{p,\triv}$ one has
$\phi(v_{p,\triv})=-1$ while $\log|x|(v_{p,\triv})=-\infty$, so no
finite $C$ satisfies \eqref{eq:homogeneous-carrier}.  The conclusion of
\Cref{lem:ch35-carrier-tail} fails accordingly: by
\Cref{ex:ch35-legendre-pair-p1} the slices for $\lambda\leq0$ are zero when $\lambda\leq-1$ and
$(1+\lambda)\log|x|$ when $-1<\lambda\leq0$. A zero coefficient means
that the term is omitted, including at the polar endpoint. Since
$\log|x|\leq0$ with strict inequality off the vertex, the inequality
$\widehat\phi_\lambda\leq\log|x|$ holds only when
$\max\{0,1+\lambda\}\geq1$, that is, only at the endpoint
$\lambda=0=\sup\phi$.  For every $\lambda<0$ the slice lies strictly
above $\log|x|$ on each divisorial ray through $p$.  Thus the carrier
hypothesis in \eqref{eq:ch35-slice-below-carrier} cannot be dropped.

By contrast, the homogeneous function
$\psi:=\log|x|$ itself has $\log|x|$ as a carrier with $C=0$, and
\Cref{thm:ch35-homogeneous-siu-general} splits it as
$\psi=\log|s_E|+0$ with $0\in\PSH_{\mathrm{hom},X}(c_1(\cO(1))-[p])=
\PSH_{\mathrm{hom},X}(0)$; the residual class is trivial and the residual
function is the zero function, in agreement with
\Cref{cor:ch35-zero-class}.  The strict form
\Cref{thm:ch35-homogeneous-siu} does \emph{not} apply here:
\eqref{eq:ch35-siu-strict-hypothesis} asks for
$\psi\leq t_0\log|s_E|$ with $t_0>1$, which for
$\psi=\log|s_E|=\log|x|$ would force $(t_0-1)\log|x|\geq0$, impossible
since $\log|x|<0$ off the vertex.  This is exactly the endpoint
$u=1$ of \Cref{rem:ch35-siu-boundary}, which is why the general form is
the one cited.
\end{example}

\subsection{Homogeneous Siu decomposition}
\label{sec:ch35-siu}

The next results convert the inequality $\psi\leq\log|s_E|$ into an
identity.  This is the statement that \Cref{chap:bj-cartier-realization}
calls homogeneous Siu division.  It belongs to the homogeneous calculus and
uses no envelope property, and we prove it in the form in which the later
chapters apply it.

Throughout the subsection $X$ is normal, $E$ is an effective
$\mathbf Q$-Cartier divisor on $X$, and $\log|s_E|$ is its section
function \eqref{eq:divisor-section-pole}.  We write
\begin{equation}
\label{eq:ch35-siu-slice}
 K_E:=\{v\in X^{\An}\mid\log|s_E|(v)=-1\},
\end{equation}
a compact subset of $X^{\An}$: the function $\log|s_E|$ is continuous
into $[-\infty,0]$ by \Cref{prop:ch35-hom-pl-semiring}, so $K_E$ is
closed in the compact space $X^{\An}$.  We assume $E\neq0$, so that
$K_E\neq\varnothing$.

\begin{lemma}
\label{lem:ch35-slice-to-cone}
Let $\eta\in N^1(X)$, let $\chi\in\PSH_{\mathrm{hom},X}(\eta)$, and let
$t\in\mathbf R_{>0}$.  Then
\begin{equation}
\label{eq:ch35-slice-to-cone}
 \chi\leq t\log|s_E|\ \text{ on }X^{\An}
 \iff
 \sup_{K_E}\chi\leq-t .
\end{equation}
\end{lemma}

\begin{proof}
$\Rightarrow$ is immediate, since $t\log|s_E|=-t$ on $K_E$, which is
nonempty: if $F$ is a component of $E$ with coefficient $a_F>0$, then
$a_F^{-1}\ord_F\in K_E$.

For $\Leftarrow$ we first check the inequality on $X^{\divis}$, where
$c:=v(E)$ is finite.  If $0<c<\infty$, then $c^{-1}v\in K_E$, so
homogeneity of $\chi$ gives
\[
 \chi(v)=c\,\chi(c^{-1}v)\leq-tc=t\log|s_E|(v).
\]
If $c=0$, then $t\log|s_E|(v)=0\geq\chi(v)$, the last inequality by
\Cref{prop:ch35-homogeneous-normalization}.

Thus $\chi\leq t\log|s_E|$ on $X^{\divis}$.  The function
$t\log|s_E|$ is upper semicontinuous on $X^{\An}$, by
\Cref{prop:ch35-hom-pl-semiring}, and $\chi$ is $\eta$-psh; so
\Cref{thm:ch35-divisorial-comparison} promotes the inequality to all of
$X^{\An}$.
\end{proof}

\begin{lemma}
\label{lem:ch35-fs-division}
Let $X$ be normal, $L\in\Pic(X)_{\mathbf Q}$,
$\psi\in\cH_{\mathrm{hom}}(L)$, and $t\in\mathbf Q_{>0}$.  If
$\psi\leq t\log|s_E|$, then
\begin{equation}
\label{eq:ch35-fs-division}
 \psi=\log|s_{tE}|+\chi,
 \qquad\chi\in\cH_{\mathrm{hom}}(L-tE).
\end{equation}
The residual is unique and equals $\psi-t\log|s_E|$ on $X^{\val}$,
where subtraction is defined.
\end{lemma}

\begin{proof}
Choose $m\in\mathbf Z_{>0}$ so divisible that $mL$ is a line bundle,
$mtE$ is an integral Cartier divisor, and
$\psi=m^{-1}\max_i\log|\sigma_i|$ with
$0\neq\sigma_i\in H^0(X,mL)$, as in
\eqref{eq:homogeneous-fs-presentation}.  Then
$\log|s_{tE}|=m^{-1}\log|s_{mtE}|$, with
$s_{mtE}\in H^0(X,\cO_X(mtE))$ the canonical section, and the hypothesis
reads
\[
 \log|\sigma_i|\leq\log|s_{mtE}|\qquad\text{for every }i,
\]
that is, $v(\sigma_i)\geq v(mtE)$ for every $v\in X^{\An}$.  Taking
$v=\ord_F$ for the prime divisors $F\subset X$ gives
$\operatorname{div}(\sigma_i)\geq mtE$ as Weil divisors on $X$.

Since $X$ is normal and $mtE$ is Cartier, the quotient
$\tau_i:=\sigma_i/s_{mtE}$ is a rational section of
$m(L-tE)$ whose divisor $\operatorname{div}(\sigma_i)-mtE$ is effective;
by normality a rational section of a line bundle with effective divisor is
regular.  Indeed, in a local trivialization the section is a rational
function $f$ on the normal variety $X$ with
$\ord_F(f)\geq0$ for every prime divisor $F$; since a noetherian normal
domain is the intersection of its localizations at height-one primes, each
of which is a discrete valuation ring, $f$ is regular
\cite[Proposition~II.6.3A]{Har}.  Hence
$\tau_i\in H^0(X,m(L-tE))$, and $\tau_i\neq0$ because
$\sigma_i\neq0$.  Define $\chi:=m^{-1}\max_i\log|\tau_i|$, an element of
$\cH_{\mathrm{hom}}(L-tE)$. The regular-section identity
$\sigma_i=\tau_i s_{mtE}$ gives, on all of $X^{\An}$,
\[
 \psi=\frac1m\max_i\{\log|\tau_i|+\log|s_{mtE}|\}
 =\chi+\log|s_{tE}|.
\]
All summands lie in $\mathbf R\cup\{-\infty\}$, so this addition is
defined even on the polar locus. On $X^{\val}$ it identifies the residual
with the stated difference. Two residuals agree there by finite
cancellation, hence agree everywhere by divisorial determination.
\end{proof}

\begin{theorem}
\label{thm:ch35-homogeneous-siu}
Assume $X$ is normal.  Let $\theta\in N^1(X)$, let $E\neq0$ be an
effective $\mathbf Q$-Cartier divisor on $X$, and let
$\psi\in\PSH_{\mathrm{hom},X}(\theta)$ satisfy
\begin{equation}
\label{eq:ch35-siu-strict-hypothesis}
 \psi\leq t_0\log|s_E|\qquad\text{for some }t_0>1 .
\end{equation}
Then
\begin{equation}
\label{eq:ch35-homogeneous-siu}
 \psi=\log|s_E|+\psi'
 \qquad\text{with }\psi'\in\PSH_{\mathrm{hom},X}(\theta-[E]),
\end{equation}
and $\psi'$ is unique.  The identity is asserted pointwise on all of
$X^{\An}$, both sides being $-\infty$ at every semivaluation supported
in $E$.  Conversely, if $\psi=\log|s_E|+\psi'$ with
$\psi'\in\PSH_{\mathrm{hom},X}(\theta-[E])$, then
$\psi\leq\log|s_E|$.
\end{theorem}

\begin{proof}
\emph{The converse.}  If $\psi=\log|s_E|+\psi'$ with $\psi'$ homogeneous
$(\theta-[E])$-psh, then $\psi'\leq0$ by
\Cref{prop:ch35-homogeneous-normalization}, so $\psi\leq\log|s_E|$; the
sum is defined everywhere because $\psi'$ never takes the value
$+\infty$.

\medskip
\noindent\emph{Step 1: a homogeneous regularization.}
By \Cref{cor:ch35-homogenization-psh} applied to $\psi$, which is its own
homogenization, there are a decreasing net $(\psi_a)$, rational line
bundles $L_a$ with $c_1(L_a)\to\theta$, and
$\psi_a\in\cH_{\mathrm{hom}}(L_a)$ with $\psi_a\searrow\psi$
pointwise.

\medskip
\noindent\emph{Step 2: the strict hypothesis makes Dini's theorem
effective.}
By \Cref{lem:ch35-slice-to-cone}, hypothesis
\eqref{eq:ch35-siu-strict-hypothesis} says $\sup_{K_E}\psi\leq-t_0<-1$.
Each $\psi_a$ is upper semicontinuous, and the net decreases to $\psi$
on the compact set $K_E$. Compactness for decreasing upper-semicontinuous
functions, \Cref{prop:dini-usc}\,(2), applied with $c=-1$, gives an
index $a_0$ with $\psi_a<-1$ on $K_E$ for all $a\geq a_0$, since
$\psi\leq-t_0<-1$ there. In particular $\sup_{K_E}\psi_a\leq-1$.  Replacing the directed
set by its cofinal subset $\{a\geq a_0\}$, we may assume this for every
$a$.  \emph{The constant $t=1$ is now fixed once and for all}; this is
the point at which the strict hypothesis is used, and
\Cref{rem:ch35-siu-boundary} explains what happens without it.

\medskip
\noindent\emph{Step 3: division.}
By \Cref{lem:ch35-slice-to-cone} again,
\begin{equation}
\label{eq:ch35-siu-slice-bound}
 \psi_a\leq\log|s_E|
 \qquad\text{on }X^{\An},\text{ for every }a .
\end{equation}
The inequality is first obtained on $X^{\divis}$, where the slice
computation of \Cref{lem:ch35-slice-to-cone} takes place, and is then valid
on all of $X^{\An}$ by \Cref{thm:ch35-divisorial-comparison}.  By \Cref{lem:ch35-fs-division} with
$t=1$,
\[
 \psi_a=\log|s_E|+\psi'_a,
 \qquad\psi'_a\in\cH_{\mathrm{hom}}(L_a-E).
\]
Here $\psi'_a$ is the section-defined residual of that lemma; it equals
the finite difference on $X^{\val}$.

\medskip
\noindent\emph{Step 4: pass to the limit.}
On $X^{\divis}$ the function $\log|s_E|$ is finite, so
$\psi'_a=\psi_a-\log|s_E|$ is a decreasing net there, the subtracted
function being independent of $a$; by
\Cref{cor:ch35-divisorial-determination}\,(1) the net is decreasing on all
of $X^{\An}$.  Each $\psi'_a$ vanishes at $v_{\triv}$, being
homogeneous psh, so the pointwise limit
$\psi':=\lim_a\psi'_a$ is generically finite.  Since
$c_1(L_a)-[E]\to\theta-[E]$, \Cref{thm:bj-basic-closure}\,(5) gives
$\psi'\in\PSH_X(\theta-[E])$, and $\psi'$ is homogeneous as an infimum
of homogeneous functions.

On $X^{\divis}$, all terms being finite,
\[
 \psi=\lim_a\psi_a=\lim_a\bigl(\psi'_a+\log|s_E|\bigr)
 =\psi'+\log|s_E| .
\]
Both $\psi$ and $\psi'+\log|s_E|$ are upper semicontinuous on
$X^{\An}$ and both are $\theta$-psh, so
\Cref{cor:ch35-divisorial-determination}\,(2) upgrades the identity to
$X^{\An}$.  At a semivaluation supported in $E$, both sides are
$-\infty$, and no difference of infinite values has been formed anywhere
in the argument.

\emph{Uniqueness.}  If $\psi'_1,\psi'_2$ both satisfy
\eqref{eq:ch35-homogeneous-siu}, they agree on $X^{\divis}$, where
$\log|s_E|$ is finite and may be cancelled; then
\Cref{cor:ch35-divisorial-determination}\,(2) gives $\psi'_1=\psi'_2$.
\end{proof}

\begin{theorem}
\label{thm:ch35-homogeneous-siu-general}
Assume $X$ is normal.  Let $\theta\in N^1(X)$, let $E$ be an effective
$\mathbf Q$-Cartier divisor on $X$, and let
$\psi\in\PSH_{\mathrm{hom},X}(\theta)$.  Then
\begin{equation}
\label{eq:ch35-homogeneous-siu-general}
 \psi\leq\log|s_E|
 \iff
 \psi=\log|s_E|+\psi'\ \text{ for some }
 \psi'\in\PSH_{\mathrm{hom},X}(\theta-[E]),
\end{equation}
and $\psi'$ is then unique.  No envelope hypothesis is required.
\end{theorem}

\begin{proof}
If $E=0$, then $\log|s_E|=0$, and the assertion follows from
$\psi\leq0$, \Cref{prop:ch35-homogeneous-normalization}.  Assume henceforth
that $E\neq0$.  The implication $\Leftarrow$ and uniqueness were proved
in \Cref{thm:ch35-homogeneous-siu}; the same theorem proves
$\Rightarrow$ whenever
$\sup_{K_E}\psi<-1$.  It remains to identify precisely the input used at
the endpoint
\begin{equation}
\label{eq:ch35-siu-critical-sup}
 \sup_{K_E}\psi=-1.
\end{equation}
That endpoint is \cite[Theorem~6.21]{BJ22}, and we use it as a foundational
input.  Here is the complete reconstruction of the approximation appearing
in its published proof.

Choose, by \Cref{cor:ch35-homogenization-psh}, a decreasing net
$\psi_a\in\cH_{\mathrm{hom}}(L_a)$ such that
$c_1(L_a)\to\theta$ and $\psi_a\searrow\psi$.  Dini's lemma on the
compact set $K_E$ gives rational numbers $t_a>0$, after passage to a
cofinal subnet, such that
\begin{equation}
\label{eq:ch35-siu-source-ta}
 t_a\longrightarrow1,
 \qquad
 \sup_{K_E}\psi_a\leq-t_a.
\end{equation}
By \Cref{lem:ch35-slice-to-cone,lem:ch35-fs-division},
\begin{equation}
\label{eq:ch35-siu-source-residual}
 \rho_a:=\psi_a-t_a\log|s_E|
 \in\cH_{\mathrm{hom}}(L_a-t_aE),
 \qquad
 c_1(L_a)-t_a[E]\longrightarrow\theta-[E].
\end{equation}
Each residual in \eqref{eq:ch35-siu-source-residual} is defined by
regular sections as in \Cref{lem:ch35-fs-division}; the displayed
subtraction is an equality on $X^{\val}$. On $X^{\divis}$ one has
$\rho_a\to\psi-\log|s_E|$, and this limit vanishes at
$v_{\triv}$, hence is generically finite.  The remaining assertion of
\cite[Theorem~6.21]{BJ22} is exactly that the residual can be admitted by
the decreasing-closure axiom in this endpoint situation; applying that
assertion to \eqref{eq:ch35-siu-source-residual} yields
$\psi-\log|s_E|\in\PSH_{\mathrm{hom},X}(\theta-[E])$.  Divisorial
determination then gives the identity on all of $X^{\An}$, as in Step~4
of \Cref{thm:ch35-homogeneous-siu}.

The last paragraph is deliberately not presented as a consequence of
\eqref{eq:ch35-siu-source-ta} alone.  Indeed, for $b\geq a$, monotonicity
of the residuals would require
\begin{equation}
\label{eq:ch35-siu-residual-monotonicity-test}
 t_b-t_a
 \leq\inf_{K_E}(\psi_a-\psi_b),
\end{equation}
which is not implied by \eqref{eq:ch35-siu-source-ta}.  Thus the
foundational input is the full endpoint of \cite[Theorem~6.21]{BJ22}, not
an unproved claim in this book that arbitrary choices of $t_a$ make
$(\rho_a)$ decreasing.  After the envelope property has independently
been established, the endpoint also follows from the fixed-$t<1$
argument and increasing-envelope closure; this later proof is written in
\cite[Proposition~17]{BJGreen24} under its smooth characteristic-zero
hypotheses.  It cannot replace the present foundational input in the proof
of the envelope theorem.
\end{proof}

\begin{remark}[Where the borderline case is delicate]
\label{rem:ch35-siu-boundary}
Write $u:=-\sup_{K_E}\psi\geq1$.  If $u>1$, the proof of
\Cref{thm:ch35-homogeneous-siu} applies verbatim.  If $u=1$, Dini's
theorem still gives $\sup_{K_E}\psi_a\to-1$, but possibly with
$\sup_{K_E}\psi_a>-1$ for every $a$; the source approximation then
divides by $\log|s_{t_aE}|$, with rational $t_a<1$ tending to $1$,
and produces the residual family
\[
 \psi'_a=\psi_a-t_a\log|s_E|=\psi_a+t_a\,v(E)
\]
as a sum of a decreasing and an increasing family.  The exact obstruction
is \eqref{eq:ch35-siu-residual-monotonicity-test}; it is why the endpoint is
imported as the full theorem rather than derived from Dini's lemma alone.
Equivalently, one may
run the argument for each fixed rational $t<1$, obtaining
\[
 \psi-t\log|s_E|\in\PSH_{\mathrm{hom},X}(\theta-t[E])
 \qquad(0<t<1,\ t\in\mathbf Q),
\]
but this family \emph{increases} to $\psi-\log|s_E|$ as $t\uparrow1$,
and increasing limits are exactly what the envelope property governs.  The
one-dimensional case shows that the conclusion is nevertheless correct: on
$X=\mathbf P^1$ a homogeneous $\theta$-psh function is recorded by its
slopes $c_r\geq0$ at the closed points, subject to
$\sum_rc_r\leq\deg\theta$ with at most countably many nonzero terms
(\Cref{ex:ch35-p1-hom-pl} and \Cref{prop:ch35-curve-classification}), and
the inequality $\psi\leq\log|s_E|$ with $E=\sum_ra_rr$ reads
$c_r\geq a_r$; the residual slopes $c_r-a_r$ are nonnegative and sum to
at most $\deg\theta-\deg E$, which is exactly membership in
$\PSH_{\mathrm{hom},X}(\theta-[E])$.  The strict theorem
\Cref{thm:ch35-homogeneous-siu} is the part proved entirely from the closure
calculus developed in this book.  The general endpoint
\Cref{thm:ch35-homogeneous-siu-general} is genuinely needed in
\Crefrange{chap:bj-cartier-realization}{chap:bj-subintegral-normalization}:
\Cref{ex:ch38-strict-room-fails,ex:ch39-reflection-borderline} exhibit the
failure of strict room in the actual reflection construction.  Those uses
therefore inherit the foundational input \cite[Theorem~6.21]{BJ22}; they are
not applications of the strict proof under a hidden stronger hypothesis.
\end{remark}

\begin{warning}{Only the homogeneous statement is unconditional}
\Cref{thm:ch35-homogeneous-siu} is stated for \emph{homogeneous}
$\theta$-psh functions and requires no envelope hypothesis.  The
corresponding statement for an arbitrary $\phi\in\PSH_X(\theta+[E])$,
namely
\[
 \phi\leq\log|s_E|+O(1)\iff\phi-\log|s_E|\in\PSH_X(\theta),
\]
is \emph{not} unconditional: its proof passes through
\eqref{eq:legendre-reconstruction-bj} and requires the envelope property for
$\theta$ in order to know that the reassembled supremum is psh.  That
version is \cite[Theorem~6.31]{BJ22}; it belongs among the deductions from
the envelope theorem collected in
\Cref{chap:bj-envelope-consequences}, not among its inputs.  Confusing the
two is the commonest way to produce a circular argument in this part of the
book.
\end{warning}

\section{Cartier and Weil b-divisors}
\label{sec:ch35-bdivisors}

Assume $X$ normal.  A model of $X$ is a normal projective variety
$\pi:Y\to X$ birational over $X$.  Models form a cofiltered system under
domination.

The word \emph{cofiltered} is a theorem, not a convention: any two models
are dominated by a third, by
\Cref{lem:appF-common-domination}, and the proof there uses no resolution of
singularities.  This is what allows every construction below to be performed
after passing to a single sufficiently high model.

\begin{definition}[b-divisors]
\label{def:algebraic-bdivisors}
A real Weil b-divisor $B$ over $X$ is a compatible family
\begin{equation}
\label{eq:weil-bdivisor}
 B=(B_Y)_Y,\qquad B_Y\in\Div(Y)_{\mathbf R},\qquad
 \mu_*B_{Y'}=B_Y
\end{equation}
for every domination $\mu:Y'\to Y$.  It is Cartier if it is determined on
one model $Y$: $B_{Y'}=\mu^*B_Y$ on every higher model.
\end{definition}

Thus Weil b-divisors form an inverse limit under pushforward, while Cartier
b-divisors form a direct limit under pullback.  A Cartier b-divisor is nef if
its trace on one, equivalently every, determination is nef.  It is
relatively semiample if the trace is semiample over $X$.

These are \Cref{def:appF-weil-bdivisor} and
\Cref{def:appF-cartier-bdivisor}, and the two limit descriptions are
\eqref{eq:appF-direct-inverse}.  Independence of the determination in the
nef case is \Cref{lem:appF-nef-determination}.  We write
$\operatorname{Div}_{\mathbf R}(\mathfrak X)$ and
$\operatorname{CDiv}_{\mathbf R}(\mathfrak X)$ for the two spaces, and
$\overline F$ for the Cartier b-divisor determined on $X$ itself by an
$\mathbf R$-Cartier divisor $F$ on $X$.  Ordering is coefficientwise:
$B\leq0$ means $B_Y\leq0$ for every model $Y$, which by
\Cref{prop:appF-coefficients} is the same as
$\operatorname{coeff}_E(B)\leq0$ for every prime divisor $E$ over
$X$.  For a Cartier b-divisor determined on $Y$, it suffices to check
$B_Y\leq0$, because pullback of an antieffective Cartier divisor is
antieffective.

\subsection{The evaluation pairing}
\label{sec:ch35-evaluation-pairing}

Each b-divisor has a coefficient function on divisorial points.
Here homogeneity on $X^{\divis}$ means compatibility with positive
rational rescaling, which preserves the rational divisorial-point locus.
The continuous Cartier extension to $X^{\val}$ will be homogeneous under
all positive real rescalings. The pairing is
\begin{equation}
\label{eq:bdivisor-coefficient-function}
                  \psi_B(t\ord_E):=t\,\ord_E(B_Y),
\end{equation}
where $Y$ carries $E$.  Pushforward compatibility makes this
well-defined, and it is homogeneous.  Coefficientwise convergence of
b-divisors is exactly pointwise convergence on $X^{\divis}$.

We now prove these three assertions and identify the image of the pairing
exactly.  Recall that a prime divisor \emph{over} $X$ is a prime divisor
$E\subset Y$ on some model $Y$, two such being identified when they
induce the same valuation $\ord_E$ of $k(X)$; this identification is
consistent with \Cref{prop:appF-coefficients}.

\begin{proposition}
\label{prop:ch35-weil-coefficient-bijection}
Let $X$ be normal and projective.  The assignment
$B\mapsto\psi_B$ of \eqref{eq:bdivisor-coefficient-function} is a
well-defined $\mathbf R$-linear bijection from
$\operatorname{Div}_{\mathbf R}(\mathfrak X)$ onto the set of functions
\begin{equation}
\label{eq:ch35-locally-finite-homogeneous}
 \mathcal W:=\left\{\psi:X^{\divis}\to\mathbf R
 \ \middle|\
 \begin{array}{l}
 \psi\ \text{homogeneous},\ \psi(v_{\triv})=0,\ \text{and for every}\\
 \text{model }Y\ \text{only finitely many prime divisors}\\
 E\subset Y\ \text{have}\ \psi(\ord_E)\neq0
 \end{array}
 \right\}.
\end{equation}
It is order preserving in the sense that $B\leq B'$ if and only if
$\psi_B\leq\psi_{B'}$, and a net $B_i$ converges to $B$ in the inverse
limit topology of \eqref{eq:appF-direct-inverse} if and only if
$\psi_{B_i}\to\psi_B$ pointwise on $X^{\divis}$.
\end{proposition}

\begin{proof}
\emph{Well-definedness.}  Let $v=t\ord_E\in X^{\divis}$ with $t>0$ and
$E$ a prime divisor on a model $Y$.  If $E$ also occurs on a model
$Y'$, choose a common domination $Y''\succeq Y,Y'$, which exists by
\Cref{lem:appF-common-domination}.  The strict transform of $E$ on
$Y''$ is a prime divisor $\widetilde E$ with
$\ord_{\widetilde E}=\ord_E$, and pushforward to $Y$ and to $Y'$
preserves the coefficient of $\widetilde E$ by
\Cref{prop:appF-coefficients}.  Hence
$\ord_E(B_Y)=\operatorname{coeff}_{\widetilde E}(B_{Y''})=\ord_E(B_{Y'})$,
so \eqref{eq:bdivisor-coefficient-function} does not depend on the model.
At $v_{\triv}$ we set $\psi_B(v_{\triv})=0$, consistently with the case
$t=0$.

\emph{Homogeneity and local finiteness.}  Homogeneity is built into
\eqref{eq:bdivisor-coefficient-function}: $\psi_B(s\cdot t\ord_E)
=st\,\ord_E(B_Y)=s\,\psi_B(t\ord_E)$.  Local finiteness is the statement
that $B_Y$ is a divisor, that is, a locally finite --- here finite, since
$Y$ is noetherian --- sum of prime divisors.

\emph{Injectivity and surjectivity.}  If $\psi_B=0$, then every
coefficient of every trace vanishes, so $B=0$; linearity is clear.
Conversely, given $\psi\in\mathcal W$, define
$B_Y:=\sum_{E\subset Y}\psi(\ord_E)\,E$, a finite sum by hypothesis.  For a
domination $\mu:Y'\to Y$, pushforward kills $\mu$-exceptional primes and
preserves the coefficient of a prime whose strict transform is taken, so
$\mu_*B_{Y'}=B_Y$; this is exactly the computation in
\Cref{prop:appF-coefficients}.  Hence $B=(B_Y)_Y$ is a Weil b-divisor with
$\psi_B=\psi$.

\emph{Order and topology.}  Both statements are coefficientwise
reformulations: $B\leq B'$ means
$\operatorname{coeff}_E(B)\leq\operatorname{coeff}_E(B')$ for every prime
$E$ over $X$, which by homogeneity is
$\psi_B\leq\psi_{B'}$ on $X^{\divis}$; and the inverse limit topology on
$\varprojlim_Y\operatorname{Div}_{\mathbf R}(Y)$ is the topology of
pointwise convergence of all coefficients, each $\operatorname{Div}_{\mathbf R}(Y)$ carrying the coefficientwise
subspace topology from $\mathbf R^{\{E\subset Y\}}$. The divisor space
itself is not finite dimensional, and a converging net need not have
one common finite support.
\end{proof}

\begin{warning}{The pairing is defined on $X^{\divis}$ only}
\Cref{prop:ch35-weil-coefficient-bijection} produces a function on
$X^{\divis}$.  For a general Weil b-divisor there is no canonical
extension of $\psi_B$ to $X^{\An}$, and none is used.  For a
\emph{Cartier} b-divisor there is a continuous real-valued extension to
$X^{\val}$, by \Cref{prop:ch35-cartier-continuous-extension}, and for the coefficient
function of a homogeneous psh function there is one, by
\Cref{cor:ch35-divisorial-determination}.  These are three different
statements and the difference is exactly the content of
\Cref{ex:appF-non-cartier-bdivisor}: a Weil b-divisor supported on an
infinite chain of exceptional primes has a perfectly good coefficient
function and no continuous extension.
\end{warning}

\begin{proposition}
\label{prop:ch35-cartier-continuous-extension}
Let $C$ be a real Cartier b-divisor with determination
$\pi:Y\to X$.  Then $\psi_C$ extends uniquely to a continuous function
\begin{equation}
\label{eq:ch35-cartier-extension}
 \psi_C:X^{\val}\longrightarrow\mathbf R,
 \qquad
 \psi_C(v)=v(C_Y),
\end{equation}
where the right-hand side is computed on $Y^{\val}\simeq X^{\val}$, and
the resulting map
$\operatorname{CDiv}_{\mathbf R}(\mathfrak X)\to C^0(X^{\val})$ is an
injective $\mathbf R$-linear map.
\end{proposition}

\begin{proof}
A birational morphism of normal projective varieties induces a bijection
$Y^{\val}\to X^{\val}$, because both sides are the set of real valuations
of the common function field $k(X)=k(Y)$ admitting a centre, and
properness supplies the centre on either side; the bijection is a homeomorphism as follows. The map
$Y^{\An}\to X^{\An}$ is closed because its source is compact and its
target is Hausdorff. Its inverse image of $X^{\val}$ is $Y^{\val}$:
a proper subvariety of Y cannot dominate X under a birational morphism.
For every closed subset C of $Y^{\An}$, its image satisfies
$\pi(C\cap Y^{\val})=\pi(C)\cap X^{\val}$. Thus the restriction is
closed as well as continuous and bijective. No compactness of
$Y^{\val}$ is assumed.  Writing $C_Y$ locally as a
difference of two effective Cartier divisors, $v\mapsto v(C_Y)$ is a
difference of two ideal evaluations and is therefore continuous and real
valued on $Y^{\val}$, by
\Cref{prop:ch35-hom-pl-semiring}.  If $Y'\succeq Y$ is higher, then
$C_{Y'}=\mu^*C_Y$ and $v(\mu^*C_Y)=v(C_Y)$ for
$v\in Y'^{\val}=Y^{\val}$, because pullback of a Cartier divisor pulls
back local equations.  Hence \eqref{eq:ch35-cartier-extension} is
independent of the determination.  It restricts to
\eqref{eq:bdivisor-coefficient-function} on $X^{\divis}$, because
$\ord_E(C_Y)$ is by definition the coefficient of $E$ in
$C_{Y'}$ for any $Y'$ carrying $E$.  Uniqueness of the extension
follows from density of $X^{\divis}$ in $X^{\val}$, which is a
consequence of \Cref{thm:triv-divisorial-density}, and injectivity from
\Cref{prop:ch35-weil-coefficient-bijection}.
\end{proof}

\subsection{The sign convention}
\label{sec:ch35-sign}

The sign is worth fixing.  If $F\geq0$ is a Cartier divisor on $X$, the
section function satisfies
\begin{equation}
\label{eq:section-bdivisor-sign}
                    \log|s_F|(v)=-v(F),
\end{equation}
whereas the coefficient function of the b-divisor $\overline F$ is
$\psi_{\overline F}(v)=v(F)$.  Hence
$\log|s_F|=-\psi_{\overline F}$ on divisorial valuations.

This is \eqref{eq:appF-section-sign}, and it is the source of the
antieffectivity convention that pervades this chapter and
\Cref{chap:bj-cartier-realization}: a \emph{nonpositive} homogeneous
function corresponds to an \emph{antieffective} b-divisor.  Explicitly, if
$\psi\in\PL^+_{\mathrm{hom}}(X)$ then $\psi\leq0$, and the Cartier
b-divisor $B$ with $\psi_B=\psi$ satisfies $B\leq0$.  Three unit tests
fix the convention beyond doubt:
\begin{equation}
\label{eq:ch35-sign-tests}
\begin{aligned}
 \psi_{\overline F}(\ord_E)&=\ord_E(F)\geq0
 &&\text{for }F\geq0\text{ effective},\\
 \log|s_F|(\ord_E)&=-\ord_E(F)\leq0,\\
 \log|\mathfrak b|(\ord_E)&=-\ord_E(\mathfrak b)\leq0
 &&\text{for }\mathfrak b\subseteq\cO_X .
\end{aligned}
\end{equation}
The conventions sheet fixes the same sign in the nontrivially valued
setting: an effective vertical divisor gives a \emph{nonnegative} model
function, because $\cO(D)$ is generated by the inverse of a local
equation.  \Cref{app:normalizations} records the translation, and the two
conventions are consistent because
$\log|s_F|$ uses the canonical section of $\cO(F)$ whereas
$\varphi_D$ uses the inverse local generator of $\cO(D)$.

\subsection{Cartier b-divisors and homogeneous PL functions}
\label{sec:ch35-cartier-pl}

\begin{proposition}\label{prop:cartier-homogeneous-pl}
The coefficient functions of real Cartier b-divisors are precisely the real
homogeneous PL functions on $X^{\divis}$.  If $B\leq0$ is rational and
relatively semiample, then $\psi_B$ is the restriction of a homogeneous
Fubini--Study function.
\end{proposition}

\begin{proof}
On a common determination, write B as a difference of relatively
semiample divisors. Their pushforwards give coherent fractional ideals on X,
whose logarithms are differences of ordinary ideal functions by
\cite[Example~6.4]{BJ22}.  This
places $\psi_B$ in $\PL_{\mathrm{hom}}(X)$.  Conversely, clear
denominators in a difference of ideal functions and take normalized blowups;
the resulting Cartier divisors glue on a common domination.  Relative
semiampleness and antieffectivity yield global generators as in
\Cref{prop:hom-fs-ideal}.

We give the argument in full, treating the rational antieffective
relatively semiample case first, since it is the one used in
\Cref{chap:bj-cartier-realization}.

\medskip
\noindent\emph{Step 1: from $B\in\operatorname{CDiv}^+_{\mathbf Q}$ to
$\PL^+_{\mathrm{hom}}$.}
Let $B$ be a rational Cartier b-divisor with
$B\leq0$ and $B_Y$ relatively semiample for a determination
$\pi:Y\to X$.  Choose $m\in\mathbf Z_{>0}$ with $mB_Y$ an integral
Cartier divisor and $\cO_Y(mB_Y)$ globally generated over $X$, that is,
$\pi^*\pi_*\cO_Y(mB_Y)\to\cO_Y(mB_Y)$ surjective.  Since
$mB_Y\leq0$, the sheaf $\cO_Y(mB_Y)$ is a subsheaf of $\cO_Y$, and it
is generically trivial because every divisor is trivial at the generic
point.  Put
$\mathfrak b:=\pi_*\cO_Y(mB_Y)\subseteq\pi_*\cO_Y=\cO_X$, the last
equality by Zariski's main theorem, $X$ being normal and $\pi$
birational and proper.  Then $\mathfrak b$ is a nonzero coherent ideal
with $\mathfrak b\cdot\cO_Y=\cO_Y(mB_Y)$, by global generation, and hence
for $v\in X^{\val}$,
\[
 v(\mathfrak b)=v\bigl(\mathfrak b\cdot\cO_Y\bigr)=-v(mB_Y)=-m\,\psi_B(v).
\]
Therefore $\psi_B=m^{-1}\log|\mathfrak b|\in\PL^+_{\mathrm{hom}}(X)$.

\medskip
\noindent\emph{Step 2: the converse in the antieffective semiample case.}
Let $\psi=m^{-1}\log|\mathfrak b|\in\PL^+_{\mathrm{hom}}(X)$ with
$0\neq\mathfrak b\subseteq\cO_X$.  Let $\pi:Y\to X$ be the normalized
blowup of $\mathfrak b$, so that
$\mathfrak b\cdot\cO_Y=\cO_Y(-E)$ for an effective Cartier divisor $E$
on $Y$, with $-E$ relatively ample, hence relatively semiample.  Let
$B$ be the Cartier b-divisor determined on $Y$ by
$B_Y:=-m^{-1}E\leq0$.  For $v\in X^{\val}$,
\[
 \psi_B(v)=v(B_Y)=-m^{-1}v(E)=-m^{-1}v(\mathfrak b)=\psi(v).
\]
Thus $\psi=\psi_B$ with $B$ rational, antieffective, and relatively
semiample.

\medskip
\noindent\emph{Step 3: the general Cartier case.}
First let $B$ be rational, with determination $\pi:Y\to X$.  Choose a
relatively ample Cartier divisor $A$ on $Y$.  For a sufficiently large
integer $q$, both
\[
 B_Y^{(1)}:=B_Y+qA,
 \qquad B_Y^{(2)}:=qA
\]
are relatively ample, hence relatively semiample, and
$B_Y=B_Y^{(1)}-B_Y^{(2)}$.  After replacing both by a common positive
multiple, put
\[
 \mathfrak b_j:=\pi_*\cO_Y(mB_Y^{(j)})\subset K(X),\qquad j=1,2.
\]
These are nonzero coherent fractional ideals. Their logarithms belong to
$\PL_{\mathrm{hom}}(X)$ by the conductor-ideal argument of
\cite[Example~6.4]{BJ22}. Relative global
generation gives
$\mathfrak b_j\cO_Y=\cO_Y(mB_Y^{(j)})$, hence
\[
 \psi_B=m^{-1}\bigl(\log|\mathfrak b_1|
                    -\log|\mathfrak b_2|\bigr)
 \quad\text{on }X^{\val}.
\]
Thus $\psi_B\in\PL_{\mathrm{hom}}(X)$.  Tensoring this rational
identification with $\mathbf R$ gives the assertion for real Cartier
b-divisors.  Conversely, a real homogeneous PL function is a finite real
linear combination of ideal functions; Step 2 realizes the ideal functions
by rational Cartier b-divisors, and a common determination realizes their
linear combination.  This is the precise point at which fractional ideals
are needed: the two semiample summands are not asserted to be
antieffective.

\medskip
\noindent\emph{Step 4: the Fubini--Study assertion.}
If $B\leq0$ is rational and relatively semiample, Step 1 gives
$\psi_B=m^{-1}\log|\mathfrak b|$.  By
\Cref{prop:hom-fs-ideal}, this lies in
$\cH_{\mathrm{hom}}(L)$ for any $\mathbf Q$-line bundle $L$ with
$mL$ a line bundle and $mL\otimes\mathfrak b$ globally generated, and
such $L$ exists, for instance any sufficiently positive multiple of an
ample bundle.  This is \cite[Theorem~6.32]{BJ22}.
\end{proof}

\begin{theorem}
\label{thm:ch35-semiample-criterion}
Let $X$ be normal, $L\in\Pic(X)_{\mathbf Q}$, and $B$ a rational
Cartier b-divisor with determination $\pi:Y\to X$.  Then
\begin{equation}
\label{eq:ch35-semiample-criterion}
 \psi_B\in\cH_{\mathrm{hom}}(L)
 \iff
 B\leq0\ \text{ and }\ \pi^*L+B_Y\ \text{ is semiample on }Y .
\end{equation}
The right-hand condition is independent of the determination.
\end{theorem}

\begin{proof}
\emph{Independence of the determination.}  If $\mu:Y'\to Y$ is a higher
model, then $\mu^*(\pi^*L+B_Y)=(\pi\mu)^*L+B_{Y'}$, and semiampleness is
preserved by pullback along a projective birational morphism and detected by
it, because a globally generated multiple pulls back to a globally generated
multiple and pushforward along $\mu$ with $\mu_*\cO_{Y'}=\cO_Y$ recovers
the sections.

\emph{The implication $\Rightarrow$.}  If
$\psi_B\in\cH_{\mathrm{hom}}(L)$, then $\psi_B\leq0$ by
\Cref{prop:ch35-homogeneous-normalization}, hence $B\leq0$.  By
\Cref{prop:hom-fs-ideal} there are $m\geq1$ and a nonzero ideal
$\mathfrak b\subseteq\cO_X$ with $mL$ a line bundle,
$mL\otimes\mathfrak b$ globally generated, and
$\psi_B=m^{-1}\log|\mathfrak b|$.  Let $\rho:Z\to X$ be the normalized
blowup of $\mathfrak b$, with $\mathfrak b\cdot\cO_Z=\cO_Z(-E)$.  By
Step 2 of the proof of \Cref{prop:cartier-homogeneous-pl},
$B$ is determined on $Z$ by $-m^{-1}E$.  Global generation of
$mL\otimes\mathfrak b$ says exactly that
$\rho^*(mL)\otimes\cO_Z(-E)=m(\rho^*L+B_Z)$ is globally generated, hence
$\rho^*L+B_Z$ is semiample.  By independence of the determination the
same holds on $Y$.

\emph{The implication $\Leftarrow$.}  Assume $B\leq0$ and
$\pi^*L+B_Y$ semiample.  Choose $m\in\mathbf Z_{>0}$ with
$m(\pi^*L+B_Y)$ a globally generated line bundle and $mB_Y$ integral
Cartier.  The projection formula gives
\[
 H^0\bigl(Y,m(\pi^*L+B_Y)\bigr)
 \simeq H^0\bigl(X,mL\otimes\pi_*\cO_Y(mB_Y)\bigr).
\]
Let $mL\otimes\mathfrak b_m$ be the image of the resulting evaluation
map into
$mL\otimes\pi_*\cO_Y(mB_Y)\subseteq mL$, where the inclusion uses
$mB_Y\leq0$ and $\pi_*\cO_Y=\cO_X$.  Thus
$\mathfrak b_m\subseteq\cO_X$ is the coherent ideal locally generated by
the coefficients of these global sections, and
$mL\otimes\mathfrak b_m$ is globally generated by construction.
Pulling the evaluation map to $Y$, global generation of
$m(\pi^*L+B_Y)$ gives
$\mathfrak b_m\cdot\cO_Y=\cO_Y(mB_Y)$, which gives
\[
 \psi_B(v)=v(B_Y)=-m^{-1}v(\mathfrak b_m)=m^{-1}\log|\mathfrak b_m|(v)
\]
for $v\in X^{\val}$.  By \Cref{prop:hom-fs-ideal},
$\psi_B\in\cH_{\mathrm{hom}}(L)$.  This is
\cite[Lemma~6.34]{BJ22}.
\end{proof}

\begin{warning}{Semiample, not nef}
\Cref{thm:ch35-semiample-criterion} characterizes membership in the
Fubini--Study cone by \emph{semiampleness} of the total class, an
algebro-geometric condition on an actual line bundle.  Replacing semiample
by nef changes the statement and its proof: the characterization of
$\PSH_{\mathrm{hom},X}(\theta)$ by nefness of $\theta+B$ is
\Cref{thm:ch35-cartier-psh-criterion}, and its proof uses the envelope
property for ample classes.  The gap between the two is the usual gap
between a linear-series condition and a numerical one, and it is the reason
\Cref{chap:bj-cartier-realization} inserts an ample buffer
\eqref{eq:bj-buffer-path} before invoking any Fubini--Study statement.
\end{warning}

\subsection{Worked examples}
\label{sec:ch35-bdivisor-examples}

\begin{proposition}
\label{prop:ch35-curve-classification}
Let $X$ be a smooth connected projective curve over $k$ and
$\theta\in N^1(X)=\mathbf R$, identified with $\deg\theta$.  Then
$\psi:X^{\An}\to\mathbf R\cup\{-\infty\}$ belongs to
$\PSH_{\mathrm{hom},X}(\theta)$ if and only if there are numbers
$c_r\geq0$, indexed by the closed points $r\in X(k)$, all but countably
many zero, with
\begin{equation}
\label{eq:ch35-curve-classification}
 \sum_{r\in X(k)}c_r\leq\deg\theta,
 \qquad
 \psi(s\,\ord_r)=-c_rs\quad(s\geq0),
\end{equation}
and $\psi(v_{r,\triv})=-\infty$ when $c_r>0$, $\psi(v_{r,\triv})=0$
when $c_r=0$.  Only countability may be asserted: it is what
$\sum_rc_r\leq\deg\theta$ forces, and it cannot be improved.  On
$X=\mathbf P^1$ with $\theta=c_1(\cO(1))$, pick distinct closed points
$r_1,r_2,\dots$, set $c_n=2^{-n}$, and let
$\psi_N:=\sum_{n\leq N}2^{-n}\log|\mathfrak m_{r_n}|
=2^{-N}\log|\mathfrak b_N|$ with
$\mathfrak b_N=\prod_{n\leq N}\mathfrak m_{r_n}^{2^{N-n}}$: then
$\deg\bigl(\cO(2^N)\otimes\mathfrak b_N\bigr)=2^N-(2^N-1)=1$, so
$\psi_N\in\cH^{\mathrm{gf}}_{\mathbf Q}(\cO(1))$, the $\psi_N$ decrease
pointwise to a $\psi$ with $\psi(v_{\triv})=0$, and
$\sum_nc_n=1=\deg\theta$.  This $\psi$ has infinitely many nonzero
slopes.

The coefficient function of a b-divisor requires more: $\psi$ is of the
form $B=-\sum_rc_r\overline r$ for a (necessarily Cartier) b-divisor
precisely when only finitely many $c_r$ are nonzero, local finiteness on
the single model $X$ being exactly that condition.  So on a curve
$\PSH_{\mathrm{hom},X}(\theta)$ is strictly larger than the set of
coefficient functions of b-divisors.  Every b-divisor over $X$ is Cartier,
determined on $X$ itself.
\end{proposition}

\begin{proof}
Every normal projective variety birational to a smooth projective curve is
that curve, so $X$ is the unique model and every b-divisor is Cartier,
determined on $X$; this proves the last assertion and shows that
$\operatorname{Div}_{\mathbf R}(\mathfrak X)=\Div(X)_{\mathbf R}$.  By
\Cref{prop:trivial-p1-tree} and its analogue for an arbitrary smooth curve,
recorded in \Cref{prop:ch08-trivial-curve-scaling}, the space
$X^{\An}$ is the union of the rays $\{s\,\ord_r\}_{s\in[0,\infty]}$ glued
at $v_{\triv}$, and $X^{\divis}$ consists of $v_{\triv}$ together with
the points $s\,\ord_r$ with $s\in\mathbf Q_{>0}$.  A homogeneous
function is therefore determined by the numbers
$c_r:=-\psi(\ord_r)$, and $c_r\geq0$ because
$\psi\leq0$ by \Cref{prop:ch35-homogeneous-normalization}.

Use the curve criterion recorded in \Cref{sec:ch34-pl-positivity}, from
\cite[Example~4.13 and Proposition~4.12]{BJ22}: for an upper-semicontinuous
function finite at the root, with each closed endpoint equal to its radial
limit, psh membership is equivalent to convexity of
$s\mapsto\psi(s\,\ord_r)$ on $[0,+\infty)$ for every $r$, together with
\[
 \deg\theta+\sum_{r\in X(k)}\frac{\rd}{\rd s}\Big|_{s=0^+}\psi(s\,\ord_r)
 \geq0 .
\]
For a homogeneous $\psi$ the functions $s\mapsto-c_rs$ are linear, hence
convex, with right derivative $-c_r$ at $s=0$; the displayed condition
becomes $\sum_rc_r\leq\deg\theta$, and it forces all but countably many
$c_r$ to vanish.  Countability is all that follows, and it is sharp: the
example in the statement has infinitely many nonzero $c_r$.  Finiteness is
strictly stronger, and it is exactly the condition that $\psi$ be the
coefficient function of a b-divisor, local finiteness on the single model
$X$ being precisely the requirement that only finitely many $c_r$ be
nonzero; that is the second paragraph of the statement, and it is where the
two classes differ.  Conversely, the stated function is upper semicontinuous: it is the
infimum over finite sets $F$ of
$\sum_{r\in F,\,c_r>0}c_r\log|\mathfrak m_r|$, each upper
semicontinuous. It has root value zero and the prescribed radial endpoint
values. The curve criterion therefore proves sufficiency. Finally the value at the endpoint
$v_{r,\triv}=\lim_{s\to\infty}s\,\ord_r$ is
$\lim_{s\to\infty}(-c_rs)$, which is $-\infty$ if $c_r>0$ and $0$ if
$c_r=0$; this is consistent with
\Cref{lem:ch35-homogeneity-reformulations}, the endpoints being fixed points
of the scaling action.
\end{proof}

\begin{example}[A homogeneous potential on the cone over
             $\mathbf P^1$, with its b-divisor]
\label{ex:ch35-p1-bdivisor}
Let $X=\mathbf P^1_k$ with homogeneous coordinates $x,y$, let
$p=[0:1]$, $q=[1:0]$, $r_0=[1:1]$, and let
$\theta=c_1(\cO(3))$, so $\deg\theta=3$.  Consider
\begin{equation}
\label{eq:ch35-p1-example-psi}
 \psi:=\tfrac32\log|\mathfrak m_p|+\tfrac12\log|\mathfrak m_q|
 +\tfrac12\log|\mathfrak m_{r_0}| .
\end{equation}
By \Cref{prop:ch35-hom-pl-semiring} this is a homogeneous PL function; by
\Cref{prop:ch35-curve-classification} it is $\theta$-psh, because its
slopes are
\[
 c_p=\tfrac32,\qquad c_q=\tfrac12,\qquad c_{r_0}=\tfrac12,
 \qquad c_r=0\ \text{otherwise},
\]
and $\tfrac32+\tfrac12+\tfrac12=\tfrac52\leq3$.  Its values are
\[
 \begin{aligned}
 \psi(s\,\ord_p)&=-\tfrac32s,
 &\psi(s\,\ord_q)&=-\tfrac12s,\\
 \psi(s\,\ord_{r_0})&=-\tfrac12s,
 &\psi(s\,\ord_r)&=0\quad(r\neq p,q,r_0),
 \end{aligned}
\]
with $\psi(v_{\triv})=0$ and
$\psi(v_{p,\triv})=\psi(v_{q,\triv})=\psi(v_{r_0,\triv})=-\infty$.

Its b-divisor is the Cartier b-divisor determined on $X$ by
\begin{equation}
\label{eq:ch35-p1-example-B}
 B=-\tfrac32\,\overline p-\tfrac12\,\overline q
   -\tfrac12\,\overline{r_0},
 \qquad
 B_X=-\tfrac32p-\tfrac12q-\tfrac12r_0,
\end{equation}
an antieffective real Cartier b-divisor of degree $-\tfrac52$.  The total
class is
\[
 \theta+B_X=3-\tfrac52=\tfrac12>0,
\]
so $\theta+B$ is nef --- on a curve, nef means nonnegative degree ---
in accordance with \Cref{thm:ch35-cartier-psh-criterion}.  Since
$\tfrac12>0$, the total class is in fact ample, so
$\psi\in\cH_{\mathrm{hom}}(L)$ for suitable $L$, by
\Cref{thm:ch35-semiample-criterion}: taking $m=2$ and
$L=\cO(3)$, the ideal $\mathfrak b=\mathfrak m_p^3\mathfrak m_q
\mathfrak m_{r_0}$ has $\deg\cO(6)\otimes\mathfrak b=6-5=1\geq0$ and
$\cO(6)\otimes\mathfrak b\simeq\cO(1)$ is globally generated, so
$\psi=\tfrac12\log|\mathfrak b|\in\cH_{\mathrm{hom}}(\cO(3))$.

Finally the boundary case: if instead $c_p=2$, $c_q=\tfrac12$,
$c_{r_0}=\tfrac12$, the slopes sum to $3=\deg\theta$, the total class
$\theta+B_X$ is $0$, still nef but no longer ample, and the associated
ideal $\mathfrak b=\mathfrak m_p^4\mathfrak m_q\mathfrak m_{r_0}$ satisfies
$\cO(6)\otimes\mathfrak b\simeq\cO_X$, which is still globally generated;
so $\psi$ remains a homogeneous Fubini--Study function.  If one increases
$c_p$ to $\tfrac52$, the sum exceeds $3$, the total class becomes
negative, and $\psi$ is no longer $\theta$-psh: on the model $X$ the
class $\theta+B_X$ has degree $-\tfrac12<0$ and no globally generated
multiple.
\end{example}

\begin{example}[A b-divisor genuinely spread over a model]
\label{ex:ch35-blowup-bdivisor}
Let $X=\mathbf P^2_k$, let $q\in X$ be a closed point, and let
$\pi:Y\to X$ be the blowup at $q$, with exceptional curve $E$.  Take
$\psi:=\log|\mathfrak m_q|$, which lies in
$\cH_{\mathrm{hom}}(\cO(1))$ by \Cref{prop:hom-fs-ideal}, since
$\cO(1)\otimes\mathfrak m_q$ is globally generated by the two linear forms
vanishing at $q$.  The normalized blowup of $\mathfrak m_q$ is
$\pi:Y\to X$ and $\mathfrak m_q\cdot\cO_Y=\cO_Y(-E)$; so by Step 2 of
the proof of \Cref{prop:cartier-homogeneous-pl}, the corresponding Cartier
b-divisor is
\[
 B\ \text{determined on }Y\text{ by }B_Y=-E,
 \qquad B_X=\pi_*(-E)=0 .
\]
Thus $B$ is a nonzero Cartier b-divisor whose trace on $X$ vanishes: it
is \emph{exceptional} in the sense of
\Cref{def:ch35-exceptional-bdivisor}.  Its coefficient function is
\[
 \psi_B(t\,\ord_E)=-t,
 \qquad
 \psi_B(\ord_F)=0\ \text{ for every prime divisor }F\subset X,
\]
and more generally $\psi_B(v)=-v(\mathfrak m_q)$ for
$v\in X^{\val}$, which is the multiplicity function of the point $q$.
The total class is
$\pi^*c_1(\cO(1))-[E]=H-E$ in the notation of
\Cref{ex:appF-blowup}, which is nef with
$(H-E)^2=1-1=0$; it is semiample but not ample.  It is \emph{not} the
pullback of $\cO(1)$ from $\mathbf P^2$, which is $H$: since
$h^0(Y,H-E)=h^0(\mathbf P^2,\cO(1)\otimes\mathfrak m_q)=2$, the linear
system $|H-E|$ is the pencil of strict transforms of the lines through
$q$ and defines a morphism $f:Y\to\mathbf P^1$ with
$H-E=f^*\cO_{\mathbf P^1}(1)$; this is what
$(H-E)^2=0$ records: a class pulled back along a morphism to a \emph{curve}
squares to zero, in contrast to $H=\pi^*\cO(1)$, which is pulled back along
the birational morphism $\pi$ and has $H^2=1$.  The class on the boundary of the nef cone
is $H-E$ itself; $E$ is not nef at all, having $(E^2)=-1$.  Consistently,
\Cref{thm:ch35-semiample-criterion} certifies
$\psi_B\in\cH_{\mathrm{hom}}(\cO(1))$.

The example also shows why the model matters: $B$ is \emph{not} determined
on $X$.  A determination must carry $E$, and $E$ does not appear on
$X$.
\end{example}

\section{Nef b-divisor classes}
\label{sec:ch35-nef-bclasses}

\subsection{The Cartier criterion}
\label{sec:ch35-cartier-criterion}

The proof of the criterion will use the following closedness statement.

\begin{lemma}
\label{lem:ch35-nef-limit}
Let $X$ be normal and projective, let $(B_i)$ be a decreasing net of
rational Cartier b-divisors converging pointwise on $X^{\divis}$ to a
Cartier b-divisor $B$, and let $\theta_i\to\theta$ in $N^1(X)$.  If
$B_i\leq0$ and $\theta_i+B_i$ is nef for every $i$, then
$\theta+B$ is nef.
\end{lemma}

\begin{proof}
Let $\pi:Y\to X$ determine $B$.  It suffices to show
$(\pi^*\theta+B_Y)\cdot\Gamma\geq0$ for every integral curve
$\Gamma\subset Y$.

If $n=1$, every normal birational model is isomorphic to $X$, and the
claim is simply closedness of the nonnegative half-line in
$N^1(X)\simeq\mathbf R$.  Assume $n\geq2$.  Let $\mu:Z\to Y$ be the
normalized blowup of the ideal of $\Gamma$, and write
$\mathfrak a_\Gamma\cO_Z=\cO_Z(-F)$.  Thus $F$ is an effective Cartier
divisor; when the blowup is trivial, this notation means the Cartier
pullback of $\Gamma$, not a nonexistent exceptional divisor.  Fix an
ample class $A$ on $Z$.  The curve-testing formula used in
\cite[Lemma~6.37]{BJ22} gives
$\mu_*(F\cdot A^{n-2})=a\,\Gamma$ in $N_1(Y)$ for some $a>0$, where
$n=\dim X$.  By the projection formula
\eqref{eq:appF-projection}, it therefore suffices to prove
\[
 (\mu^*\pi^*\theta+B_Z)\cdot F\cdot A^{n-2}\geq0 .
\]
For each $i$ we claim
$(\mu^*\pi^*\theta_i+B_{i,Z})\cdot F\cdot A^{n-2}\geq0$.  The reason is not
that $B_{i,Z}$ is a pullback: $Z$ was chosen to dominate $Y$, a
determination of the \emph{limit} $B$, and nothing makes it dominate a
determination of $B_i$, whose trace on $Z$ is then a pushforward --- and
pushforward does not preserve nefness.  Instead, choose by
\Cref{lem:appF-common-domination} a model $\rho:W\to Z$ which also
dominates a determination of $B_i$.  On $W$ the trace $B_{i,W}$
\emph{is} the pullback of that determination, so
$\rho^*\mu^*\pi^*\theta_i+B_{i,W}$ is nef by
\Cref{lem:appF-nef-determination}; and $\rho^*F\cdot(\rho^*A)^{n-2}$ is the
intersection of an effective divisor with nef classes, whence
\[
 \bigl(\rho^*\mu^*\pi^*\theta_i+B_{i,W}\bigr)\cdot\rho^*F\cdot
 (\rho^*A)^{n-2}\geq0 .
\]
Applying \eqref{eq:appF-projection} three times, with $\rho_*[W]=[Z]$ and
$\rho_*B_{i,W}=B_{i,Z}$, rewrites the left-hand side as
$(\mu^*\pi^*\theta_i+B_{i,Z})\cdot F\cdot A^{n-2}$.

It remains to pass to the limit.  Since $B\leq B_i\leq0$ for all
$i$ --- the first inequality because the net is decreasing with limit
$B$ --- every $B_{i,Z}$ is supported on the finitely many prime
components $(E_\beta)$ of $B_Z$.  Hence
\[
 \begin{aligned}
 B_{i,Z}\cdot F\cdot A^{n-2}
 &=\sum_\beta\operatorname{coeff}_{E_\beta}(B_i)\,
   (E_\beta\cdot F\cdot A^{n-2})\\
 &\longrightarrow
 \sum_\beta\operatorname{coeff}_{E_\beta}(B)\,
   (E_\beta\cdot F\cdot A^{n-2})\\
 &=B_Z\cdot F\cdot A^{n-2},
 \end{aligned}
\]
a finite sum of convergent scalar factors, using that
$\operatorname{coeff}_{E_\beta}(B_i)\to\operatorname{coeff}_{E_\beta}(B)$
by \Cref{prop:ch35-weil-coefficient-bijection}.  Together with
$\theta_i\to\theta$ in the finite-dimensional space $N^1(X)$, this gives
the required inequality in the limit.  This is
\cite[Lemma~6.37]{BJ22}.
\end{proof}

Before restricting to smooth models we record the criterion that replaces
semiampleness by nefness.  It is stated for normal $X$, with the envelope
hypothesis displayed, because it is the only statement in this chapter that
has one.

\begin{theorem}
\label{thm:ch35-cartier-psh-criterion}
Let $X$ be normal and projective, and assume that every ample class on
$X$ has the envelope property $\operatorname{EP}$ of
\Cref{def:bj-envelope-property}.  Let $\theta\in N^1(X)$ and let $B$ be a
rational Cartier b-divisor with determination $\pi:Y\to X$.  Then
\begin{equation}
\label{eq:ch35-cartier-psh-criterion}
 \psi_B\in\PSH_{\mathrm{hom},X}(\theta)
 \iff
 B\leq0\ \text{ and }\ \pi^*\theta+B_Y\ \text{ is nef}.
\end{equation}
Here the left-hand side means that $\psi_B|_{X^{\divis}}$ is the
restriction of a (necessarily unique, by
\Cref{cor:ch35-divisorial-determination}) homogeneous $\theta$-psh
function.
\end{theorem}

\begin{proof}
This is \cite[Theorem~6.35]{BJ22}.  The hypothesis holds when $X$ is
smooth and $\operatorname{char}k=0$, by \cite[Corollary~5.21]{BJ22}, and it
holds unconditionally for integral projective unibranch $X$ once
\Cref{thm:bj-final-envelope} is available; but that theorem is proved in
\Crefrange{chap:bj-cartier-realization}{chap:bj-subintegral-normalization}
using the present chapter, so \Cref{thm:ch35-cartier-psh-criterion} is never
invoked there.  We prove the implication $\Rightarrow$, which is
unconditional, and the reduction of $\Leftarrow$ to the ample case, and we
cite the ample case.

Assume $\psi_B\in\PSH_{\mathrm{hom},X}(\theta)$.  Then
$\psi_B\leq0$ by \Cref{prop:ch35-homogeneous-normalization}, so
$B\leq0$ by \Cref{prop:ch35-weil-coefficient-bijection}.  By
\Cref{cor:ch35-homogenization-psh} there is a decreasing net
$\psi_i\in\cH_{\mathrm{hom}}(L_i)$ with
$c_1(L_i)\to\theta$ and $\psi_i\searrow\psi_B$ pointwise.  Let
$B_i$ be the rational Cartier b-divisor with $\psi_{B_i}=\psi_i$, which
exists by \Cref{prop:cartier-homogeneous-pl}.  By
\Cref{thm:ch35-semiample-criterion}, $B_i\leq0$ and
$L_i+B_i$ is semiample, hence $c_1(L_i)+B_i$ is nef.  The net
$(B_i)$ is decreasing and converges pointwise on $X^{\divis}$ to
$B$, by \Cref{prop:ch35-weil-coefficient-bijection}.  Now
\Cref{lem:ch35-nef-limit} gives that $\theta+B$ is nef, that is,
$\pi^*\theta+B_Y$ is nef.

For the converse, we reduce to the ample case.  Choose an effective Cartier
divisor $F$ on $X$ with $\omega:=\theta+[F]$ ample, possible because
$N^1(X)$ is spanned by ample classes and $\theta$ may be corrected by a
sufficiently positive multiple of a very ample divisor.  Put
$B':=B-\overline F$, a rational Cartier b-divisor with the same
determination $Y$, and note
\[
 \pi^*\omega+B'_Y=\pi^*\theta+\pi^*[F]+B_Y-\pi^*[F]=\pi^*\theta+B_Y,
\]
which is nef by hypothesis; also $B'\leq B\leq0$.  Moreover
$\psi_{\overline F}(v)=v(F)=-\log|s_F|(v)$ by
\eqref{eq:section-bdivisor-sign}, so
$\psi_{B'}=\psi_B+\log|s_F|$ on $X^{\val}$.  Granting the ample case for
$(\omega,B')$, we obtain
$\psi_{B'}\in\PSH_{\mathrm{hom},X}(\omega)$, and then
$\psi_{B'}\leq\log|s_F|$ because $\psi_B\leq0$.
\Cref{thm:ch35-homogeneous-siu-general} yields
$\psi_B=\psi_{B'}-\log|s_F|\in\PSH_{\mathrm{hom},X}(\omega-[F])
=\PSH_{\mathrm{hom},X}(\theta)$, as required.

The ample case itself is \cite[Theorem~6.35]{BJ22}: with
$\theta=\omega$ ample, choose a $\pi$-ample Cartier divisor $A$ on $Y$. An effective Cartier divisor
$F$ on $X$ can be chosen with $\pi^*F\geq A$: a sufficiently positive
hypersurface contains the images of all finitely many positive prime
components of $A$, and multiplying it makes its pullback coefficients large
enough. Then $H:=A-\pi^*F\leq0$ remains $\pi$-ample. Choose
$L\in\Pic(X)_{\mathbf Q}$ with $c_1(L)-\omega$ ample, so that
$\pi^*L+(B_Y+\varepsilon H)$ is ample for small
$\varepsilon\in\mathbf Q_{>0}$; \Cref{thm:ch35-semiample-criterion} then
gives $\psi_B+\varepsilon\psi_H\in\cH_{\mathrm{hom}}(L)$, and letting
$\varepsilon\downarrow0$ produces an \emph{increasing} net converging to
$\psi_B$ on $X^{\divis}$.  It is exactly at this point that the envelope
property for the ample class $c_1(L)$ is used, through
\cite[Theorem~5.11\,(iii)]{BJ22}, to obtain a homogeneous psh extension with the prescribed divisorial
limit. Values of $\psi_B$ on $X^{\An}$ have not been assumed in advance.
Letting $L$ vary and applying
\Cref{thm:bj-basic-closure}\,(5) gives
$\psi_B\in\PSH_{\mathrm{hom},X}(\omega)$.
\end{proof}

\begin{warning}{Two criteria, two hypotheses}
\Cref{thm:ch35-semiample-criterion} is unconditional and characterizes
$\cH_{\mathrm{hom}}(L)$; \Cref{thm:ch35-cartier-psh-criterion}
characterizes $\PSH_{\mathrm{hom},X}(\theta)$ and uses the envelope
property for ample classes.  The difference is one increasing limit.  Since
\Crefrange{chap:bj-cartier-realization}{chap:bj-subintegral-normalization}
are proving the envelope property, they may use the first and must not use
the second.  \Cref{thm:bj-real-cartier-realization} is precisely the
substitute: it realizes an antieffective real Cartier b-divisor with nef
total class as a homogeneous psh function \emph{without} an increasing
limit, by inserting an ample buffer and taking a decreasing net of rational
stages.
\end{warning}

\subsection{Numerical b-divisor classes}
\label{sec:ch35-numerical-bclasses}

We state the original numerical theorem in its actual range.  Assume in this
section that $X$ is smooth and $\operatorname{char}k=0$, and let smooth
models be cofinal.  Set
\begin{equation}
\label{eq:bdivisor-class-space}
 N_b^1(X):=\varprojlim_Y N^1(Y),
\end{equation}
where transition maps are pushforwards.  The nef cone
$\Nef_b(X)$ is the closure, in the projective-limit topology, of the nef
Cartier b-divisor classes.  Nets, rather than sequences, occur in this
definition.

The two hypotheses are exactly those isolated in
\Cref{sec:appF-numerical-bclasses}: one needs a cofinal category of models on
which divisor pushforward stays inside the chosen numerical space, and a
topology in which to take the inverse limit.  Resolution of singularities in
characteristic zero supplies the first, by making smooth models cofinal, and
the projective-limit topology supplies the second.  In the notation of
\eqref{eq:appF-numerical-limits}, $N^1_b(X)=N^1_{\mathrm{Weil}}(\mathfrak X)$
and \Cref{def:appF-nef-weil-class} is the definition of $\Nef_b(X)$ used
here.  Beyond this section we do not use $N^1_b(X)$; the envelope proof
works with Cartier b-divisors and one fixed model at a time.

\begin{definition}[Exceptional b-divisor]
\label{def:ch35-exceptional-bdivisor}
A real Weil b-divisor $B$ over $X$ is \emph{exceptional} if
$B_X=0$.  Equivalently, by
\Cref{prop:ch35-weil-coefficient-bijection}, its coefficient function
satisfies $\psi_B(\ord_E)=0$ for every prime divisor $E\subset X$.  We
write
\[
 \mathcal E:=\{\psi\in\mathcal W\mid
 \psi(\ord_E)=0\text{ for every prime divisor }E\subset X\},
\]
topologized by pointwise convergence.  Membership in $\mathcal W$ is
essential: it includes the modelwise finite-support condition that makes
the function the coefficient function of a Weil b-divisor.
\end{definition}

For $\alpha=(\alpha_Y)_Y\in N_b^1(X)$, there is a unique exceptional
b-divisor $B_\alpha$, with $(B_\alpha)_X=0$, such that
\begin{equation}
\label{eq:exceptional-bdivisor-class}
 [(B_\alpha)_Y]=\alpha_Y-\pi^*\alpha_X.
\end{equation}
The negativity lemma gives uniqueness.  Write
$\psi_\alpha:=\psi_{B_\alpha}$.

\begin{lemma}
\label{lem:ch35-bclass-splitting}
Under the hypotheses of this section, the map
\begin{equation}
\label{eq:ch35-bclass-splitting}
 \alpha\longmapsto(\alpha_X,\psi_\alpha)
\end{equation}
is an isomorphism of topological real vector spaces
$N^1_b(X)\xrightarrow{\ \sim\ }N^1(X)\times\mathcal E$.
\end{lemma}

\begin{proof}
\emph{Existence and uniqueness of $B_\alpha$.}  Fix
$\alpha\in N^1_b(X)$ and a smooth model $\pi:Y\to X$.  The class
$\gamma:=\alpha_Y-\pi^*\alpha_X\in N^1(Y)$ satisfies
$\pi_*\gamma=\alpha_X-\alpha_X=0$, because divisor pushforward is the
transition map in \eqref{eq:bdivisor-class-space} and
$\pi_*\pi^*=\operatorname{id}$ on $N^1(X)$ by the projection formula
\eqref{eq:appF-projection} with $\pi_*[Y]=[X]$.

\emph{Existence.}  Represent $\gamma$ by an $\mathbf R$-divisor $D$ on
$Y$ and set $F:=D-\pi^*\pi_*D$.  Then
$\pi_*F=\pi_*D-\pi_*D=0$, so $F$ is $\pi$-exceptional, and
$[F]=\gamma-\pi^*[\pi_*D]=\gamma$ since $[\pi_*D]=\pi_*\gamma=0$.

\emph{Uniqueness.}  If $F,F'$ are $\pi$-exceptional $\mathbf R$-divisors
with $[F]=[F']$, then $G:=F-F'$ is $\pi$-exceptional and numerically
trivial on $Y$; in particular $G\cdot C=0$ for every $\pi$-contracted
curve $C$, so both $G$ and $-G$ are $\pi$-nef.  The negativity
lemma \cite[Lemma~3.39]{KM08} applied to each of them gives $G\geq0$ and
$G\leq0$, whence $G=0$.

Compatibility of the resulting divisors under pushforward follows from
uniqueness on each model: if $\mu:Y'\to Y$, then $\mu_*F_{Y'}$ is
$\pi$-exceptional with class $\mu_*(\alpha_{Y'}-\mu^*\pi^*\alpha_X)
=\alpha_Y-\pi^*\alpha_X$, so it equals $F_Y$.  Hence $B_\alpha$
exists, is unique, and is exceptional.

\emph{Injectivity.}  If $\alpha_X=0$ and $\psi_\alpha=0$, then
$B_\alpha=0$ by \Cref{prop:ch35-weil-coefficient-bijection}, so
$\alpha_Y=\pi^*\alpha_X=0$ for every $Y$, that is, $\alpha=0$.

\emph{Surjectivity.}  Given $\theta\in N^1(X)$ and $\psi\in\mathcal E$,
let $B$ be the exceptional Weil b-divisor with $\psi_B=\psi$ and set
$\alpha_Y:=\pi^*\theta+[B_Y]$.  Compatibility under pushforward is
$\mu_*\bigl(\mu^*\pi^*\theta+[B_{Y'}]\bigr)=\pi^*\theta+[B_Y]$, again by
the projection formula and $\mu_*B_{Y'}=B_Y$.  Thus
$\alpha\in N^1_b(X)$ with $\alpha_X=\theta$ and
$B_\alpha=B$.

\emph{Topologies.}  The projective-limit topology on $N^1_b(X)$ is that of
convergence of every trace $\alpha_Y$; by
\eqref{eq:exceptional-bdivisor-class} this is equivalent to convergence of
$\alpha_X$ in $N^1(X)$ together with convergence of the exceptional
divisor classes on every model.  On a fixed $Y$ the exceptional divisors
form a finite-dimensional space on which, by the uniqueness just proved, the
class map is injective and hence a linear homeomorphism onto its image; so
convergence of classes is equivalent to coefficientwise convergence, that
is, to pointwise convergence of $\psi_\alpha$ on $X^{\divis}$, by
\Cref{prop:ch35-weil-coefficient-bijection}.  This is
\cite[Lemma~6.39]{BJ22}.
\end{proof}

\begin{theorem}\label{thm:bj-nef-bdivisor-correspondence}
Under the smooth characteristic-zero hypotheses above, the map
\begin{equation}
\label{eq:nef-bdivisor-psh-map}
 \alpha\longmapsto(\alpha_X,\psi_\alpha)
\end{equation}
identifies $\Nef_b(X)$ with the pairs $(\theta,\psi)$ such that
\begin{equation}
\label{eq:nef-bdivisor-psh-pairs}
 \psi\in\PSH_{\mathrm{hom},X}(\theta),\qquad
 \psi(\ord_E)=0
 \quad\text{for every prime divisor }E\subset X.
\end{equation}
Every nef b-divisor class is the limit of a decreasing net of nef Cartier
b-divisors.
\end{theorem}

\begin{proof}[Proof architecture and historical dependency]
For a nef Cartier b-divisor, relative semiampleness after a small ample
perturbation produces homogeneous FS approximants to its coefficient
function.  Passing to a general nef class uses projective-limit compactness
and decreasing Cartier approximation.  Conversely, a decreasing homogeneous
GF regularization of $\psi$ yields nef Cartier b-divisor classes; their
traces converge on every model and reconstruct $\alpha$.

In the original proof, the intermediate Cartier criterion
$\psi_B\in\PSH_{\mathrm{hom}}(\theta)$ iff
$B\leq0$ and $\pi^*\theta+B_Y$ is nef uses the envelope property for
ample classes.  That property was available here because $X$ is smooth in
characteristic zero.  Thus the theorem is not being asserted in arbitrary
characteristic by this proof.  See \cite[Theorems~6.35 and~6.40,
Corollary~6.41]{BJ22}.  \Cref{chap:bj-cartier-realization} supplies a
different realization lemma
for the later general argument.

We now give the argument in the four steps it decomposes into.  Throughout,
$\operatorname{EP}$ for ample classes on $X$ is available, by
\cite[Corollary~5.21]{BJ22}, because $X$ is smooth over a field of
characteristic zero; so
\Cref{thm:ch35-cartier-psh-criterion} may be used.

\medskip
\noindent\emph{Step 1: nef Cartier b-divisors.}
Let $B$ be a rational nef Cartier b-divisor with a smooth determination
$\pi:Y\to X$, and let $\alpha:=[B]$, that is,
$\alpha_{Y'}:=[B_{Y'}]$ for all $Y'$.  Set
$\theta:=\alpha_X=[B_X]$.  The exceptional b-divisor of
\eqref{eq:exceptional-bdivisor-class} is the Cartier b-divisor determined on
$Y$ by $B_Y-\pi^*B_X$, which is $\pi$-exceptional, so
\[
 \psi_\alpha=\psi_{B}-\psi_{\overline{B_X}}\in\PL_{\mathrm{hom},\mathbf R}(X)
\]
by \Cref{prop:cartier-homogeneous-pl}, and the total class of the
exceptional b-divisor in the class $\theta$ is
$\pi^*\theta+[B_Y-\pi^*B_X]=[B_Y]$, which is nef by hypothesis.  Applying
\Cref{thm:ch35-cartier-psh-criterion} to the b-divisor
$B-\overline{B_X}$ --- which is $\leq0$, since it is exceptional and
its total class is nef, whence its trace is antieffective by the negativity
lemma --- gives $\psi_\alpha\in\PSH_{\mathrm{hom},X}(\theta)$.  By
construction $\psi_\alpha(\ord_E)=0$ for prime divisors $E\subset X$.

\medskip
\noindent\emph{Step 2: general nef classes.}
Rational nef Cartier classes are dense among real nef Cartier classes.
On a fixed smooth determination, add $\varepsilon A$ with A ample to a
real nef divisor, then approximate its finitely many coefficients by
rationals inside the open ample cone; let $\varepsilon\downarrow0$.
Continuity of pullback and pushforward gives convergence on every model.
Thus for $\alpha\in\Nef_b(X)$ we may choose a net $(B_i)$ of rational
nef Cartier b-divisors with $\alpha_i:=[B_i]\to\alpha$ in $N^1_b(X)$; by
\Cref{lem:ch35-bclass-splitting} this means
$\theta_i:=(\alpha_i)_X\to\theta:=\alpha_X$ in $N^1(X)$ and
$\psi_{\alpha_i}\to\psi_\alpha$ pointwise on $X^{\divis}$.  Fix
$\omega\in\Amp(X)$.  For $i$ large, $\theta+\omega-\theta_i$ is ample,
so by Step 1 and \Cref{thm:bj-basic-closure}\,(3),
\[
 \psi_{\alpha_i}\in\PSH_{\mathrm{hom},X}(\theta_i)
 \subseteq\PSH_{\mathrm{hom},X}(\theta+\omega).
\]
The space $\PSH_{\mathrm{hom},X}(\theta+\omega)$ is compact for the
topology of pointwise convergence on $X^{\divis}$, by
\cite[Corollary~6.22]{BJ22}, whose hypothesis --- $\operatorname{EP}$ for
ample classes --- is available here.  Passing to a subnet, we may assume
$\psi_{\alpha_i}\to\psi\in\PSH_{\mathrm{hom},X}(\theta+\omega)$
pointwise on $X^{\divis}$; since also
$\psi_{\alpha_i}\to\psi_\alpha$ there, we get
$\psi_\alpha=\psi\in\PSH_{\mathrm{hom},X}(\theta+\omega)$.  As
$\omega\in\Amp(X)$ was arbitrary and
$\theta+\omega\to\theta$ along a net of ample classes,
\Cref{thm:bj-basic-closure}\,(5) gives
$\psi_\alpha\in\PSH_{\mathrm{hom},X}(\theta)$.  Vanishing on the prime
divisors of $X$ passes to the limit.

\medskip
\noindent\emph{Step 3: the converse.}
Let $\theta\in N^1(X)$ and $\psi\in\PSH_{\mathrm{hom},X}(\theta)$ with
$\psi(\ord_E)=0$ for every prime divisor $E\subset X$.  Its restriction
to $X^{\divis}$ lies in $\mathcal E$: on each model every nonexceptional
prime is the strict transform of a prime on X, where the coefficient
vanishes. Only the finitely many exceptional prime divisors can have
nonzero coefficients. This verifies the modelwise finite-support condition
in \Cref{def:ch35-exceptional-bdivisor}, so by
\Cref{prop:ch35-weil-coefficient-bijection} there is an exceptional real
Weil b-divisor $B$ with $\psi_B=\psi$; define
$\alpha\in N^1_b(X)$ by $\alpha_Y:=\pi^*\theta+[B_Y]$, as in the
surjectivity part of \Cref{lem:ch35-bclass-splitting}.

By \Cref{cor:ch35-homogenization-psh} write
$\psi_i\searrow\psi$ with $\psi_i\in\cH_{\mathrm{hom}}(L_i)$ and
$c_1(L_i)\to\theta$.  Since $\psi\leq\psi_i\leq0$ and
$\psi(\ord_E)=0$ for prime divisors $E\subset X$, also
$\psi_i(\ord_E)=0$ for such $E$; hence the Cartier b-divisor $B_i$
with $\psi_{B_i}=\psi_i$, furnished by
\Cref{prop:cartier-homogeneous-pl}, satisfies $(B_i)_X=0$.  Choose a
$\mathbf Q$-Cartier divisor $D_i$ on $X$ with
$[D_i]=c_1(L_i)$ and set $B_i':=\overline{D_i}+B_i$.  By
\Cref{thm:ch35-semiample-criterion}, $L_i+B_i$ is semiample, so
$B_i'$ is a semiample, hence nef, Cartier b-divisor.  Its class satisfies
$(\,[B_i']\,)_X=c_1(L_i)\to\theta$ and
$\psi_{[B_i']}=\psi_{B_i}=\psi_i\to\psi$ pointwise on $X^{\divis}$, so
\Cref{lem:ch35-bclass-splitting} gives $[B_i']\to\alpha$ in
$N^1_b(X)$.  Therefore $\alpha\in\Nef_b(X)$.

\medskip
\noindent\emph{Step 4: monotone approximation.}
This is \Cref{cor:ch35-monotone-approximation} below.
\end{proof}

\begin{corollary}
\label{cor:ch35-monotone-approximation}
Under the hypotheses of
\Cref{thm:bj-nef-bdivisor-correspondence}, every
$\alpha\in\Nef_b(X)$ is the limit of a \emph{decreasing} net $(B_i)$ of
nef Cartier b-divisors, $[B_i]\to\alpha$ in $N^1_b(X)$.
\end{corollary}

\begin{proof}
Write $\alpha_X=[D]$ for an $\mathbf R$-divisor $D$ on $X$ and
choose an effective ample $\mathbf R$-divisor $H$ on $X$ such that
$L:=D+H$ is an ample $\mathbf Q$-divisor. Concretely, since $X$ is smooth,
round the finitely many coefficients of $D$ upwards to obtain an integral
Cartier divisor $D^+$, and choose an effective very ample Cartier divisor $A$.
For a sufficiently large integer $q$, both $L:=D^++qA$ and
$H:=(D^+-D)+qA$ are ample, with $H$ effective and $L$ integral.  By
\Cref{thm:bj-nef-bdivisor-correspondence},
$\psi_\alpha\in\PSH_{\mathrm{hom},X}(\alpha_X)$, and adding the
homogeneous psh function $\psi_{-\overline H}=\log|s_H|$ gives the
following sum. Here, for the effective real Cartier divisor
$H=\sum_jh_jF_j$, define
$\log|s_H|:=\sum_{j:h_j>0}h_j\log|s_{F_j}|$ on $X^{\An}$.
It is homogeneous $[H]$-psh by positive linear-combination closure and
has divisorial coefficients $-H$. No line bundle with class $H$ or FS
presentation for an irrational divisor is asserted. Thus
\[
 \psi_\alpha+\log|s_H|\in\PSH_{\mathrm{hom},X}(c_1(L)).
\]
By \Cref{cor:ch35-homogenization-psh} in its ample form, there is a
decreasing net $(\psi_i)$ in $\cH_{\mathrm{hom}}(L)$ with
$\psi_i\searrow\psi_\alpha+\log|s_H|$ pointwise on $X^{\divis}$.  Write
$\psi_i=\psi_{C_i}$ with $C_i\in\operatorname{CDiv}_{\mathbf Q}(\mathfrak X)$
by \Cref{prop:cartier-homogeneous-pl}.  By
\Cref{thm:ch35-semiample-criterion}, $B_i:=C_i+\overline L$ is a semiample,
hence nef, Cartier b-divisor.  Let $B_\alpha$ be the exceptional Weil
b-divisor attached to $\alpha$.  Since
$\log|s_H|=\psi_{-\overline H}$, the coefficient pairing gives
\[
 C_i\searrow B_\alpha-\overline H,
 \qquad
 B_i=C_i+\overline L\searrow B_\alpha+\overline D
\]
coefficientwise on every model.  In particular $(B_i)$ is decreasing.
The class of the limiting real Weil b-divisor
$B_\alpha+\overline D$ is precisely $\alpha$, by
\eqref{eq:exceptional-bdivisor-class}; continuity of the passage from a
trace divisor to its numerical class, therefore gives $[B_i]\to\alpha$. Indeed
$-H\leq(C_i)_X\leq0$, so the base traces have the fixed finite support
of $H$. Their coefficients converge, hence their numerical classes do too.
On any model only its finitely many exceptional primes contribute in
addition; \Cref{lem:ch35-bclass-splitting} then gives the required
projective-limit convergence.
Notice that the base traces of the approximants need only converge to
$[D]=\alpha_X$; they are not claimed to equal it term by term.  This
recovers
\cite[Theorem~A]{DF20}; see \cite[Corollary~6.41]{BJ22}.
\end{proof}

\begin{remark}[Relation to the intersection theory of
             \Cref{app:intersection-bdivisors}]
\label{rem:ch35-df-relation}
\Cref{cor:ch35-monotone-approximation} is the monotone approximation input
\cite[Theorem~2.1]{DF20} on which
\Cref{thm:appF-df-nef-intersection} rests, and the derivation just given is
logically independent of that appendix: the appendix imports the statement
from \cite{DF20} and builds the intersection product on it, whereas here it
is deduced from the psh correspondence.  Neither derivation is circular,
because \Cref{app:intersection-bdivisors} nowhere uses
\Cref{thm:bj-nef-bdivisor-correspondence}.  The reader who wants the
intersection numbers $(\alpha_1\cdots\alpha_n)_{\mathfrak X}$ of nef
b-classes, and the volume \eqref{eq:appF-nef-volume}, should read the
appendix; the reader who wants the analytic meaning of a nef b-class should
read the present section.  The standard supply of nef b-classes is the
positive part of a divisorial Zariski decomposition
\cite{Bou04,Nak04,BFJ09}.
\end{remark}

\begin{example}[Nef b-classes on a curve and on a surface]
\label{ex:ch35-nef-bclass-examples}
On a smooth projective curve $X$, every model is $X$, so
$N^1_b(X)=N^1(X)=\mathbf R$, $\mathcal E=0$, and
$\Nef_b(X)=\mathbf R_{\geq0}$.  \Cref{thm:bj-nef-bdivisor-correspondence}
then says that for $\theta\geq0$ the only homogeneous $\theta$-psh
function vanishing on all prime divisors of $X$ is $0$, which is
\Cref{prop:ch35-curve-classification} with all $c_r=0$.  The content of the
theorem is therefore entirely two-dimensional and higher.

On $X=\mathbf P^2$, take $\alpha:=[B]$ for the b-divisor $B$ of
\Cref{ex:ch35-blowup-bdivisor}, twisted into a nef class: put
$\theta=c_1(\cO(1))$ and $\alpha_Y:=\pi^*\theta-[E]=H-E$ on the blowup
$Y$ at $q$, and $\alpha_{Y'}$ the corresponding pullback on higher
models.  Then $\alpha_X=\theta$, $\psi_\alpha=\log|\mathfrak m_q|$, and
\eqref{eq:nef-bdivisor-psh-pairs} is satisfied:
$\log|\mathfrak m_q|\in\PSH_{\mathrm{hom},X}(\theta)$ and it vanishes on
every prime divisor of $\mathbf P^2$, since a curve in $\mathbf P^2$ has
$\ord_C(\mathfrak m_q)=0$ unless $C$ passes through $q$, and even then
$\ord_C(\mathfrak m_q)=0$ because a prime divisor on $X$ itself is not
exceptional.  The corresponding b-volume is
$(\alpha^2)_{\mathfrak X}=(H-E)^2=0$ by
\eqref{eq:appF-blowup-table}, so by
\Cref{cor:appF-nef-volume} no probability normalization is available for
this class.
\end{example}

\section{Translation between the homogeneous and inhomogeneous categories}
\label{sec:ch35-translation}

\subsection{The two functors}
\label{sec:ch35-two-functors}

We now assemble the results of
\Cref{sec:ch35-homogenization,sec:ch35-legendre} into a single statement
with its domain of definition made explicit.  This is the form in which the
translation is used in
\Crefrange{chap:bj-cartier-realization}{chap:bj-subintegral-normalization},
and it is the step in the plan's proof architecture that the audit singles
out.

\begin{definition}[The Legendre category]
\label{def:ch35-legendre-category}
Fix $\theta\in N^1(X)$.  A \emph{concave homogeneous family} in the class
$\theta$ is a pair $(b,\psi)$ consisting of a real number
$b\in\mathbf R$ and a family
$\psi=(\psi_\lambda)_{\lambda\leq b}$ such that
\begin{enumerate}
\item $\psi_\lambda\in\PSH_{\mathrm{hom},X}(\theta)$ for every
 $\lambda\leq b$;
\item for every $v\in X^{\divis}$, the function
 $\lambda\mapsto\psi_\lambda(v)$ is concave and nonincreasing on
 $(-\infty,b]$.
\end{enumerate}
We write $\mathfrak L_X(\theta)$ for the set of such pairs, ordered by
$(b,\psi)\leq(b',\psi')$ when $b\leq b'$ and
$\psi_\lambda\leq\psi'_\lambda$ for all $\lambda\leq b$.
\end{definition}

Condition (2) is imposed on $X^{\divis}$ only.  By
\Cref{cor:ch35-divisorial-determination} this is no loss: a concave
inequality between homogeneous psh functions that holds on
$X^{\divis}$ holds on $X^{\An}$.  Condition (1) already forces
$\psi_\lambda\leq0$, by
\Cref{prop:ch35-homogeneous-normalization}.

\begin{theorem}
\label{thm:ch35-translation}
Fix $\theta\in N^1(X)$.  Define
\[
 \Lambda:\PSH_X(\theta)\longrightarrow\mathfrak L_X(\theta),
 \qquad
 \Lambda(\phi):=\bigl(\sup\phi,\ (\widehat\phi_\lambda)_{\lambda\leq\sup\phi}\bigr),
\]
and, for $(b,\psi)\in\mathfrak L_X(\theta)$,
\[
 \Sigma(b,\psi):=\sup_{\lambda\leq b}\{\psi_\lambda+\lambda\}
 \qquad\text{pointwise on }X^{\An}.
\]
Then:
\begin{enumerate}
\item $\Lambda$ is well defined: for
 $\phi\in\PSH_X(\theta)$, the pair $\Lambda(\phi)$ satisfies
 \Cref{def:ch35-legendre-category}, and moreover
 $\lambda\mapsto\widehat\phi_\lambda(v)$ is continuous on
 $(-\infty,\sup\phi]$ for every $v\in X^{\divis}$;
\item $\Sigma\circ\Lambda=\operatorname{id}$ on
 $\PSH_X(\theta)$; in particular $\Lambda$ is injective and
 $\Sigma$ is surjective onto $\PSH_X(\theta)$ when restricted to the
 image of $\Lambda$;
\item $\Lambda\circ\Sigma=\operatorname{id}$ on the image of
 $\Lambda$; more precisely, if
 $(b,\psi)\in\mathfrak L_X(\theta)$ and
 $\phi:=\Sigma(b,\psi)$ happens to satisfy
 $\phi\in\PSH_X(\theta)$ with $\sup\phi=b$, then
 $\psi_\lambda=\widehat\phi_\lambda$ for every $\lambda<b$;
\item $\Lambda$ and $\Sigma$ are order preserving, and
 $\Lambda(\phi+c)=\bigl(\sup\phi+c,\ (\widehat\phi_{\lambda-c})_{\lambda}\bigr)$
 for $c\in\mathbf R$, while
 $\Lambda(t\cdot\phi)=\bigl(t\sup\phi,\ (\widehat\phi_{\lambda/t})_\lambda\bigr)$
 for $t>0$.
\end{enumerate}
\par
\end{theorem}

\begin{proof}
(1)  Membership in $\PSH_{\mathrm{hom},X}(\theta)$ for
$\lambda\leq b:=\sup\phi$ is
\Cref{prop:ch35-slice-trichotomy}\,(2); concavity and monotonicity in
$\lambda$ are \Cref{lem:ch35-slice-basic}\,(6); continuity on
$(-\infty,b]$ at divisorial points is
\Cref{lem:ch35-slice-continuity}.

(2)  This is exactly
\Cref{thm:homogeneous-legendre-decomposition}\,(3).

(3)  Assume $\phi=\Sigma(b,\psi)$ lies in
$\PSH_X(\theta)$ with $\sup\phi=b$.  Fix $v\in X^{\An}$ and
$t>0$.  Apply the scaling action to the defining identity: by
\Cref{lem:ch35-action-basic}\,(3)--(4) and homogeneity of each
$\psi_\lambda$,
\[
 (t\cdot\phi)(v)
 =\sup_{\lambda\leq b}\bigl(t\cdot(\psi_\lambda+\lambda)\bigr)(v)
 =\sup_{\lambda\leq b}\{\psi_\lambda(v)+t\lambda\}.
\]
Thus, for each fixed $v$, the function
$F(t):=(t\cdot\phi)(v)$ is the supremum of the affine functions
$t\mapsto\psi_\lambda(v)+t\lambda$, that is, $F=\mathcal Lg$ in the
notation \eqref{eq:appD-concave-transform}, where
$g(\lambda):=\psi_\lambda(v)$ for $\lambda\leq b$ and
$g(\lambda):=-\infty$ for $\lambda>b$.  The function $g$ is concave
and nonincreasing by hypothesis, and it is finite on $(-\infty,b]$ when $v\in X^{\divis}$, by
\Cref{thm:bj-basic-closure}\,(7); being finite and concave, it is
continuous on the open half-line $(-\infty,b)$.  It need \emph{not} be
upper semicontinuous at $\lambda=b$: a finite concave function on
$(-\infty,b]$ may drop at the endpoint --- take $g=0$ on
$(-\infty,0)$ and $g(0)=-1$, which is concave --- and
\Cref{def:ch35-legendre-category}\,(2) does not forbid it.  This costs
nothing here, because the conclusion is asserted only for $\lambda<b$:
the closure $\operatorname{cl}g$ is proper, concave and nonincreasing,
it has the same concave transform as $g$, and it agrees with $g$ on
the relative interior $(-\infty,b)$ of its domain.  So
\Cref{thm:appD-concave-duality} applied to $\operatorname{cl}g$ gives
inversion with $t\geq0$. It is equivalent here to $t>0$. If a supporting
slope of $g$ at an interior parameter $\lambda<b$ is negative, its dual
parameter is already positive. If that slope is zero, concavity and
nonincreasingness imply $g(\mu)=g(\lambda)$ for $\mu\leq\lambda$
and $g(\mu)\leq g(\lambda)$ for all $\mu\leq b$. Hence
$g(\lambda)+t\lambda\leq F(t)\leq g(\lambda)+tb$, so letting
$t\downarrow0$ recovers the same value. Thus
\[
 \psi_\lambda(v)=g(\lambda)=\inf_{t>0}\{F(t)-t\lambda\}
 =\widehat\phi_\lambda(v)
\]
for every $\lambda<b$ and every $v\in X^{\divis}$.  Both
$\psi_\lambda$ and $\widehat\phi_\lambda$ are $\theta$-psh, so
\Cref{cor:ch35-divisorial-determination} upgrades the identity to
$X^{\An}$.  This is the converse half of
\cite[Theorem~6.24]{BJ22}.

(4)  Order preservation of $\Lambda$ is
\Cref{lem:ch35-slice-basic}\,(5) together with monotonicity of
$\phi\mapsto\sup\phi$; order preservation of $\Sigma$ is immediate from
the definition.  The two equivariance formulas are
\Cref{lem:ch35-slice-basic}\,(3) and (4), together with
$\sup(\phi+c)=\sup\phi+c$ and
$\sup(t\cdot\phi)=t\sup\phi$.
\end{proof}

\subsection{The exact domain of the backward direction}
\label{sec:ch35-domain}

\begin{warning}{The backward construction has a smaller domain}
The pointwise expression
\[
 \Sigma(b,\psi)=\sup_{\lambda\leq b}\{\psi_\lambda+\lambda\}
\]
is defined for every $(b,\psi)\in\mathfrak L_X(\theta)$, but it need not
be upper semicontinuous or psh.  Moreover a concave family may have a
downward jump at $\lambda=b$, whereas a slice family coming from a psh
function is closed at the endpoint by
\Cref{lem:ch35-slice-continuity}.  Precisely, the image of $\Lambda$
consists of those $(b,\psi)$ for which
\[
 \Sigma(b,\psi)\in\PSH_X(\theta),\qquad
 \sup\Sigma(b,\psi)=b,\qquad
 \psi_b=\inf_{\lambda<b}\psi_\lambda
 \quad\text{on }X^{\divis}.
\]
Indeed, \Cref{thm:ch35-translation}\,(3) recovers every slice with
$\lambda<b$, and the displayed endpoint condition, together with
\Cref{lem:ch35-slice-continuity}, recovers the last slice.

If $\theta$ has the envelope property, it only follows automatically
that the upper-semicontinuous regularization
$\Sigma(b,\psi)^*$ is $\theta$-psh.  It does not follow that the raw
supremum equals its regularization, nor that the endpoint is closed.
Thus the envelope property alone does not enlarge the formal inverse beyond
the domain just displayed.  The later envelope proof uses $\Sigma$ only
on the image of $\Lambda$, so no circularity is introduced.
\end{warning}

\subsection{Chapter dictionary}
\label{sec:ch35-dictionary}

The Legendre slices in \eqref{eq:homogeneous-legendre-slice} give the
translation between categories:
\begin{description}
\item[Potential.]  The inhomogeneous function
  $\phi\in\PSH_X(\theta)$ corresponds to the homogeneous family
  $(\widehat\phi_\lambda)_{\lambda\leq\sup\phi}$.
\item[Weighted presentation.]  A finite maximum of weighted sections
  corresponds, after homogenization, to weight-zero coefficient systems.
\item[Radial bound.]  The inequality $\phi\leq p+C$ corresponds to control
  of the Legendre tail by the carrier $p$.
\item[Analytic values.]  Values on all of $X^{\An}$ are retained by the
  homogeneous divisorial coefficients together with the parameter $\lambda$.
\end{description}

The last row is not a claim that arbitrary functions are determined by
b-divisors.  Psh functions are determined by their divisorial trace; the
homogeneous Legendre family is what retains the inhomogeneous information.
This distinction is central to the reflection arguments of
\Crefrange{chap:bj-cartier-realization}{chap:bj-reflection}.

The full dictionary established in this chapter is the following.  Each row
is a proved equivalence, with its hypotheses stated in the corresponding
item.

\begin{description}
\item[Positive homogeneous PL.]
  $\psi\in\PL^+_{\mathrm{hom}}(X)$ corresponds to an antieffective,
  relatively semiample Cartier b-divisor when $X$ is normal
  (\Cref{prop:cartier-homogeneous-pl}).
\item[Homogeneous PL.]
  $\psi\in\PL_{\mathrm{hom}}(X)$ corresponds to a Cartier b-divisor when
  $X$ is normal (\Cref{prop:cartier-homogeneous-pl}).
\item[Homogeneous FS.]
  $\psi_B\in\cH_{\mathrm{hom}}(L)$ is equivalent to
  $B\leq0$ and $L+B$ semiample when $X$ is normal
  (\Cref{thm:ch35-semiample-criterion}).
\item[Homogeneous psh, Cartier level.]
  $\psi_B\in\PSH_{\mathrm{hom},X}(\theta)$ is equivalent to
  $B\leq0$ and $\theta+B$ nef when $X$ is normal and the envelope property
  holds for ample classes
  (\Cref{thm:ch35-cartier-psh-criterion}).
\item[Homogeneous psh, numerical level.]
  The pairs $\psi\in\PSH_{\mathrm{hom},X}(\theta)$ with
  $\psi(\ord_E)=0$ for every prime divisor $E\subset X$ correspond to nef
  b-divisor classes $\alpha$ with $\alpha_X=\theta$ when $X$ is smooth and
  $\operatorname{char}k=0$
  (\Cref{thm:bj-nef-bdivisor-correspondence}).
\item[Weil level.]
  Weil b-divisors correspond to modelwise finitely supported homogeneous
  functions on $X^{\divis}$ when $X$ is normal
  (\Cref{prop:ch35-weil-coefficient-bijection}).
\item[Legendre level.]
  Each $\phi\in\PSH_X(\theta)$ determines its concave homogeneous family
  $(\widehat\phi_\lambda)$ under the standing hypotheses
  \eqref{eq:ch35-standing} (\Cref{thm:ch35-translation}).
\end{description}

Three separations are recorded alongside the dictionary and should be
carried forward.

\emph{Cartier versus Weil.}  Pullback and intersection are automatic for
Cartier b-divisors and are not automatic for Weil ones; this is the firewall
of \Cref{app:intersection-bdivisors}, and it is why
\Cref{chap:bj-cartier-realization} realizes \emph{Cartier} b-divisors and
does not attempt an arbitrary Weil realization.

\emph{Divisor versus class.}  A Weil b-divisor is coefficient data on every
model; a numerical b-divisor class is inverse-limit data in
$N^1$.  \Cref{lem:ch35-bclass-splitting} relates them, and the relation
uses the negativity lemma and therefore normality and a suitable model
category.

\emph{Trace versus function.}  A homogeneous psh function is an
extended-real function on $X^{\An}$; its restriction to
$X^{\divis}$ determines it, by
\Cref{cor:ch35-divisorial-determination}, but the divisorial topology does not give convergence at every analytic point.
For general psh functions this is \Cref{ex:ch35-evaluation-discontinuous};
the same failure occurs for the homogeneous sequence
$m^{-1}\log|\mathfrak m_p|$ on $\mathbf P^1$: its divisorial values
tend to zero, while its value at $v_{p,\triv}$ stays $-\infty$.

\section{Source note}

Homogeneous FS functions, homogenization, and Legendre reconstruction are
from \cite[\S\S6.1--6.3]{BJ22}.  The b-divisor material follows
\cite[\S\S6.4--6.5]{BJ22}, with background from \cite{BFJ09,DF20}.  We have
kept three objects separate: a Weil b-divisor is coefficient data on every
model, a b-divisor class is numerical inverse-limit data, and a homogeneous
psh function is an extended-real analytic function determined by its
divisorial values.

More precisely, the attributions are as follows.  The scaling action on
functions and the notion of homogeneity are
\cite[\S1.1.4]{BJ22}; homogeneous PL and Fubini--Study functions, including
\Cref{prop:hom-fs-ideal} and \Cref{cor:ch35-ample-exhausts}, are
\cite[\S6.1]{BJ22}, with the tropical spectrum interpretation of
\Cref{rem:ch35-elimination-tropical} going back to
\cite{JM12}.  Homogenization, the elimination lemma
\Cref{lem:ch35-elimination}, and \Cref{prop:homogenization-gf} are
\cite[\S6.2]{BJ22}; we have given the Fourier--Motzkin argument explicitly
and referred the biconjugation step to
\Cref{thm:appD-fenchel-moreau} rather than to convex analysis in general.
The Legendre decomposition
\Cref{thm:homogeneous-legendre-decomposition}, the continuity
\Cref{lem:ch35-slice-continuity}, the maximal slice
\Cref{prop:ch35-phimax}, and the general homogeneous Siu decomposition
\Cref{thm:ch35-homogeneous-siu-general} are
\cite[\S6.3, Theorem~6.21]{BJ22}.  The strict Siu statement
\Cref{thm:ch35-homogeneous-siu} is the book-local argument that isolates the
endpoint left to that imported theorem.  The idea of encoding an
inhomogeneous potential by a concave family of homogeneous ones is due to
Ross and Witt Nystr\"om \cite{RWN14} in the complex setting, and the
comparison is made in \Cref{chap:rays-filtrations}.  The identification of
Cartier b-divisors with homogeneous PL functions,
\Cref{prop:cartier-homogeneous-pl}, and the semiample and nef criteria
\Cref{thm:ch35-semiample-criterion} and
\Cref{thm:ch35-cartier-psh-criterion}, are \cite[\S6.4]{BJ22}.  The nef
b-class correspondence \Cref{thm:bj-nef-bdivisor-correspondence} and the
monotone approximation \Cref{cor:ch35-monotone-approximation} are
\cite[\S6.5]{BJ22}; the approximation theorem was first proved by a
different and more direct route in \cite{DF20}, whose intersection theory of
nef b-classes is developed in \Cref{app:intersection-bdivisors}.  The
underlying b-divisor formalism goes back to \cite{BFJ09}, and the
divisorial Zariski decomposition that furnishes the standard examples of nef
b-classes is \cite{Bou04,Nak04}.  For the analogous transcendental theory
over $\mathbf C$ see \cite{Xiabdiv,Xiabdiv2}.

The comparison theorem \Cref{thm:ch35-divisorial-comparison} and its
corollaries are \cite[Theorem~4.22 and Corollary~4.23]{BJ22}.  Within its
proof we isolated the valuative criterion of effectivity
\Cref{prop:ch35-effectivity} and derived the finite-attainment step from it.
This makes visible the point at which equality on divisorial points is
upgraded to equality everywhere; and
\Cref{ex:ch35-usc-not-determined} shows that the upgrade is false without
plurisubharmonicity.
